% !TEX program = lualatex
% ---------------------------------------------------------------------------------------------
%  Readback note 42: Box partitions of arithmetic unipotent groups: frame counts, the covolume, amplitude sums
%
%  Build:  python3 tools/rb.py note <slug>
%
%  One unit of a Lean artifact, rewritten as ordinary mathematics. The prose is hand-written. Lean names,
%  line ranges and counts come from data.tex; quoted Lean statements are copied from the artifact by the build.
%
%  Rules of the text: standard terminology; every symbol is explained where it first appears and listed again
%  in the notation at the end; no letter has two meanings; a proof is complete or says what it takes for granted;
%  Lean names appear only in the LEAN lines and in the concordance.
%
%  The signature rule (the build enforces it): a function, relation or predicate that the artifact defines is
%  written only through its notation macro, which the build makes from the spec (`notation`) with one slot for
%  each explicit argument of the Lean definition. No argument is fixed and dropped, also not "for brevity" and
%  not for one section. Every argument and every hypothesis of a stated declaration stands inside its
%  statement, or in a Setting that the statement cites by number; only implicit and instance arguments (the
%  ambient objects) may stand in the conventions.
% ---------------------------------------------------------------------------------------------
\documentclass[11pt,a4paper]{article}
\ifdefined\pdfvariable\pdfvariable suppressoptionalinfo \numexpr 1+2+4+8\relax\fi
\usepackage[a4paper,hmargin=35mm,top=22mm,bottom=27mm,footskip=12mm,nohead]{geometry}

\usepackage{fontspec}
\directlua{
  luaotfload.add_fallback("gfallback", {"STIXTwoMath-Regular.otf:mode=node;", "DejaVuSans.ttf:mode=node;"})
  luaotfload.add_fallback("gmonofallback", {"DejaVuSansMono.ttf:mode=node;", "STIXTwoMath-Regular.otf:mode=node;"})
}
\setmainfont{NewCM10}[Extension=.otf,UprightFont=*-Book,ItalicFont=*-BookItalic,
  BoldFont=*-Bold,BoldItalicFont=*-BoldItalic,SmallCapsFont=*-Book,SmallCapsFeatures={Letters=SmallCaps},
  RawFeature={fallback=gfallback}]
\setsansfont{NewCMSans10}[Extension=.otf,UprightFont=*-Book,ItalicFont=*-BookOblique,
  BoldFont=*-Bold,BoldItalicFont=*-BoldOblique,RawFeature={fallback=gfallback}]
\setmonofont{JetBrainsMono}[Extension=.otf,UprightFont=*-Regular,ItalicFont=*-Italic,
  BoldFont=*-SemiBold,BoldItalicFont=*-SemiBoldItalic,Scale=MatchLowercase,
  RawFeature={-calt;-liga;fallback=gmonofallback}]
\usepackage{amsmath}
\usepackage{amsthm}
\usepackage[math-style=TeX]{unicode-math}
\setmathfont{NewCMMath-Book.otf}
\emergencystretch=2.5em\relax

\usepackage{xcolor}
\usepackage{tikz}
\usetikzlibrary{calc,arrows.meta,positioning,backgrounds}
\usepackage[protrusion=true,expansion=true]{microtype}
\usepackage{enumitem}
\usepackage{needspace}
\usepackage{booktabs}
\usepackage{array}
\usepackage{longtable}
\usepackage{fancyhdr}
\usepackage{lastpage}
\usepackage{multicol}
\usepackage[hidelinks,pdfusetitle]{hyperref}

\definecolor{ink}{HTML}{0F172A}
\definecolor{inksoft}{HTML}{475569}
\definecolor{inkfaint}{HTML}{94A3B8}
\definecolor{rulegrey}{HTML}{CBD5E1}
\definecolor{arith}{HTML}{C2410C}
\definecolor{est}{HTML}{7C3AED}
\definecolor{ana}{HTML}{0891B2}
\definecolor{ok}{HTML}{047857}
\definecolor{leancol}{HTML}{3730A3}

\title{Box partitions of arithmetic unipotent groups: frame counts, the covolume, amplitude sums}
\author{grh-compact}
\date{}

% ------------------------------------------------------------------------------------- macros
\newcommand{\leanfont}{\ttfamily\fontsize{8.6}{10.5}\selectfont\color{leancol}}
\newcommand{\lean}[1]{\csname LN@#1\endcsname}
\newcommand{\fillin}[1]{{\color{arith}\itshape[#1]}}  % a place of the template that is not written yet; the build reports every one that is left
\newcommand{\leanquote}[1]{\csname LQ@#1\endcsname}
% generated by rb.py note; do not edit
\newcommand{\ConcDecls}{68}
\newcommand{\Pin}{72fead10}
\newcommand{\PinText}{that of commit {\leanfont 72fead10}}
\newcommand{\ArtifactName}{grh-compact}
\expandafter\newcommand\csname LN@isStablebesselS\endcsname{\href{https://github.com/lsernau/grh-lean/blob/72fead102f663d5c7a407964a13dc0cf9d977207/grh-compact/src/New/Filling/AmplitudeSums.lean\#L16-L17}{{\leanfont Filling.\allowbreak{}isStable\_besselS}}}
\expandafter\newcommand\csname LN@measurePreservingpiorbitMeasure\endcsname{\href{https://github.com/lsernau/grh-lean/blob/72fead102f663d5c7a407964a13dc0cf9d977207/grh-compact/src/New/Filling/AmplitudeSums.lean\#L25-L28}{{\leanfont Filling.\allowbreak{}measurePreserving\_pi\_orbitMeasure}}}
\expandafter\newcommand\csname LN@tsumintegralapply\endcsname{\href{https://github.com/lsernau/grh-lean/blob/72fead102f663d5c7a407964a13dc0cf9d977207/grh-compact/src/New/Filling/AmplitudeSums.lean\#L30-L35}{{\leanfont Filling.\allowbreak{}tsum\_integral\_apply}}}
\expandafter\newcommand\csname LN@besselFnmulofmemclassUnip\endcsname{\href{https://github.com/lsernau/grh-lean/blob/72fead102f663d5c7a407964a13dc0cf9d977207/grh-compact/src/New/Filling/AmplitudeSums.lean\#L41-L44}{{\leanfont Filling.\allowbreak{}besselFn\_mul\_of\_mem\_classUnip}}}
\expandafter\newcommand\csname LN@besselKernelbesselUmul\endcsname{\href{https://github.com/lsernau/grh-lean/blob/72fead102f663d5c7a407964a13dc0cf9d977207/grh-compact/src/New/Filling/AmplitudeSums.lean\#L46-L49}{{\leanfont Filling.\allowbreak{}besselKernel\_besselU\_mul}}}
\expandafter\newcommand\csname LN@lintegralclassCutoffmulprod\endcsname{\href{https://github.com/lsernau/grh-lean/blob/72fead102f663d5c7a407964a13dc0cf9d977207/grh-compact/src/New/Filling/AmplitudeSums.lean\#L54-L58}{{\leanfont Filling.\allowbreak{}lintegral\_classCutoff\_mul\_prod}}}
\expandafter\newcommand\csname LN@amplitudeAbslttopandcutle\endcsname{\href{https://github.com/lsernau/grh-lean/blob/72fead102f663d5c7a407964a13dc0cf9d977207/grh-compact/src/New/Filling/AmplitudeSums.lean\#L64-L106}{{\leanfont Filling.\allowbreak{}amplitudeAbs\_lt\_top\_and\_cut\_le}}}
\expandafter\newcommand\csname LN@summableamplitude\endcsname{\href{https://github.com/lsernau/grh-lean/blob/72fead102f663d5c7a407964a13dc0cf9d977207/grh-compact/src/New/Filling/AmplitudeSums.lean\#L110-L111}{{\leanfont Filling.\allowbreak{}summable\_amplitude}}}
\expandafter\newcommand\csname LN@hasSumamplitudeSumchangeofweight\endcsname{\href{https://github.com/lsernau/grh-lean/blob/72fead102f663d5c7a407964a13dc0cf9d977207/grh-compact/src/New/Filling/AmplitudeSums.lean\#L119-L158}{{\leanfont Filling.\allowbreak{}hasSum\_amplitudeSum\_change\_of\_weight}}}
\expandafter\newcommand\csname LN@amplitudeSummove\endcsname{\href{https://github.com/lsernau/grh-lean/blob/72fead102f663d5c7a407964a13dc0cf9d977207/grh-compact/src/New/Filling/AmplitudeSums.lean\#L163-L179}{{\leanfont Filling.\allowbreak{}amplitudeSum\_move}}}
\expandafter\newcommand\csname LN@amplitudeSumbesselUmul\endcsname{\href{https://github.com/lsernau/grh-lean/blob/72fead102f663d5c7a407964a13dc0cf9d977207/grh-compact/src/New/Filling/AmplitudeSums.lean\#L181-L187}{{\leanfont Filling.\allowbreak{}amplitudeSum\_besselU\_mul}}}
\expandafter\newcommand\csname LN@amplitudeSumeqzeroofnotagrees\endcsname{\href{https://github.com/lsernau/grh-lean/blob/72fead102f663d5c7a407964a13dc0cf9d977207/grh-compact/src/New/Filling/AmplitudeSums.lean\#L192-L194}{{\leanfont Filling.\allowbreak{}amplitudeSum\_eq\_zero\_of\_not\_agrees}}}
\expandafter\newcommand\csname LN@amplitudeSumbesselSone\endcsname{\href{https://github.com/lsernau/grh-lean/blob/72fead102f663d5c7a407964a13dc0cf9d977207/grh-compact/src/New/Filling/AmplitudeSums.lean\#L201-L220}{{\leanfont Filling.\allowbreak{}amplitudeSum\_besselS\_one}}}
\expandafter\newcommand\csname LQ@amplitudeSumbesselSone\endcsname{\leanquoteline{0}{lemma amplitudeSum\_besselS\_one i j : S.amplitudeSum (σ := .full) (S.ratMiddle E).besselS c 1 c' (i,}\leanquoteline{2}{j) = ite (c = c' ∧ i = j) S.unipCovolume 0}}
\expandafter\newcommand\csname LN@liouvillegram\endcsname{\href{https://github.com/lsernau/grh-lean/blob/72fead102f663d5c7a407964a13dc0cf9d977207/grh-compact/src/New/Filling/BoxPartition.lean\#L45-L55}{{\leanfont Filling.\allowbreak{}liouville\_gram}}}
\expandafter\newcommand\csname LN@linkMetric\endcsname{\href{https://github.com/lsernau/grh-lean/blob/72fead102f663d5c7a407964a13dc0cf9d977207/grh-compact/src/New/Filling/BoxPartition.lean\#L59-L60}{{\leanfont Filling.\allowbreak{}linkMetric}}}
\expandafter\newcommand\csname LN@linkDiscSq\endcsname{\href{https://github.com/lsernau/grh-lean/blob/72fead102f663d5c7a407964a13dc0cf9d977207/grh-compact/src/New/Filling/BoxPartition.lean\#L64-L65}{{\leanfont Filling.\allowbreak{}linkDiscSq}}}
\expandafter\newcommand\csname LN@framePairlasteqzFormPairmul\endcsname{\href{https://github.com/lsernau/grh-lean/blob/72fead102f663d5c7a407964a13dc0cf9d977207/grh-compact/src/New/Filling/BoxPartition.lean\#L72-L74}{{\leanfont Filling.\allowbreak{}framePair\_last\_eq\_zFormPair\_mul}}}
\expandafter\newcommand\csname LN@classCoordeqfromBlocks\endcsname{\href{https://github.com/lsernau/grh-lean/blob/72fead102f663d5c7a407964a13dc0cf9d977207/grh-compact/src/New/Filling/BoxPartition.lean\#L76-L78}{{\leanfont Filling.\allowbreak{}classCoord\_eq\_fromBlocks}}}
\expandafter\newcommand\csname LN@HasIntegralLifts\endcsname{\href{https://github.com/lsernau/grh-lean/blob/72fead102f663d5c7a407964a13dc0cf9d977207/grh-compact/src/New/Filling/BoxPartition.lean\#L86-L88}{{\leanfont Filling.\allowbreak{}HasIntegralLifts}}}
\expandafter\newcommand\csname LN@finranklink\endcsname{\href{https://github.com/lsernau/grh-lean/blob/72fead102f663d5c7a407964a13dc0cf9d977207/grh-compact/src/New/Filling/BoxPartition.lean\#L90-L96}{{\leanfont Filling.\allowbreak{}finrank\_link}}}
\expandafter\newcommand\csname LN@gramclassCoordX\endcsname{\href{https://github.com/lsernau/grh-lean/blob/72fead102f663d5c7a407964a13dc0cf9d977207/grh-compact/src/New/Filling/BoxPartition.lean\#L98-L100}{{\leanfont Filling.\allowbreak{}gram\_classCoordX₀}}}
\expandafter\newcommand\csname LN@classCoordgram\endcsname{\href{https://github.com/lsernau/grh-lean/blob/72fead102f663d5c7a407964a13dc0cf9d977207/grh-compact/src/New/Filling/BoxPartition.lean\#L102-L104}{{\leanfont Filling.\allowbreak{}classCoord\_gram}}}
\expandafter\newcommand\csname LN@linkMetricposDef\endcsname{\href{https://github.com/lsernau/grh-lean/blob/72fead102f663d5c7a407964a13dc0cf9d977207/grh-compact/src/New/Filling/BoxPartition.lean\#L106-L108}{{\leanfont Filling.\allowbreak{}linkMetric\_posDef}}}
\expandafter\newcommand\csname LN@gramlinkFormeqzerooflelink\endcsname{\href{https://github.com/lsernau/grh-lean/blob/72fead102f663d5c7a407964a13dc0cf9d977207/grh-compact/src/New/Filling/BoxPartition.lean\#L110-L113}{{\leanfont Filling.\allowbreak{}gram\_linkForm\_eq\_zero\_of\_le\_link}}}
\expandafter\newcommand\csname LN@detgramlinkMetricnezero\endcsname{\href{https://github.com/lsernau/grh-lean/blob/72fead102f663d5c7a407964a13dc0cf9d977207/grh-compact/src/New/Filling/BoxPartition.lean\#L115-L118}{{\leanfont Filling.\allowbreak{}det\_gram\_linkMetric\_ne\_zero}}}
\expandafter\newcommand\csname LN@cardeqfinrankcolSpan\endcsname{\href{https://github.com/lsernau/grh-lean/blob/72fead102f663d5c7a407964a13dc0cf9d977207/grh-compact/src/New/Filling/BoxPartition.lean\#L121-L123}{{\leanfont Filling.\allowbreak{}card\_eq\_finrank\_colSpan}}}
\expandafter\newcommand\csname LN@gramlinkMatrixmul\endcsname{\href{https://github.com/lsernau/grh-lean/blob/72fead102f663d5c7a407964a13dc0cf9d977207/grh-compact/src/New/Filling/BoxPartition.lean\#L125-L134}{{\leanfont Filling.\allowbreak{}gram\_linkMatrix\_mul}}}
\expandafter\newcommand\csname LN@classCoordmulframePair\endcsname{\href{https://github.com/lsernau/grh-lean/blob/72fead102f663d5c7a407964a13dc0cf9d977207/grh-compact/src/New/Filling/BoxPartition.lean\#L136-L138}{{\leanfont Filling.\allowbreak{}classCoord\_mul\_framePair}}}
\expandafter\newcommand\csname LN@classCoordmulisotropicFramePair\endcsname{\href{https://github.com/lsernau/grh-lean/blob/72fead102f663d5c7a407964a13dc0cf9d977207/grh-compact/src/New/Filling/BoxPartition.lean\#L142-L143}{{\leanfont Filling.\allowbreak{}classCoord\_mul\_isotropicFramePair}}}
\expandafter\newcommand\csname LN@isIntMatrixframePair\endcsname{\href{https://github.com/lsernau/grh-lean/blob/72fead102f663d5c7a407964a13dc0cf9d977207/grh-compact/src/New/Filling/BoxPartition.lean\#L145-L146}{{\leanfont Filling.\allowbreak{}isIntMatrix\_framePair}}}
\expandafter\newcommand\csname LN@legenIndexdetoffixesNormAtLeast\endcsname{\href{https://github.com/lsernau/grh-lean/blob/72fead102f663d5c7a407964a13dc0cf9d977207/grh-compact/src/New/Filling/BoxPartition.lean\#L158-L212}{{\leanfont Filling.\allowbreak{}le\_genIndex\_det\_of\_fixesNormAtLeast}}}
\expandafter\newcommand\csname LN@existsmemclassUnip\endcsname{\href{https://github.com/lsernau/grh-lean/blob/72fead102f663d5c7a407964a13dc0cf9d977207/grh-compact/src/New/Filling/BoxPartition.lean\#L227-L250}{{\leanfont Filling.\allowbreak{}exists\_mem\_classUnip}}}
\expandafter\newcommand\csname LN@existseichlerclassUnip\endcsname{\href{https://github.com/lsernau/grh-lean/blob/72fead102f663d5c7a407964a13dc0cf9d977207/grh-compact/src/New/Filling/BoxPartition.lean\#L252-L263}{{\leanfont Filling.\allowbreak{}exists\_eichler\_classUnip}}}
\expandafter\newcommand\csname LN@tsumleofencardrangele\endcsname{\href{https://github.com/lsernau/grh-lean/blob/72fead102f663d5c7a407964a13dc0cf9d977207/grh-compact/src/New/Filling/BoxPartition.lean\#L270-L274}{{\leanfont Filling.\allowbreak{}tsum\_le\_of\_encard\_range\_le}}}
\expandafter\newcommand\csname LN@tsumleofcolouring\endcsname{\href{https://github.com/lsernau/grh-lean/blob/72fead102f663d5c7a407964a13dc0cf9d977207/grh-compact/src/New/Filling/BoxPartition.lean\#L276-L282}{{\leanfont Filling.\allowbreak{}tsum\_le\_of\_colouring}}}
\expandafter\newcommand\csname LN@encardintegralle\endcsname{\href{https://github.com/lsernau/grh-lean/blob/72fead102f663d5c7a407964a13dc0cf9d977207/grh-compact/src/New/Filling/BoxPartition.lean\#L287-L292}{{\leanfont Filling.\allowbreak{}encard\_integral\_le}}}
\expandafter\newcommand\csname LN@submoduleIndexsmulintVectors\endcsname{\href{https://github.com/lsernau/grh-lean/blob/72fead102f663d5c7a407964a13dc0cf9d977207/grh-compact/src/New/Filling/BoxPartition.lean\#L296-L298}{{\leanfont Filling.\allowbreak{}submoduleIndex\_smul\_intVectors}}}
\expandafter\newcommand\csname LN@submoduleIndexsmulnezero\endcsname{\href{https://github.com/lsernau/grh-lean/blob/72fead102f663d5c7a407964a13dc0cf9d977207/grh-compact/src/New/Filling/BoxPartition.lean\#L300-L302}{{\leanfont Filling.\allowbreak{}submoduleIndex\_smul\_ne\_zero}}}
\expandafter\newcommand\csname LN@IndexBelow\endcsname{\href{https://github.com/lsernau/grh-lean/blob/72fead102f663d5c7a407964a13dc0cf9d977207/grh-compact/src/New/Filling/BoxPartition.lean\#L305-L305}{{\leanfont Filling.\allowbreak{}IndexBelow}}}
\expandafter\newcommand\csname LN@xRows\endcsname{\href{https://github.com/lsernau/grh-lean/blob/72fead102f663d5c7a407964a13dc0cf9d977207/grh-compact/src/New/Filling/BoxPartition.lean\#L313-L314}{{\leanfont Filling.\allowbreak{}xRows}}}
\expandafter\newcommand\csname LN@paddedLeading\endcsname{\href{https://github.com/lsernau/grh-lean/blob/72fead102f663d5c7a407964a13dc0cf9d977207/grh-compact/src/New/Filling/BoxPartition.lean\#L318-L318}{{\leanfont Filling.\allowbreak{}paddedLeading}}}
\expandafter\newcommand\csname LN@detpaddedLeading\endcsname{\href{https://github.com/lsernau/grh-lean/blob/72fead102f663d5c7a407964a13dc0cf9d977207/grh-compact/src/New/Filling/BoxPartition.lean\#L320-L323}{{\leanfont Filling.\allowbreak{}det\_paddedLeading}}}
\expandafter\newcommand\csname LN@frameMiddle\endcsname{\href{https://github.com/lsernau/grh-lean/blob/72fead102f663d5c7a407964a13dc0cf9d977207/grh-compact/src/New/Filling/BoxPartition.lean\#L327-L327}{{\leanfont Filling.\allowbreak{}frameMiddle}}}
\expandafter\newcommand\csname LN@latticeDisc\endcsname{\href{https://github.com/lsernau/grh-lean/blob/72fead102f663d5c7a407964a13dc0cf9d977207/grh-compact/src/New/Filling/BoxPartition.lean\#L331-L331}{{\leanfont Filling.\allowbreak{}latticeDisc}}}
\expandafter\newcommand\csname LN@frameValueIndex\endcsname{\href{https://github.com/lsernau/grh-lean/blob/72fead102f663d5c7a407964a13dc0cf9d977207/grh-compact/src/New/Filling/BoxPartition.lean\#L335-L337}{{\leanfont Filling.\allowbreak{}frameValueIndex}}}
\expandafter\newcommand\csname LN@xRowsIndex\endcsname{\href{https://github.com/lsernau/grh-lean/blob/72fead102f663d5c7a407964a13dc0cf9d977207/grh-compact/src/New/Filling/BoxPartition.lean\#L341-L342}{{\leanfont Filling.\allowbreak{}xRowsIndex}}}
\expandafter\newcommand\csname LN@saturationIndex\endcsname{\href{https://github.com/lsernau/grh-lean/blob/72fead102f663d5c7a407964a13dc0cf9d977207/grh-compact/src/New/Filling/BoxPartition.lean\#L348-L349}{{\leanfont Filling.\allowbreak{}saturationIndex}}}
\expandafter\newcommand\csname LN@pairInvariant\endcsname{\href{https://github.com/lsernau/grh-lean/blob/72fead102f663d5c7a407964a13dc0cf9d977207/grh-compact/src/New/Filling/BoxPartition.lean\#L353-L354}{{\leanfont Filling.\allowbreak{}pairInvariant}}}
\expandafter\newcommand\csname LN@termElem\endcsname{\href{https://github.com/lsernau/grh-lean/blob/72fead102f663d5c7a407964a13dc0cf9d977207/grh-compact/src/New/Filling/BoxPartition.lean\#L357-L357}{{\leanfont Filling.\allowbreak{}termElem}}}
\expandafter\newcommand\csname LN@minorNorm\endcsname{\href{https://github.com/lsernau/grh-lean/blob/72fead102f663d5c7a407964a13dc0cf9d977207/grh-compact/src/New/Filling/BoxPartition.lean\#L361-L362}{{\leanfont Filling.\allowbreak{}minorNorm}}}
\expandafter\newcommand\csname LN@classGauge\endcsname{\href{https://github.com/lsernau/grh-lean/blob/72fead102f663d5c7a407964a13dc0cf9d977207/grh-compact/src/New/Filling/BoxPartition.lean\#L366-L367}{{\leanfont Filling.\allowbreak{}classGauge}}}
\expandafter\newcommand\csname LN@IsNearRep\endcsname{\href{https://github.com/lsernau/grh-lean/blob/72fead102f663d5c7a407964a13dc0cf9d977207/grh-compact/src/New/Filling/BoxPartition.lean\#L371-L371}{{\leanfont Filling.\allowbreak{}IsNearRep}}}
\expandafter\newcommand\csname LN@invPowWeight\endcsname{\href{https://github.com/lsernau/grh-lean/blob/72fead102f663d5c7a407964a13dc0cf9d977207/grh-compact/src/New/Filling/BoxPartition.lean\#L375-L375}{{\leanfont Filling.\allowbreak{}invPowWeight}}}
\expandafter\newcommand\csname LN@termWeight\endcsname{\href{https://github.com/lsernau/grh-lean/blob/72fead102f663d5c7a407964a13dc0cf9d977207/grh-compact/src/New/Filling/BoxPartition.lean\#L380-L383}{{\leanfont Filling.\allowbreak{}termWeight}}}
\expandafter\newcommand\csname LN@pairWeight\endcsname{\href{https://github.com/lsernau/grh-lean/blob/72fead102f663d5c7a407964a13dc0cf9d977207/grh-compact/src/New/Filling/BoxPartition.lean\#L387-L389}{{\leanfont Filling.\allowbreak{}pairWeight}}}
\expandafter\newcommand\csname LN@IsLabelledTerm\endcsname{\href{https://github.com/lsernau/grh-lean/blob/72fead102f663d5c7a407964a13dc0cf9d977207/grh-compact/src/New/Filling/BoxPartition.lean\#L393-L395}{{\leanfont Filling.\allowbreak{}IsLabelledTerm}}}
\expandafter\newcommand\csname LN@minorNormpos\endcsname{\href{https://github.com/lsernau/grh-lean/blob/72fead102f663d5c7a407964a13dc0cf9d977207/grh-compact/src/New/Filling/BoxPartition.lean\#L401-L402}{{\leanfont Filling.\allowbreak{}HasLevel.\allowbreak{}minorNorm\_pos}}}
\expandafter\newcommand\csname LN@inv\endcsname{\href{https://github.com/lsernau/grh-lean/blob/72fead102f663d5c7a407964a13dc0cf9d977207/grh-compact/src/New/Filling/BoxPartition.lean\#L404-L405}{{\leanfont Filling.\allowbreak{}HasLevel.\allowbreak{}inv}}}
\expandafter\newcommand\csname LN@paddedLeadingmulinv\endcsname{\href{https://github.com/lsernau/grh-lean/blob/72fead102f663d5c7a407964a13dc0cf9d977207/grh-compact/src/New/Filling/BoxPartition.lean\#L407-L408}{{\leanfont Filling.\allowbreak{}HasLevel.\allowbreak{}paddedLeading\_mul\_inv}}}
\expandafter\newcommand\csname LN@existssharedframevector\endcsname{\href{https://github.com/lsernau/grh-lean/blob/72fead102f663d5c7a407964a13dc0cf9d977207/grh-compact/src/New/Filling/BoxPartition.lean\#L413-L418}{{\leanfont Filling.\allowbreak{}HasLevel.\allowbreak{}exists\_shared\_frame\_vector}}}
\expandafter\newcommand\csname LN@xRowsconjTransposemuleq\endcsname{\href{https://github.com/lsernau/grh-lean/blob/72fead102f663d5c7a407964a13dc0cf9d977207/grh-compact/src/New/Filling/BoxPartition.lean\#L425-L429}{{\leanfont Filling.\allowbreak{}xRows\_conjTranspose\_mul\_eq}}}
\expandafter\newcommand\csname LN@frameVectorsubeqsum\endcsname{\href{https://github.com/lsernau/grh-lean/blob/72fead102f663d5c7a407964a13dc0cf9d977207/grh-compact/src/New/Filling/BoxPartition.lean\#L443-L447}{{\leanfont Filling.\allowbreak{}frameVector\_sub\_eq\_sum}}}
\expandafter\newcommand\csname LN@isBoxPartitionepsilon\endcsname{\href{https://github.com/lsernau/grh-lean/blob/72fead102f663d5c7a407964a13dc0cf9d977207/grh-compact/src/New/Filling/BoxPartition.lean\#L453-L505}{{\leanfont Filling.\allowbreak{}isBoxPartition\_epsilon}}}
\expandafter\newcommand\csname LN@tsumclassCutoffeqone\endcsname{\href{https://github.com/lsernau/grh-lean/blob/72fead102f663d5c7a407964a13dc0cf9d977207/grh-compact/src/New/Filling/BoxPartition.lean\#L507-L508}{{\leanfont Filling.\allowbreak{}tsum\_classCutoff\_eq\_one}}}
\expandafter\newcommand\csname LN@tsumclassCutoffboxle\endcsname{\href{https://github.com/lsernau/grh-lean/blob/72fead102f663d5c7a407964a13dc0cf9d977207/grh-compact/src/New/Filling/BoxPartition.lean\#L513-L556}{{\leanfont Filling.\allowbreak{}tsum\_classCutoff\_box\_le}}}
\expandafter\newcommand\csname LQ@tsumclassCutoffboxle\endcsname{\leanquoteline{0}{lemma tsum\_classCutoff\_box\_le \{c s α y y' x\} \{k γ δ w : α → \_\} (hi : ∀ \{p q\}, k p = k q → δ p = δ}\leanquoteline{2}{q → p = q) (H : ∀ p ∈ w.support, γ p ∈ S.termsOver c s (k p) ∧ (δ p)⁻¹ * γ p ∈ S.classUnip (k p) ∧}\leanquoteline{2}{(S.ratToInf (γ p))⁻¹ * y ∈ S.boxedPoints ∧ w p ≤ S.classCutoff (k p) ((S.ratToInf (δ p))⁻¹ * y') *}\leanquoteline{2}{x) \{V : Submodule ℝ \_\} (hV : let F p v := frameVector ((S.ratToInf (γ p) v)⁻¹ * y v); ∀ᵉ}\leanquoteline{4}{(p ∈ w.support) (q ∈ w.support), F p - F q ∈ V) :}\leanquoteline{2}{∑' p, w p ≤ .ofReal (S.b \textasciicircum{} (6 * 𝔰.n \textasciicircum{} 2 * finrank ℝ V)) * x}}
\expandafter\newcommand\csname LN@existsclassUnipmulmemboxedPoints\endcsname{\href{https://github.com/lsernau/grh-lean/blob/72fead102f663d5c7a407964a13dc0cf9d977207/grh-compact/src/New/Filling/BoxPartition.lean\#L558-L561}{{\leanfont Filling.\allowbreak{}exists\_classUnip\_mul\_mem\_boxedPoints}}}
\expandafter\newcommand\csname LN@measurableclassCutoff\endcsname{\href{https://github.com/lsernau/grh-lean/blob/72fead102f663d5c7a407964a13dc0cf9d977207/grh-compact/src/New/Filling/BoxPartition.lean\#L563-L564}{{\leanfont Filling.\allowbreak{}measurable\_classCutoff}}}
\expandafter\newcommand\csname LN@lintegralclassCutoffequnipCovolume\endcsname{\href{https://github.com/lsernau/grh-lean/blob/72fead102f663d5c7a407964a13dc0cf9d977207/grh-compact/src/New/Filling/UnipotentCovolume.lean\#L24-L107}{{\leanfont Filling.\allowbreak{}lintegral\_classCutoff\_eq\_unipCovolume}}}
\expandafter\newcommand\csname LQ@lintegralclassCutoffequnipCovolume\endcsname{\leanquoteline{0}{lemma Filling.lintegral\_classCutoff\_eq\_unipCovolume \{E\} [Field E] [CMType E] [StablyUnram E] \{𝔰\}}\leanquoteline{2}{(S : Filling E 𝔰) c [S.IsAdmissible] (y : \_ → \_) :}\leanquoteline{2}{∫⁻ u, S.classCutoff c fun v ↦ u v * y v ∂(.pi fun \_ ↦ 𝔰.besselHaar) = .ofReal S.unipCovolume}}
\newcommand{\ConcordanceRows}{\crow{isStablebesselS}{\href{https://github.com/lsernau/grh-lean/blob/72fead102f663d5c7a407964a13dc0cf9d977207/grh-compact/src/New/Filling/AmplitudeSums.lean\#L16-L17}{{\leanfont Filling.\allowbreak{}isStable\_besselS}}}{AmplitudeSums.lean}{16--17}
\crow{measurePreservingpiorbitMeasure}{\href{https://github.com/lsernau/grh-lean/blob/72fead102f663d5c7a407964a13dc0cf9d977207/grh-compact/src/New/Filling/AmplitudeSums.lean\#L25-L28}{{\leanfont Filling.\allowbreak{}measurePreserving\_pi\_orbitMeasure}}}{AmplitudeSums.lean}{25--28}
\crow{tsumintegralapply}{\href{https://github.com/lsernau/grh-lean/blob/72fead102f663d5c7a407964a13dc0cf9d977207/grh-compact/src/New/Filling/AmplitudeSums.lean\#L30-L35}{{\leanfont Filling.\allowbreak{}tsum\_integral\_apply}}}{AmplitudeSums.lean}{30--35}
\crow{besselFnmulofmemclassUnip}{\href{https://github.com/lsernau/grh-lean/blob/72fead102f663d5c7a407964a13dc0cf9d977207/grh-compact/src/New/Filling/AmplitudeSums.lean\#L41-L44}{{\leanfont Filling.\allowbreak{}besselFn\_mul\_of\_mem\_classUnip}}}{AmplitudeSums.lean}{41--44}
\crow{besselKernelbesselUmul}{\href{https://github.com/lsernau/grh-lean/blob/72fead102f663d5c7a407964a13dc0cf9d977207/grh-compact/src/New/Filling/AmplitudeSums.lean\#L46-L49}{{\leanfont Filling.\allowbreak{}besselKernel\_besselU\_mul}}}{AmplitudeSums.lean}{46--49}
\crow{lintegralclassCutoffmulprod}{\href{https://github.com/lsernau/grh-lean/blob/72fead102f663d5c7a407964a13dc0cf9d977207/grh-compact/src/New/Filling/AmplitudeSums.lean\#L54-L58}{{\leanfont Filling.\allowbreak{}lintegral\_classCutoff\_mul\_prod}}}{AmplitudeSums.lean}{54--58}
\crow{amplitudeAbslttopandcutle}{\href{https://github.com/lsernau/grh-lean/blob/72fead102f663d5c7a407964a13dc0cf9d977207/grh-compact/src/New/Filling/AmplitudeSums.lean\#L64-L106}{{\leanfont Filling.\allowbreak{}amplitudeAbs\_lt\_top\_and\_cut\_le}}}{AmplitudeSums.lean}{64--106}
\crow{summableamplitude}{\href{https://github.com/lsernau/grh-lean/blob/72fead102f663d5c7a407964a13dc0cf9d977207/grh-compact/src/New/Filling/AmplitudeSums.lean\#L110-L111}{{\leanfont Filling.\allowbreak{}summable\_amplitude}}}{AmplitudeSums.lean}{110--111}
\crow{hasSumamplitudeSumchangeofweight}{\href{https://github.com/lsernau/grh-lean/blob/72fead102f663d5c7a407964a13dc0cf9d977207/grh-compact/src/New/Filling/AmplitudeSums.lean\#L119-L158}{{\leanfont Filling.\allowbreak{}hasSum\_amplitudeSum\_change\_of\_weight}}}{AmplitudeSums.lean}{119--158}
\crow{amplitudeSummove}{\href{https://github.com/lsernau/grh-lean/blob/72fead102f663d5c7a407964a13dc0cf9d977207/grh-compact/src/New/Filling/AmplitudeSums.lean\#L163-L179}{{\leanfont Filling.\allowbreak{}amplitudeSum\_move}}}{AmplitudeSums.lean}{163--179}
\crow{amplitudeSumbesselUmul}{\href{https://github.com/lsernau/grh-lean/blob/72fead102f663d5c7a407964a13dc0cf9d977207/grh-compact/src/New/Filling/AmplitudeSums.lean\#L181-L187}{{\leanfont Filling.\allowbreak{}amplitudeSum\_besselU\_mul}}}{AmplitudeSums.lean}{181--187}
\crow{amplitudeSumeqzeroofnotagrees}{\href{https://github.com/lsernau/grh-lean/blob/72fead102f663d5c7a407964a13dc0cf9d977207/grh-compact/src/New/Filling/AmplitudeSums.lean\#L192-L194}{{\leanfont Filling.\allowbreak{}amplitudeSum\_eq\_zero\_of\_not\_agrees}}}{AmplitudeSums.lean}{192--194}
\crow{amplitudeSumbesselSone}{\href{https://github.com/lsernau/grh-lean/blob/72fead102f663d5c7a407964a13dc0cf9d977207/grh-compact/src/New/Filling/AmplitudeSums.lean\#L201-L220}{{\leanfont Filling.\allowbreak{}amplitudeSum\_besselS\_one}}}{AmplitudeSums.lean}{201--220}
\crow{liouvillegram}{\href{https://github.com/lsernau/grh-lean/blob/72fead102f663d5c7a407964a13dc0cf9d977207/grh-compact/src/New/Filling/BoxPartition.lean\#L45-L55}{{\leanfont Filling.\allowbreak{}liouville\_gram}}}{BoxPartition.lean}{45--55}
\crow{linkMetric}{\href{https://github.com/lsernau/grh-lean/blob/72fead102f663d5c7a407964a13dc0cf9d977207/grh-compact/src/New/Filling/BoxPartition.lean\#L59-L60}{{\leanfont Filling.\allowbreak{}linkMetric}}}{BoxPartition.lean}{59--60}
\crow{linkDiscSq}{\href{https://github.com/lsernau/grh-lean/blob/72fead102f663d5c7a407964a13dc0cf9d977207/grh-compact/src/New/Filling/BoxPartition.lean\#L64-L65}{{\leanfont Filling.\allowbreak{}linkDiscSq}}}{BoxPartition.lean}{64--65}
\crow{framePairlasteqzFormPairmul}{\href{https://github.com/lsernau/grh-lean/blob/72fead102f663d5c7a407964a13dc0cf9d977207/grh-compact/src/New/Filling/BoxPartition.lean\#L72-L74}{{\leanfont Filling.\allowbreak{}framePair\_last\_eq\_zFormPair\_mul}}}{BoxPartition.lean}{72--74}
\crow{classCoordeqfromBlocks}{\href{https://github.com/lsernau/grh-lean/blob/72fead102f663d5c7a407964a13dc0cf9d977207/grh-compact/src/New/Filling/BoxPartition.lean\#L76-L78}{{\leanfont Filling.\allowbreak{}classCoord\_eq\_fromBlocks}}}{BoxPartition.lean}{76--78}
\crow{HasIntegralLifts}{\href{https://github.com/lsernau/grh-lean/blob/72fead102f663d5c7a407964a13dc0cf9d977207/grh-compact/src/New/Filling/BoxPartition.lean\#L86-L88}{{\leanfont Filling.\allowbreak{}HasIntegralLifts}}}{BoxPartition.lean}{86--88}
\crow{finranklink}{\href{https://github.com/lsernau/grh-lean/blob/72fead102f663d5c7a407964a13dc0cf9d977207/grh-compact/src/New/Filling/BoxPartition.lean\#L90-L96}{{\leanfont Filling.\allowbreak{}finrank\_link}}}{BoxPartition.lean}{90--96}
\crow{gramclassCoordX}{\href{https://github.com/lsernau/grh-lean/blob/72fead102f663d5c7a407964a13dc0cf9d977207/grh-compact/src/New/Filling/BoxPartition.lean\#L98-L100}{{\leanfont Filling.\allowbreak{}gram\_classCoordX₀}}}{BoxPartition.lean}{98--100}
\crow{classCoordgram}{\href{https://github.com/lsernau/grh-lean/blob/72fead102f663d5c7a407964a13dc0cf9d977207/grh-compact/src/New/Filling/BoxPartition.lean\#L102-L104}{{\leanfont Filling.\allowbreak{}classCoord\_gram}}}{BoxPartition.lean}{102--104}
\crow{linkMetricposDef}{\href{https://github.com/lsernau/grh-lean/blob/72fead102f663d5c7a407964a13dc0cf9d977207/grh-compact/src/New/Filling/BoxPartition.lean\#L106-L108}{{\leanfont Filling.\allowbreak{}linkMetric\_posDef}}}{BoxPartition.lean}{106--108}
\crow{gramlinkFormeqzerooflelink}{\href{https://github.com/lsernau/grh-lean/blob/72fead102f663d5c7a407964a13dc0cf9d977207/grh-compact/src/New/Filling/BoxPartition.lean\#L110-L113}{{\leanfont Filling.\allowbreak{}gram\_linkForm\_eq\_zero\_of\_le\_link}}}{BoxPartition.lean}{110--113}
\crow{detgramlinkMetricnezero}{\href{https://github.com/lsernau/grh-lean/blob/72fead102f663d5c7a407964a13dc0cf9d977207/grh-compact/src/New/Filling/BoxPartition.lean\#L115-L118}{{\leanfont Filling.\allowbreak{}det\_gram\_linkMetric\_ne\_zero}}}{BoxPartition.lean}{115--118}
\crow{cardeqfinrankcolSpan}{\href{https://github.com/lsernau/grh-lean/blob/72fead102f663d5c7a407964a13dc0cf9d977207/grh-compact/src/New/Filling/BoxPartition.lean\#L121-L123}{{\leanfont Filling.\allowbreak{}card\_eq\_finrank\_colSpan}}}{BoxPartition.lean}{121--123}
\crow{gramlinkMatrixmul}{\href{https://github.com/lsernau/grh-lean/blob/72fead102f663d5c7a407964a13dc0cf9d977207/grh-compact/src/New/Filling/BoxPartition.lean\#L125-L134}{{\leanfont Filling.\allowbreak{}gram\_linkMatrix\_mul}}}{BoxPartition.lean}{125--134}
\crow{classCoordmulframePair}{\href{https://github.com/lsernau/grh-lean/blob/72fead102f663d5c7a407964a13dc0cf9d977207/grh-compact/src/New/Filling/BoxPartition.lean\#L136-L138}{{\leanfont Filling.\allowbreak{}classCoord\_mul\_framePair}}}{BoxPartition.lean}{136--138}
\crow{classCoordmulisotropicFramePair}{\href{https://github.com/lsernau/grh-lean/blob/72fead102f663d5c7a407964a13dc0cf9d977207/grh-compact/src/New/Filling/BoxPartition.lean\#L142-L143}{{\leanfont Filling.\allowbreak{}classCoord\_mul\_isotropicFramePair}}}{BoxPartition.lean}{142--143}
\crow{isIntMatrixframePair}{\href{https://github.com/lsernau/grh-lean/blob/72fead102f663d5c7a407964a13dc0cf9d977207/grh-compact/src/New/Filling/BoxPartition.lean\#L145-L146}{{\leanfont Filling.\allowbreak{}isIntMatrix\_framePair}}}{BoxPartition.lean}{145--146}
\crow{legenIndexdetoffixesNormAtLeast}{\href{https://github.com/lsernau/grh-lean/blob/72fead102f663d5c7a407964a13dc0cf9d977207/grh-compact/src/New/Filling/BoxPartition.lean\#L158-L212}{{\leanfont Filling.\allowbreak{}le\_genIndex\_det\_of\_fixesNormAtLeast}}}{BoxPartition.lean}{158--212}
\crow{existsmemclassUnip}{\href{https://github.com/lsernau/grh-lean/blob/72fead102f663d5c7a407964a13dc0cf9d977207/grh-compact/src/New/Filling/BoxPartition.lean\#L227-L250}{{\leanfont Filling.\allowbreak{}exists\_mem\_classUnip}}}{BoxPartition.lean}{227--250}
\crow{existseichlerclassUnip}{\href{https://github.com/lsernau/grh-lean/blob/72fead102f663d5c7a407964a13dc0cf9d977207/grh-compact/src/New/Filling/BoxPartition.lean\#L252-L263}{{\leanfont Filling.\allowbreak{}exists\_eichler\_classUnip}}}{BoxPartition.lean}{252--263}
\crow{tsumleofencardrangele}{\href{https://github.com/lsernau/grh-lean/blob/72fead102f663d5c7a407964a13dc0cf9d977207/grh-compact/src/New/Filling/BoxPartition.lean\#L270-L274}{{\leanfont Filling.\allowbreak{}tsum\_le\_of\_encard\_range\_le}}}{BoxPartition.lean}{270--274}
\crow{tsumleofcolouring}{\href{https://github.com/lsernau/grh-lean/blob/72fead102f663d5c7a407964a13dc0cf9d977207/grh-compact/src/New/Filling/BoxPartition.lean\#L276-L282}{{\leanfont Filling.\allowbreak{}tsum\_le\_of\_colouring}}}{BoxPartition.lean}{276--282}
\crow{encardintegralle}{\href{https://github.com/lsernau/grh-lean/blob/72fead102f663d5c7a407964a13dc0cf9d977207/grh-compact/src/New/Filling/BoxPartition.lean\#L287-L292}{{\leanfont Filling.\allowbreak{}encard\_integral\_le}}}{BoxPartition.lean}{287--292}
\crow{submoduleIndexsmulintVectors}{\href{https://github.com/lsernau/grh-lean/blob/72fead102f663d5c7a407964a13dc0cf9d977207/grh-compact/src/New/Filling/BoxPartition.lean\#L296-L298}{{\leanfont Filling.\allowbreak{}submoduleIndex\_smul\_intVectors}}}{BoxPartition.lean}{296--298}
\crow{submoduleIndexsmulnezero}{\href{https://github.com/lsernau/grh-lean/blob/72fead102f663d5c7a407964a13dc0cf9d977207/grh-compact/src/New/Filling/BoxPartition.lean\#L300-L302}{{\leanfont Filling.\allowbreak{}submoduleIndex\_smul\_ne\_zero}}}{BoxPartition.lean}{300--302}
\crow{IndexBelow}{\href{https://github.com/lsernau/grh-lean/blob/72fead102f663d5c7a407964a13dc0cf9d977207/grh-compact/src/New/Filling/BoxPartition.lean\#L305-L305}{{\leanfont Filling.\allowbreak{}IndexBelow}}}{BoxPartition.lean}{305--305}
\crow{xRows}{\href{https://github.com/lsernau/grh-lean/blob/72fead102f663d5c7a407964a13dc0cf9d977207/grh-compact/src/New/Filling/BoxPartition.lean\#L313-L314}{{\leanfont Filling.\allowbreak{}xRows}}}{BoxPartition.lean}{313--314}
\crow{paddedLeading}{\href{https://github.com/lsernau/grh-lean/blob/72fead102f663d5c7a407964a13dc0cf9d977207/grh-compact/src/New/Filling/BoxPartition.lean\#L318-L318}{{\leanfont Filling.\allowbreak{}paddedLeading}}}{BoxPartition.lean}{318--318}
\crow{detpaddedLeading}{\href{https://github.com/lsernau/grh-lean/blob/72fead102f663d5c7a407964a13dc0cf9d977207/grh-compact/src/New/Filling/BoxPartition.lean\#L320-L323}{{\leanfont Filling.\allowbreak{}det\_paddedLeading}}}{BoxPartition.lean}{320--323}
\crow{frameMiddle}{\href{https://github.com/lsernau/grh-lean/blob/72fead102f663d5c7a407964a13dc0cf9d977207/grh-compact/src/New/Filling/BoxPartition.lean\#L327-L327}{{\leanfont Filling.\allowbreak{}frameMiddle}}}{BoxPartition.lean}{327--327}
\crow{latticeDisc}{\href{https://github.com/lsernau/grh-lean/blob/72fead102f663d5c7a407964a13dc0cf9d977207/grh-compact/src/New/Filling/BoxPartition.lean\#L331-L331}{{\leanfont Filling.\allowbreak{}latticeDisc}}}{BoxPartition.lean}{331--331}
\crow{frameValueIndex}{\href{https://github.com/lsernau/grh-lean/blob/72fead102f663d5c7a407964a13dc0cf9d977207/grh-compact/src/New/Filling/BoxPartition.lean\#L335-L337}{{\leanfont Filling.\allowbreak{}frameValueIndex}}}{BoxPartition.lean}{335--337}
\crow{xRowsIndex}{\href{https://github.com/lsernau/grh-lean/blob/72fead102f663d5c7a407964a13dc0cf9d977207/grh-compact/src/New/Filling/BoxPartition.lean\#L341-L342}{{\leanfont Filling.\allowbreak{}xRowsIndex}}}{BoxPartition.lean}{341--342}
\crow{saturationIndex}{\href{https://github.com/lsernau/grh-lean/blob/72fead102f663d5c7a407964a13dc0cf9d977207/grh-compact/src/New/Filling/BoxPartition.lean\#L348-L349}{{\leanfont Filling.\allowbreak{}saturationIndex}}}{BoxPartition.lean}{348--349}
\crow{pairInvariant}{\href{https://github.com/lsernau/grh-lean/blob/72fead102f663d5c7a407964a13dc0cf9d977207/grh-compact/src/New/Filling/BoxPartition.lean\#L353-L354}{{\leanfont Filling.\allowbreak{}pairInvariant}}}{BoxPartition.lean}{353--354}
\crow{termElem}{\href{https://github.com/lsernau/grh-lean/blob/72fead102f663d5c7a407964a13dc0cf9d977207/grh-compact/src/New/Filling/BoxPartition.lean\#L357-L357}{{\leanfont Filling.\allowbreak{}termElem}}}{BoxPartition.lean}{357--357}
\crow{minorNorm}{\href{https://github.com/lsernau/grh-lean/blob/72fead102f663d5c7a407964a13dc0cf9d977207/grh-compact/src/New/Filling/BoxPartition.lean\#L361-L362}{{\leanfont Filling.\allowbreak{}minorNorm}}}{BoxPartition.lean}{361--362}
\crow{classGauge}{\href{https://github.com/lsernau/grh-lean/blob/72fead102f663d5c7a407964a13dc0cf9d977207/grh-compact/src/New/Filling/BoxPartition.lean\#L366-L367}{{\leanfont Filling.\allowbreak{}classGauge}}}{BoxPartition.lean}{366--367}
\crow{IsNearRep}{\href{https://github.com/lsernau/grh-lean/blob/72fead102f663d5c7a407964a13dc0cf9d977207/grh-compact/src/New/Filling/BoxPartition.lean\#L371-L371}{{\leanfont Filling.\allowbreak{}IsNearRep}}}{BoxPartition.lean}{371--371}
\crow{invPowWeight}{\href{https://github.com/lsernau/grh-lean/blob/72fead102f663d5c7a407964a13dc0cf9d977207/grh-compact/src/New/Filling/BoxPartition.lean\#L375-L375}{{\leanfont Filling.\allowbreak{}invPowWeight}}}{BoxPartition.lean}{375--375}
\crow{termWeight}{\href{https://github.com/lsernau/grh-lean/blob/72fead102f663d5c7a407964a13dc0cf9d977207/grh-compact/src/New/Filling/BoxPartition.lean\#L380-L383}{{\leanfont Filling.\allowbreak{}termWeight}}}{BoxPartition.lean}{380--383}
\crow{pairWeight}{\href{https://github.com/lsernau/grh-lean/blob/72fead102f663d5c7a407964a13dc0cf9d977207/grh-compact/src/New/Filling/BoxPartition.lean\#L387-L389}{{\leanfont Filling.\allowbreak{}pairWeight}}}{BoxPartition.lean}{387--389}
\crow{IsLabelledTerm}{\href{https://github.com/lsernau/grh-lean/blob/72fead102f663d5c7a407964a13dc0cf9d977207/grh-compact/src/New/Filling/BoxPartition.lean\#L393-L395}{{\leanfont Filling.\allowbreak{}IsLabelledTerm}}}{BoxPartition.lean}{393--395}
\crow{minorNormpos}{\href{https://github.com/lsernau/grh-lean/blob/72fead102f663d5c7a407964a13dc0cf9d977207/grh-compact/src/New/Filling/BoxPartition.lean\#L401-L402}{{\leanfont Filling.\allowbreak{}HasLevel.\allowbreak{}minorNorm\_pos}}}{BoxPartition.lean}{401--402}
\crow{inv}{\href{https://github.com/lsernau/grh-lean/blob/72fead102f663d5c7a407964a13dc0cf9d977207/grh-compact/src/New/Filling/BoxPartition.lean\#L404-L405}{{\leanfont Filling.\allowbreak{}HasLevel.\allowbreak{}inv}}}{BoxPartition.lean}{404--405}
\crow{paddedLeadingmulinv}{\href{https://github.com/lsernau/grh-lean/blob/72fead102f663d5c7a407964a13dc0cf9d977207/grh-compact/src/New/Filling/BoxPartition.lean\#L407-L408}{{\leanfont Filling.\allowbreak{}HasLevel.\allowbreak{}paddedLeading\_mul\_inv}}}{BoxPartition.lean}{407--408}
\crow{existssharedframevector}{\href{https://github.com/lsernau/grh-lean/blob/72fead102f663d5c7a407964a13dc0cf9d977207/grh-compact/src/New/Filling/BoxPartition.lean\#L413-L418}{{\leanfont Filling.\allowbreak{}HasLevel.\allowbreak{}exists\_shared\_frame\_vector}}}{BoxPartition.lean}{413--418}
\crow{xRowsconjTransposemuleq}{\href{https://github.com/lsernau/grh-lean/blob/72fead102f663d5c7a407964a13dc0cf9d977207/grh-compact/src/New/Filling/BoxPartition.lean\#L425-L429}{{\leanfont Filling.\allowbreak{}xRows\_conjTranspose\_mul\_eq}}}{BoxPartition.lean}{425--429}
\crow{frameVectorsubeqsum}{\href{https://github.com/lsernau/grh-lean/blob/72fead102f663d5c7a407964a13dc0cf9d977207/grh-compact/src/New/Filling/BoxPartition.lean\#L443-L447}{{\leanfont Filling.\allowbreak{}frameVector\_sub\_eq\_sum}}}{BoxPartition.lean}{443--447}
\crow{isBoxPartitionepsilon}{\href{https://github.com/lsernau/grh-lean/blob/72fead102f663d5c7a407964a13dc0cf9d977207/grh-compact/src/New/Filling/BoxPartition.lean\#L453-L505}{{\leanfont Filling.\allowbreak{}isBoxPartition\_epsilon}}}{BoxPartition.lean}{453--505}
\crow{tsumclassCutoffeqone}{\href{https://github.com/lsernau/grh-lean/blob/72fead102f663d5c7a407964a13dc0cf9d977207/grh-compact/src/New/Filling/BoxPartition.lean\#L507-L508}{{\leanfont Filling.\allowbreak{}tsum\_classCutoff\_eq\_one}}}{BoxPartition.lean}{507--508}
\crow{tsumclassCutoffboxle}{\href{https://github.com/lsernau/grh-lean/blob/72fead102f663d5c7a407964a13dc0cf9d977207/grh-compact/src/New/Filling/BoxPartition.lean\#L513-L556}{{\leanfont Filling.\allowbreak{}tsum\_classCutoff\_box\_le}}}{BoxPartition.lean}{513--556}
\crow{existsclassUnipmulmemboxedPoints}{\href{https://github.com/lsernau/grh-lean/blob/72fead102f663d5c7a407964a13dc0cf9d977207/grh-compact/src/New/Filling/BoxPartition.lean\#L558-L561}{{\leanfont Filling.\allowbreak{}exists\_classUnip\_mul\_mem\_boxedPoints}}}{BoxPartition.lean}{558--561}
\crow{measurableclassCutoff}{\href{https://github.com/lsernau/grh-lean/blob/72fead102f663d5c7a407964a13dc0cf9d977207/grh-compact/src/New/Filling/BoxPartition.lean\#L563-L564}{{\leanfont Filling.\allowbreak{}measurable\_classCutoff}}}{BoxPartition.lean}{563--564}
\crow{lintegralclassCutoffequnipCovolume}{\href{https://github.com/lsernau/grh-lean/blob/72fead102f663d5c7a407964a13dc0cf9d977207/grh-compact/src/New/Filling/UnipotentCovolume.lean\#L24-L107}{{\leanfont Filling.\allowbreak{}lintegral\_classCutoff\_eq\_unipCovolume}}}{UnipotentCovolume.lean}{24--107}}
% the notations of the defined functions: one slot for each explicit argument of the Lean declaration
\newcommand{\lmet}[1]{Q^{\mathrm{link}}_{#1}}
\newcommand{\ldsq}[3]{\mathrm{ld}_{#1}(#2,#3)}
\newcommand{\IdxB}[1]{\mathrm{Ib}(#1)}
\newcommand{\xrows}[1]{#1^{\mathrm{x}}}
\newcommand{\padl}[2]{\operatorname{pad}_{#1}(#2)}
\newcommand{\fmid}[1]{(#1\mathfrak{f})^{\mathrm{mid}}}
\newcommand{\latd}[2]{\operatorname{disc}_{#1,#2}}
\newcommand{\fvi}[5]{\mathrm{o}_{#1,#2}(#3;#4,#5)}
\newcommand{\xri}[3]{\mathrm{i}_{#1,#2}(#3)}
\newcommand{\sati}[2]{\mathfrak{d}_{#1}(#2)}
\newcommand{\pinv}[5]{\mathrm{pi}_{#1}(#2,#3;#4,#5)}
\newcommand{\telem}[3]{\mathrm{el}_{#1}(#2;#3)}
\newcommand{\mnorm}[3]{\mathrm{mn}_{#1}(#2;#3)}
\newcommand{\cgauge}[5]{\mathrm{cg}_{#1}(#2;#3,#4;#5)}
\newcommand{\IsNRep}[4]{#4\in\operatorname{NRep}_{#1}(#2,#3)}
\newcommand{\IsLab}[5]{#5\in\operatorname{Lab}_{#1}(#2;#3,#4)}
\newcommand{\UnitMembers}{497}
\newcommand{\UnitInNote}{68}
\newcommand{\UnitProofs}{179}
\newcommand{\UnitSteps}{1371}
\newcommand{\UnitAllSteps}{1479}

% generated by notes/filling-setup/shared/mkspec.py from shared/notation-148ab.json; do not edit.
% \input after data.tex: a macro that the build made from the spec is kept, the others are defined here (same TeX).
\providecommand{\fs}{\mathfrak{s}}                      % the shape (rho, kappa, k_0)
\providecommand{\sT}{\mathsf{T}}                        % a finite set of finite places
\providecommand{\cL}{\mathcal{L}}                       % a level
\providecommand{\AEf}{\mathbb{A}_{E,f}}
\providecommand{\AdE}{\mathbb{A}_{E}}
\providecommand{\hO}{\hat{O}_E}
\providecommand{\GL}{\mathrm{GL}}
\providecommand{\full}{\mathrm{full}}
\providecommand{\cut}{\mathrm{cut}}
\providecommand{\rh}[1]{\varrho(#1)}                    % Shape.ϱ
\providecommand{\kk}[1]{k(#1)}                          % Shape.k = #κ
\providecommand{\nn}[1]{n(#1)}                          % Shape.n = ϱ + 1
\providecommand{\rr}[1]{r(#1)}                          % Shape.r = 2ϱ + k + 1
\providecommand{\Ix}[1]{\operatorname{Ind}(#1)}                          % Shape.BesselIndex
\providecommand{\Mx}[1]{\operatorname{Mid}(#1)}                          % Shape.MiddleIndex
\providecommand{\ninf}[1]{f(#1)}                        % numInfPlaces E
\providecommand{\ET}[1]{E_{#1}}                         % LocProd T
\providecommand{\bS}[1]{b(#1)}                          % S.b
\providecommand{\NS}[1]{N(#1)}                          % S.N
\providecommand{\qS}[2]{q_{#2}(#1)}                     % S.levelPrime i
\providecommand{\piS}[1]{\bar\pi(#1)}                   % S.piBar
\providecommand{\vpS}[1]{\varpi(#1)}                    % S.ϖ
\providecommand{\ChS}[2]{\mathrm{Ch}_{#2}(#1)}          % S.chosenPlaces i
\providecommand{\KcolS}[1]{\mathsf{K}^{\mathrm{col}}(#1)}        % S.columnLevel
\providecommand{\BsysS}[2]{\mathfrak{B}_{#2}(#1)}       % S.besselSys v
\providecommand{\mass}[2]{\operatorname{mass}(#1,#2)}   % mass E r
% Lib: forms, groups and frames (letters of the notes of unit 132)
\providecommand{\fr}[1]{\mathfrak{f}_{#1}}              % the frame vector of index a (e_a for a < rho, z for a = rho)
\providecommand{\Trg}[2]{T^{\mathrm{rat}}(#1;#2)}       % Shape.middleGram N: the form diag(1,...,1,-N)
\providecommand{\JT}[1]{J(#1)}                          % Shape.Form.splitGram
\providecommand{\GT}[1]{G(#1)}                          % Shape.Form.G
\providecommand{\Gone}[1]{G_{1}(#1)}                    % Shape.G₁
\providecommand{\UT}[1]{\mathcal{U}(#1)}                % besselU
\providecommand{\UoT}[1]{\mathcal{U}^{0}(#1)}           % besselU0
\providecommand{\HT}[1]{\mathcal{H}(#1)}                % besselH
\providecommand{\ST}[1]{\mathcal{S}(#1)}                % besselS
\providecommand{\Hl}[2]{\mathcal{H}_{#1}(#2)}           % GT l
\providecommand{\Ql}[2]{\mathcal{Q}_{#1}(#2)}           % unitriGroup l
\providecommand{\sd}[1]{\operatorname{sd}(#1)}          % simpleRootSum
\providecommand{\ZT}[2]{Z_{#1}(#2)}                     % Shape.Form.frameGram
\providecommand{\Zr}[1]{Z(#1)}                          % realFrameGram
\providecommand{\lmin}[2]{\mathrm{D}_{#1}(#2)}          % leadingMinor Z l
\providecommand{\weyl}[2]{w_{#1}(#2)}                   % weylElem l p
\providecommand{\gix}[1]{[#1]}                          % genIndex (the generalized index)
\providecommand{\den}[1]{\operatorname{den}(#1)}        % denomIndex
\providecommand{\Kf}[2]{\mathsf{K}_{f}(#1,#2)}                   % Kf E n
\providecommand{\Gf}[2]{G_{f}(#1,#2)}                   % Gf E n
\providecommand{\volO}[1]{\operatorname{vol}(\hat O_{#1})} % volÔ E
\providecommand{\rdE}[1]{\operatorname{rd}_{#1}}        % rootDiscr E
\providecommand{\dE}[1]{d_{#1}}                         % discr E
\providecommand{\CC}{\mathbb{C}}
\providecommand{\RR}{\mathbb{R}}
\providecommand{\NN}{\mathbb{N}}
\providecommand{\vol}{\operatorname{vol}}
% ----- the definitions of chunks a and b of unit 148 (and, for this note, its own additions)
\providecommand{\Tlatg}[2]{T^{\mathrm{lat}}(#1;#2)}                  % latMiddleGram (𝔰, N)
\providecommand{\mbc}[2]{\Lambda^{\circ}(#1;#2)}                     % middleBaseChange (𝔰, p)
\providecommand{\blb}[2]{\Lambda(#1;#2)}                             % besselLatticeBasis (𝔰, p)
\providecommand{\fdisc}[2]{\mathsf{g}(#1;#2)}                        % frameDisc (N, Z)
\providecommand{\fpair}[1]{[\mathfrak{f}\mid #1\mathfrak{f}]}        % framePair (γ)
\providecommand{\ifpair}[1]{[\mathfrak{e}\mid #1\mathfrak{e}]}       % isotropicFramePair (γ)
\providecommand{\lmat}[1]{\mathrm{L}(#1)}                            % linkMatrix (γ)
\providecommand{\link}[1]{\operatorname{link}(#1)}                   % link (γ)
\providecommand{\damp}[2]{\operatorname{dp}_{#1}(#2)}                % dampedPower (k, r)
\providecommand{\hw}[3]{\Lambda^{\mathbb{C}}_{#1}(#2;#3)}            % headWeight (K, i, p)
\providecommand{\Sieve}[8]{\operatorname{Sieve}(#1,#2,#3,#4,#5,#6,#7,#8)} % SieveSystem (u, ι, n, J, ψ, En, Λ, β)
\providecommand{\fhw}[3]{\Lambda_{#1}(#2;#3)}                        % finHeadWeight (K, i, p)
\providecommand{\occ}[1]{c_{G}(#1)}                                  % Shape.openCellConst (𝔰)
\providecommand{\bandc}[1]{c_{\mathrm{band}}(#1)}                    % Shape.bandConst (𝔰)
\providecommand{\bsize}[1]{\operatorname{bs}(#1)}                    % besselSize (y)
\providecommand{\InCut}[2]{#1\in\operatorname{Cut}(#2)}              % InCut (y, b)
\providecommand{\frbox}[1]{\operatorname{FBox}(#1)}                  % frameBox (ρ)
\providecommand{\gdr}[1]{\mathsf{g}^{\mathbb{R}}(#1)}                % gramDiscReal (Z)
\providecommand{\cellm}[4]{\operatorname{cm}(#1;#2,#3,#4)}           % cellMass (y, deg, ρ, d)
\providecommand{\bcorr}[3]{\mathrm{J}(#1,#2;#3)}                     % besselCorrelation (μ, W, γ)
\providecommand{\SchurLe}[2]{\operatorname{Sch}(#1)\le #2}           % SchurLe (W, D)
\providecommand{\HasSB}[4]{\operatorname{SchB}(#1,#2;#3,#4)}         % HasSchurBound (W, rep, B, D)
\providecommand{\ncorr}[3]{\mathrm{J}^{\circ}(#1;#2,#3)}             % normCorrelation (y, W, p)
\providecommand{\BandL}[4]{\operatorname{BL}(#1,#2,#3,#4)}           % HasBandLimits (b, W, deg, J)
\providecommand{\CutPair}[4]{\operatorname{CP}(#1,#2,#3,#4)}         % IsCutPair (b, f, rep, W)
\providecommand{\BSys}[3]{\operatorname{BSys}(#1,#2,#3)}             % BesselSystem (𝔰, b, f)
\providecommand{\Fill}[2]{\operatorname{Fill}(#1,#2)}                % Filling (E, 𝔰)
\providecommand{\gS}[1]{g(#1)}                                       % Filling.g (S)
\providecommand{\Trat}[2]{T^{\mathrm{rat}}_{#1}(#2)}                 % Filling.ratMiddle (S, R)
\providecommand{\Tlat}[2]{T^{\mathrm{lat}}_{#1}(#2)}                 % Filling.latMiddle (S, R)
\providecommand{\GS}[2]{G_{#1}(#2)}                                  % Filling.G (S, R)
\providecommand{\GFin}[1]{G^{\mathrm{lat}}_{#1}}                     % Filling.GFin (S)
\providecommand{\GRes}[1]{\mathsf{G}(#1)}                            % Filling.GRes (S)
\providecommand{\Bc}[1]{\mathrm{Bc}(#1)}                             % Filling.BesselCoord (S)
\providecommand{\Cl}[1]{\mathrm{Cl}(#1)}                             % Filling.Classes (S)
\providecommand{\Zf}[2]{Z_{#1}(#2)}                                  % Filling.frameProducts (S, γ)
\providecommand{\KS}[1]{\mathsf{K}(#1)}                              % Filling.latStabilizer (S)
\providecommand{\redS}[1]{\mathrm{red}_{#1}}                         % Filling.reduction (S)
\providecommand{\KgS}[1]{\mathsf{K}_g(#1)}                           % Filling.principalLevel (S)
\providecommand{\redim}[2]{\overline{#2}^{\,#1}}                     % Filling.reducedImage (S, A)
\providecommand{\kapS}[2]{\kappa_{#2}(#1)}                           % Filling.kappa (S, v)
\providecommand{\bscale}[3]{\mathrm{sc}_{#1,#2}(#3)}                 % Filling.basisScale (S, v, i)
\providecommand{\sdc}[1]{\psi_{#1}}                                  % Filling.selfDualChar (S)
\providecommand{\fscale}[3]{\mathrm{fs}_{#1}(#2,#3)}                 % Filling.frameScaleSq (S, v, a)
\providecommand{\fsprod}[1]{\mathrm{fs}(#1)}                         % Filling.frameScaleProd (S)
\providecommand{\ucov}[1]{\mathrm{cov}(#1)}                          % Filling.unipCovolume (S)
\providecommand{\siegC}[1]{c_{*}(#1)}                                % Filling.siegelConst (S)
\providecommand{\tdisc}[2]{\mathsf{g}_{#1}(#2)}                      % Filling.termDisc (S, γ)
\providecommand{\bsGL}[2]{\mathrm{Sc}_{#1,#2}}                       % Filling.basisScaleGL (S, v)
\providecommand{\latB}[1]{\Lambda(#1)}                               % Filling.latBasis (S)
\providecommand{\tolb}[1]{\mathrm{lb}_{#1}}                          % Filling.toLatBasis (S)
\providecommand{\radl}[1]{\mathrm{r}^{\mathbb{A}}_{#1}}              % Filling.ratToAdelic (S)
\providecommand{\ainf}[1]{\mathrm{a}^{\infty}_{#1}}                  % Filling.adelicToInf (S)
\providecommand{\afin}[1]{\mathrm{a}^{f}_{#1}}                       % Filling.adelicToFin (S)
\providecommand{\rinf}[1]{\mathrm{r}^{\infty}_{#1}}                  % Filling.ratToInf (S)
\providecommand{\rfin}[1]{\mathrm{r}^{f}_{#1}}                       % Filling.ratToFin (S)
\providecommand{\incH}[1]{\mathrm{inc}_{#1}}                         % Filling.inclHFin (S)
\providecommand{\hc}[2]{h_{#1}(#2)}                                  % Filling.h (S, c)
\providecommand{\pos}[4]{\mathrm{pos}_{#1}(#2,#3;#4)}                % Filling.termPosition (S, c, c', γ)
\providecommand{\terms}[4]{\mathrm{Tm}_{#1}(#2,#3,#4)}               % Filling.termsOver (S, c, s, c')
\providecommand{\clU}[2]{\mathcal{U}_{#1,#2}}                        % Filling.classUnip (S, c)
\providecommand{\clUo}[2]{\mathcal{U}^{0}_{#1,#2}}                   % Filling.classUnip₀ (S, c)
\providecommand{\bxs}[1]{\operatorname{bx}(#1)}                      % Filling.boxSize (S)
\providecommand{\IsBP}[3]{#3\in\operatorname{BP}(#1,#2)}             % Filling.IsBoxPartition (S, c, ω)
\providecommand{\besfn}[4]{W^{#1}_{#2}(#3)_{#4}}                     % Filling.besselFn (σ, S, y, i)
\providecommand{\besker}[4]{\mathcal{W}^{#1}_{#2}(#3;#4)}            % Filling.besselKernel (σ, S, γ, y)
\providecommand{\cdim}[1]{\mathrm{cd}(#1)}                           % Filling.columnDim (S)
\providecommand{\bdeg}[1]{\deg(#1)}                                  % Filling.besselDegree (S)
\providecommand{\mainE}[1]{\mathsf{d}(#1)}                           % Filling.mainEnergy (S)
\providecommand{\uord}[1]{\mathrm{uo}(#1)}                           % Filling.unipOrder (S)
\providecommand{\packC}[1]{a_{*}(#1)}                                % Filling.packetConst (S)
\providecommand{\besE}[1]{\mathsf{d}_{1}(#1)}                        % Filling.besselEnergy (S)
\providecommand{\tailS}[1]{\mathrm{tl}(#1)}                          % Filling.tailScale (S)
\providecommand{\codimF}[1]{\mathrm{cf}(#1)}                         % Filling.codimFactor (S)
\providecommand{\IsAdm}[1]{#1\text{ is admissible}}                  % Filling.IsAdmissible (S)
\providecommand{\cutoff}[2]{\beta_{#1,#2}}                           % Filling.classCutoff (S, c)
\providecommand{\amp}[5]{a^{#1}_{#2}(#3,#4;#5)}                      % Filling.amplitude (σ, S, c, c', γ)
\providecommand{\ampabs}[5]{|a|^{#1}_{#2}(#3,#4;#5)}                 % Filling.amplitudeAbs (σ, S, c, c', γ)
\providecommand{\ampsum}[6]{A^{#1}_{#2}(#3;#4,#5,#6)}                % Filling.amplitudeSum (σ, S, T, c, s, c')
\providecommand{\gid}[1]{\hat{g}(#1)}                                % Filling.gIdele (S)
\providecommand{\uchar}[2]{\psi^{\mathcal{U}}_{#1}(#2)}              % Filling.unipotentChar (S, u)
\providecommand{\Agr}[2]{#2\in\operatorname{Agr}(#1)}                % Filling.Agrees (S, s)
\providecommand{\IsStab}[2]{#2\in\operatorname{Stab}(#1)}            % Filling.IsStable (S, T)
\providecommand{\tlink}[2]{\mathrm{Tm}^{\mathrm{lk}}_{#1}(#2)}       % Filling.termsOfLink (S, Λ)
\providecommand{\clB}[2]{\mathrm{B}_{#1}(#2)}                        % Filling.classBasis (S, c)
\providecommand{\ioia}[2]{\operatorname{adj}_{#1}^{#2}}              % Filling.idOrInvAdjoint (S, b)
\providecommand{\rag}[2]{\mathrm{red}_{#1,#2}}                       % Filling.reduceAboveG (S, x)
\providecommand{\vfix}[3]{\mathrm{Fix}_{#1,#2}(#3)}                  % Filling.vectorFixer (S, w, d)
\providecommand{\FixN}[2]{\operatorname{fix}(#1)\ge #2}              % Filling.FixesNormAtLeast (s, x)
\providecommand{\fixset}[3]{\mathrm{Fix}_{#1}(#2,#3)}                % Filling.fixingSet (S, Sg, d)
\providecommand{\zfp}[1]{\mathrm{L}_{z}(#1)}                         % Filling.zFormPair (T)
\providecommand{\clink}[2]{\operatorname{link}^{\circ}(#1,#2)}       % commonLink (γ, γ')
\providecommand{\HCFV}[4]{\operatorname{CF}_{#4}(#1;#2,#3)}          % HasCommonFlagVector (γ, c, c', a)
\providecommand{\IsNF}[2]{#2\in\operatorname{NF}_{#1}}               % IsNormalForm (j, γ)
\providecommand{\ratG}[1]{J_{#1}}                                    % Filling.ratGram (S)
\providecommand{\lform}[1]{Q^{\pm}_{#1}}                             % Filling.linkForm (S)
\providecommand{\clC}[2]{\mathrm{C}_{#1}(#2)}                        % Filling.classCoord (S, c)
\providecommand{\clCX}[2]{\mathrm{C}^{\circ}_{#1}(#2)}               % Filling.classCoordX₀ (S, c)
\providecommand{\cw}[2]{\operatorname{cw}_{#2}(#1)}                  % Filling.countWeight (S, i)
\providecommand{\HasLev}[3]{#3\in\operatorname{Lv}_{#2}(#1)}         % Filling.HasLevel (S, j, γ)
\providecommand{\Cong}[5]{#4\sim_{#1;#2,#3}#5}                       % Filling.CongruentTerms (S, c, c', γ, γ')
\providecommand{\tdnorm}[2]{\|\mathsf{g}\|_{#1}(#2)}                 % Filling.termDiscNorm (S, γ)
\providecommand{\bxp}[1]{\operatorname{BxP}(#1)}                     % Filling.boxedPoints (S)
\providecommand{\IsNr}[2]{#2\in\operatorname{Near}(#1)}              % Filling.IsNear (S, γ)
\providecommand{\ChA}[2]{#2\in\operatorname{ChA}(#1)}                % Filling.CharsAgree (S, γ)
\providecommand{\fadm}[1]{(#1)_{\mathbb{A}_{f}}}                     % Filling.finAdeleMat (M)
\providecommand{\slb}[2]{\mathrm{m}_{#1}(#2)}                        % Filling.sublatticeBasis (H, V)
\providecommand{\sld}[2]{\operatorname{disc}_{#1}(#2)}               % Filling.sublatticeDisc (u, V)
\providecommand{\lmet}[1]{Q^{\mathrm{link}}_{#1}}                    % Filling.linkMetric (S)
\providecommand{\ldsq}[3]{\mathrm{ld}_{#1}(#2,#3)}                   % Filling.linkDiscSq (S, v, e)
\providecommand{\HIL}[3]{#3\in\operatorname{IL}(#1;#2)}              % Filling.HasIntegralLifts (c, γ, m)
\providecommand{\IdxB}[1]{\mathrm{Ib}(#1)}                           % Filling.IndexBelow (j)
\providecommand{\xrows}[1]{#1^{\mathrm{x}}}                          % Filling.xRows (Z)
\providecommand{\padl}[2]{\operatorname{pad}_{#1}(#2)}               % Filling.paddedLeading (j, Z)
\providecommand{\fmid}[1]{(#1\mathfrak{f})^{\mathrm{mid}}}           % Filling.frameMiddle (γ)
\providecommand{\latd}[2]{\operatorname{disc}_{#1,#2}}               % Filling.latticeDisc (S, c)
\providecommand{\fvi}[5]{\mathrm{o}_{#1,#2}(#3;#4,#5)}               % Filling.frameValueIndex (S, c, k, a, V)
\providecommand{\xri}[3]{\mathrm{i}_{#1,#2}(#3)}                     % Filling.xRowsIndex (S, j, k)
\providecommand{\sati}[2]{\mathfrak{d}_{#1}(#2)}                     % Filling.saturationIndex (S, k)
\providecommand{\pinv}[5]{\mathrm{pi}_{#1}(#2,#3;#4,#5)}             % Filling.pairInvariant (S, c, c', γ, γ')
\providecommand{\telem}[3]{\mathrm{el}_{#1}(#2;#3)}                  % Filling.termElem (S, c, p)
\providecommand{\mnorm}[3]{\mathrm{mn}_{#1}(#2;#3)}                  % Filling.minorNorm (S, j, γ₁)
\providecommand{\cgauge}[5]{\mathrm{cg}_{#1}(#2;#3,#4;#5)}           % Filling.classGauge (S, j, c, c', γ)
\providecommand{\IsNRep}[4]{#4\in\operatorname{NRep}_{#1}(#2,#3)}    % Filling.IsNearRep (S, c, c', γ)
\providecommand{\ipw}[4]{\mathrm{ipw}_{#1}(#2;#3,#4)}                % Filling.invPowWeight (S, c, n, p)
\providecommand{\termw}[7]{\mathrm{w}_{#1}(#2,#3;#4;#5,#6;#7)}       % Filling.termWeight (S, j, i, c, p, p', V)
\providecommand{\pairw}[6]{\mathrm{pw}_{#1}(#2;#3,#4;#5,#6)}         % Filling.pairWeight (S, j, c, c', γ, γ')
\providecommand{\IsLab}[5]{#5\in\operatorname{Lab}_{#1}(#2;#3,#4)}   % Filling.IsLabelledTerm (S, j, c, t, p)
\providecommand{\energy}[3]{\mathrm{En}_{#1}(#2;#3)}                 % Filling.energy (S, A, Ω)
\providecommand{\linkE}[5]{\mathrm{LE}^{#1}_{#2}(#3;#4,#5)}          % Filling.linkEnergy (σ, S, T, Ω', R)
\providecommand{\csize}[3]{\operatorname{cs}_{#1}(#2;#3)}            % Filling.classSize (S, j, x)
\providecommand{\bcut}[2]{\beta^{0}_{#1,#2}}                         % Filling.boxCutoff (S, c)
\providecommand{\schw}[4]{\mathrm{sw}_{#1}(#2;#3,#4)}                % Filling.schurWeight (S, c, γ, y)
\providecommand{\schi}[4]{\mathrm{SI}_{#1}(#2,#3;#4)}                % Filling.schurIntegral (S, c, c', γ')
\providecommand{\charT}[2]{\psi_{#1,#2}}                             % Filling.charOnPlaces (S, T)
\providecommand{\fgain}[3]{\beta_{#1,#2}(#3)}                        % Filling.frequencyGain (S, T, ξ)
\providecommand{\ratpt}[3]{\mathrm{pt}_{#1}(#2;#3)}                  % Filling.ratPoint (T, l, t)
\providecommand{\Fits}[6]{\operatorname{Fits}(#1;#2,#3,#4,#5,#6)}    % Filling.Fits (S, V, l, kw, Y, η)
\providecommand{\stdK}[1]{\mathsf{K}^{\mathrm{std}}_{g}(#1)}         % Filling.stdPrincipalLevel (S)
\providecommand{\tailv}[3]{\operatorname{tv}(#1,#2,#3)}              % tailVector (π, t, x)


\linespread{1.08}                       % 11 pt on 14.7 pt
\setlength{\parindent}{0pt}
\setlength{\parskip}{5.5pt plus 1.5pt minus 0.5pt}
\setlength{\abovedisplayskip}{7.5pt plus 2pt minus 3pt}
\setlength{\belowdisplayskip}{7.5pt plus 2pt minus 3pt}
\setlength{\abovedisplayshortskip}{3pt plus 2pt}
\setlength{\belowdisplayshortskip}{6pt plus 2pt minus 2pt}
\emergencystretch=1.5em
\clubpenalty=10000 \widowpenalty=10000 \displaywidowpenalty=10000
\relpenalty=7000 \binoppenalty=9000
\raggedbottom
\pagestyle{fancy}\fancyhf{}\renewcommand{\headrulewidth}{0pt}
\fancyfoot[C]{\sffamily\fontsize{7.6}{9}\selectfont\color{inkfaint}\thepage\,/\,\pageref*{LastPage}}
\setlist{topsep=2pt,partopsep=0pt,parsep=0pt,itemsep=4pt,leftmargin=2.1em,labelsep=0.55em}

\newcommand{\Q}{\mathbb{Q}}
\newcommand{\Z}{\mathbb{Z}}
\newcommand{\F}{\mathbb{F}}
\newcommand{\len}{\operatorname{len}}
\newcommand{\cok}{\operatorname{cok}}
\newcommand{\End}{\operatorname{End}}
\newcommand{\Hom}{\operatorname{Hom}}
\newcommand{\Gal}{\operatorname{Gal}}
\newcommand{\diag}{\operatorname{diag}}
\newcommand{\kw}[1]{\emph{#1}}
\newcommand{\m}[1]{\mbox{$#1$}}        % a short formula that must not be broken across lines
\newcommand{\head}[1]{\par\addvspace{17pt plus 4pt minus 2pt}\Needspace{6\baselineskip}%
  \noindent{\fontsize{13.2}{16}\selectfont\bfseries\color{ink}#1}\par\nopagebreak\vspace{-1pt}}
\newcommand{\sect}[2]{\refstepcounter{section}\label{#2}\par\addvspace{19pt plus 4pt minus 2pt}\Needspace{7\baselineskip}%
  \noindent\llap{{\fontsize{13.2}{16}\selectfont\bfseries\color{inkfaint}\thesection}\hspace{9pt}}%
  {\fontsize{13.2}{16}\selectfont\bfseries\color{ink}#1}\par\nopagebreak\vspace{-1pt}}
\newcommand{\subhead}[1]{\par\addvspace{9pt plus 2pt}\Needspace{4\baselineskip}\noindent\textbf{#1}\enspace\ignorespaces}
% where a statement lives in the Lean archive
\newcommand{\leantag}{{\sffamily\bfseries\fontsize{6.3}{8}\selectfont\color{leancol}\addfontfeatures{LetterSpace=7}LEAN}}
\newcommand{\sep}{\unskip\nobreak\hspace{5pt}{\color{inkfaint}\textperiodcentered}\hspace{5pt plus 2pt}\ignorespaces}
\newcommand{\leanline}[1]{\par\nopagebreak\vspace{-2.5pt}{\fontsize{8.9}{12.4}\selectfont\color{inksoft}\raggedright
  \noindent\llap{\leantag\hspace{9pt}}#1\par}\vspace{1pt}}
% a displayed statement in a tinted box: #1 colour, #2 text
\newcommand{\statement}[2]{\par\addvspace{7pt}\begingroup\setlength{\fboxsep}{8.5pt}\setlength{\fboxrule}{0.6pt}%
  \noindent\fcolorbox{#1!30!white}{#1!3!white}{\parbox{\dimexpr\linewidth-2\fboxsep-2\fboxrule\relax}{#2}}%
  \endgroup\par\addvspace{7pt}}
\newcommand{\nt}[2]{\par\noindent\hangindent=1.1em\hangafter=1 #1\enspace{\color{inksoft}#2}\par\vspace{1pt}}
\newcommand{\legdot}[1]{\tikz[baseline=-0.62ex]\fill[#1] (0,0) circle (2.5pt);\hspace{3pt}}
\newcommand{\crow}[4]{\csname CS@#1\endcsname & #2 & \textcolor{inksoft}{#3} & \textcolor{inksoft}{#4}\\}
\newcommand{\leanquoteline}[2]{\par\noindent\hangindent=\dimexpr#1\fontdimen2\font+2.2em\relax\hangafter=1
  \hspace*{#1\fontdimen2\font}#2}
\newcommand{\quoted}[1]{\par\vspace{3pt}{\ttfamily\fontsize{7.5}{10.4}\selectfont\color{leancol}\raggedright
  \setlength{\parskip}{0pt}\leanquote{#1}\par}\vspace{4pt}}
\newcommand{\bx}[3]{{\sffamily\bfseries\fontsize{7.3}{9.2}\selectfont\color{#1!82!black}#2}\\[1.5pt]#3}
\newcommand{\cs}[2]{\expandafter\newcommand\csname CS@#1\endcsname{#2}}

% theorems
\newtheoremstyle{plainx}{7pt plus 2pt}{7pt plus 2pt}{\itshape}{}{\bfseries}{.}{0.5em}{}
\newtheoremstyle{uprightx}{7pt plus 2pt}{7pt plus 2pt}{}{}{\bfseries}{.}{0.5em}{}
\theoremstyle{plainx}
\newtheorem{theorem}{Theorem}[section]
\newtheorem{lemma}[theorem]{Lemma}
\newtheorem{proposition}[theorem]{Proposition}
\newtheorem{corollary}[theorem]{Corollary}
\theoremstyle{uprightx}
\newtheorem{definition}[theorem]{Definition}
\newtheorem{setting}[theorem]{Setting}          % hypotheses shared by several statements; each of them cites it by number
\newtheorem{remark}[theorem]{Remark}
\newtheorem{example}[theorem]{Example}
\renewcommand{\qedsymbol}{{\setlength{\fboxsep}{0pt}\setlength{\fboxrule}{0.5pt}\fbox{\rule{0pt}{0.46em}\hspace{0.46em}}}}
% what holds throughout the note: the ambient objects only (implicit and instance arguments), never a hypothesis
\newenvironment{conventions}{\par\addvspace{5pt}\noindent\textbf{Conventions.}\enspace\ignorespaces}{\par\addvspace{3pt}}
\makeatletter
\renewenvironment{proof}[1][Proof]{\par\pushQED{\qed}\normalfont\topsep0pt\partopsep0pt
  \trivlist\item[\hskip\labelsep\itshape #1.]\ignorespaces}{\popQED\endtrivlist\@endpefalse}
\makeatother

% ----- notions of the background library and of the record S that this note writes by hand (no declaration of the
% ----- unit stands behind them); the definitions of the unit come from macros148.tex and data.tex
\newcommand{\fm}[1]{\operatorname{fm}(#1)}                 % frameMatrix
\newcommand{\colL}[1]{\langle #1\rangle_{\hO}}             % colLattice Ô M: the Ô-span of the columns of M
\newcommand{\Wsys}[3]{W^{#1}_{#2,#3}}                      % (S.besselSys v).W σ
\newcommand{\Isys}[2]{I_{#2}(#1)}                          % (S.besselSys v).I
\newcommand{\muq}[1]{\mu^{\mathcal{U}\backslash}_{#1}}     % besselQuotMeasure
\newcommand{\muU}[1]{\mu^{\mathcal{U}}_{#1}}               % Shape.besselHaar
\newcommand{\Fw}[1]{\mathbb{F}_{#1}}                       % ResField w
\newcommand{\cj}[1]{\overline{#1}}
\newcommand{\NmE}{\operatorname{N}_{E/\mathbb{Q}}}
\newcommand{\omS}[2]{\omega_{#1,#2}}                       % the chosen box partition (Hilbert's choice in classCutoff)
\newcommand{\bfl}[2]{\beta^{\flat}_{#1,#2}}                % the chosen cutoff of U^0_c(g) in classCutoff
\newcommand{\spanE}{\operatorname{span}}
\newcommand{\Mid}{\operatorname{mid}}
\newcommand{\RePt}{\operatorname{Re}}
\newcommand{\ImPt}{\operatorname{Im}}
\newcommand{\lev}{\operatorname{lev}}
\newcommand{\rev}{\operatorname{rev}}
\newcommand{\col}{\operatorname{col}}
\newcommand{\Tr}{\operatorname{Tr}}
\newcommand{\sP}{\mathsf{P}}
\newcommand{\sA}{\mathsf{A}}
\newcommand{\sB}{\mathsf{B}}
\newcommand{\sC}{\mathsf{C}}
\newcommand{\sw}{\mathsf{w}}
\newcommand{\rF}[1]{#1_{\mathbb{F}}}                       % entrywise reduction to the residue field F (proof of Theorem A)

% where each declaration of the concordance is treated: \cs{key}{place}, the place naming a \label or a lettered item
\cs{isStablebesselS}{Lemma~\ref{l:stableS}}
\cs{measurePreservingpiorbitMeasure}{Lemma~\ref{l:unfold}}
\cs{tsumintegralapply}{Lemma~\ref{l:sumint}}
\cs{besselFnmulofmemclassUnip}{Lemma~\ref{l:Winv}}
\cs{besselKernelbesselUmul}{Lemma~\ref{l:Winv}}
\cs{lintegralclassCutoffmulprod}{Lemma~\ref{l:cutprod}}
\cs{amplitudeAbslttopandcutle}{Theorem D}
\cs{summableamplitude}{Corollary~\ref{c:summable}}
\cs{hasSumamplitudeSumchangeofweight}{Theorem E}
\cs{amplitudeSummove}{Proposition~\ref{p:moves}}
\cs{amplitudeSumbesselUmul}{Proposition~\ref{p:moves}}
\cs{amplitudeSumeqzeroofnotagrees}{Theorem F}
\cs{amplitudeSumbesselSone}{Theorem F}
\cs{liouvillegram}{Lemma~\ref{l:liouville}}
\cs{linkMetric}{Definition~\ref{d:linkmetric}}
\cs{linkDiscSq}{Definition~\ref{d:linkmetric}}
\cs{framePairlasteqzFormPairmul}{Lemma~\ref{l:formulas}}
\cs{classCoordeqfromBlocks}{Lemma~\ref{l:formulas}}
\cs{HasIntegralLifts}{Definition~\ref{d:lifts}}
\cs{finranklink}{Lemma~\ref{l:lagr}}
\cs{gramclassCoordX}{Lemma~\ref{l:formulas}}
\cs{classCoordgram}{Lemma~\ref{l:formulas}}
\cs{linkMetricposDef}{Lemma~\ref{l:posdef}}
\cs{gramlinkFormeqzerooflelink}{Lemma~\ref{l:lagr}}
\cs{detgramlinkMetricnezero}{Lemma~\ref{l:posdef}}
\cs{cardeqfinrankcolSpan}{Lemma~\ref{l:formulas}}
\cs{gramlinkMatrixmul}{Lemma~\ref{l:gramlink}}
\cs{classCoordmulframePair}{Lemma~\ref{l:formulas}}
\cs{classCoordmulisotropicFramePair}{Lemma~\ref{l:formulas}}
\cs{isIntMatrixframePair}{Lemma~\ref{l:formulas}}
\cs{legenIndexdetoffixesNormAtLeast}{Theorem A}
\cs{existsmemclassUnip}{Lemma~\ref{l:eichler}}
\cs{existseichlerclassUnip}{Lemma~\ref{l:eichlerrow}}
\cs{tsumleofencardrangele}{Lemma~\ref{l:colour}}
\cs{tsumleofcolouring}{Lemma~\ref{l:colourcoset}}
\cs{encardintegralle}{Lemma~\ref{l:intcount}}
\cs{submoduleIndexsmulintVectors}{Lemma~\ref{l:scaled}}
\cs{submoduleIndexsmulnezero}{Lemma~\ref{l:scaled}}
\cs{IndexBelow}{Definition~\ref{d:pad}}
\cs{xRows}{Definition~\ref{d:pad}}
\cs{paddedLeading}{Definition~\ref{d:pad}}
\cs{detpaddedLeading}{Lemma~\ref{l:paddet}}
\cs{frameMiddle}{Definition~\ref{d:pad}}
\cs{latticeDisc}{Definition~\ref{d:indices}}
\cs{frameValueIndex}{Definition~\ref{d:indices}}
\cs{xRowsIndex}{Definition~\ref{d:indices}}
\cs{saturationIndex}{Definition~\ref{d:indices}}
\cs{pairInvariant}{Definition~\ref{d:gauge}}
\cs{termElem}{Definition~\ref{d:gauge}}
\cs{minorNorm}{Definition~\ref{d:gauge}}
\cs{classGauge}{Definition~\ref{d:gauge}}
\cs{IsNearRep}{Definition~\ref{d:gauge}}
\cs{invPowWeight}{Definition~\ref{d:weights}}
\cs{termWeight}{Definition~\ref{d:weights}}
\cs{pairWeight}{Definition~\ref{d:weights}}
\cs{IsLabelledTerm}{Definition~\ref{d:weights}}
\cs{minorNormpos}{Lemma~\ref{l:levels}}
\cs{inv}{Lemma~\ref{l:levels}}
\cs{paddedLeadingmulinv}{Lemma~\ref{l:levels}}
\cs{existssharedframevector}{Lemma~\ref{l:levels}}
\cs{xRowsconjTransposemuleq}{Lemma~\ref{l:xrows}}
\cs{frameVectorsubeqsum}{Lemma~\ref{l:framediff}}
\cs{isBoxPartitionepsilon}{Proposition~\ref{p:boxexists}}
\cs{tsumclassCutoffeqone}{Lemma~\ref{l:sumone}}
\cs{tsumclassCutoffboxle}{Theorem B}
\cs{existsclassUnipmulmemboxedPoints}{Lemma~\ref{l:sumone}}
\cs{measurableclassCutoff}{Lemma~\ref{l:sumone}}
\cs{lintegralclassCutoffequnipCovolume}{Theorem C}

\begin{document}
\color{ink}

{\raggedright
{\sffamily\bfseries\fontsize{7.2}{9}\selectfont\color{leancol}\addfontfeatures{LetterSpace=8}READBACK NOTE 42\par}
\vspace{5pt}
{\fontsize{21.5}{25.5}\selectfont\bfseries Box partitions of arithmetic unipotent groups: frame counts, the covolume, amplitude sums\par}
\vspace{5pt}
{\fontsize{12.4}{16}\selectfont\itshape\color{inksoft} One unit of the Lean artifact \ArtifactName, read back into
mathematics\par}}
\vspace{8pt}
{\color{rulegrey}\hrule height 0.6pt}
\vspace{7pt}
\thispagestyle{empty}

\textbf{Summary.} The filling theorem of the artifact (Theorem~B of the write-up, ``enough forms'') rests on a
second moment of Poincar\'e series along a Bessel subgroup of a unitary group. Unfolded, this moment is a sum of
\emph{amplitudes} over rational elements, weighted by cutoffs for the arithmetic unipotent groups
$\clU{S}{c}$. This note treats the three files that build the first layer of estimates: the congruence that
agreeing unipotent characters force on pairings of frames (Theorem~A); cutoffs supported in boxes and the
resulting count of frames (Theorem~B); the exact covolume of $\clU{S}{c}$ (Theorem~C); and the amplitude sums
themselves: convergence, comparison of full and cut Bessel functions, change of weight, vanishing for
characters that do not agree, and the identity term (Theorems D, E, F).

\statement{leancol}{\fontsize{9.6}{13.2}\selectfont\textbf{What this note is.} The artifact \ArtifactName{} is a
body of Lean code whose final theorems state the Riemann hypothesis and its analogue for Dirichlet $L$-functions.
That is a claim: its authors report that Lean's kernel accepts the code, a replay of that check by the owner of
this write-up is recorded elsewhere, and nothing was compiled for this note; no mathematician has read the proof
through. This note treats the second part of the set-up of one unit (the folder of the filling theorem):
\UnitInNote{} of its \UnitMembers{} declarations, those of three files. The objects of the first part, the note
\emph{The filling data: Bessel subgroups, links of frames, Siegel's lemma and strong approximation} (Note~S below),
are recalled in Section~\ref{s:setting}, and its results are quoted as \emph{Inputs}, like the results of other
units and of the background library. Every statement of this note was compared, argument by argument, with the
Lean text, and the proofs are complete relative to the Inputs. The Lean text read is \PinText; Appendix~A is the
concordance of \ConcDecls{} declarations.}

\textbf{The notes of this unit} are cited by letters: S (the filling data), B (box partitions; this note),
U (Poincar\'e series and the sieve), E (the energy of amplitude sums), F (the filling theorem), P (the pairing of
links), T (the top floor at a finite place) and G (from amplitude sums to local correlations).

\sect{The statements}{s:intro}
The notation is that of Note~S and is recalled in Section~\ref{s:setting}. In a few words: $E$ is a CM field,
$\fs=(\varrho,\kappa,k_0)$ a shape, and $S$ filling data; $V$ is the hermitian space with basis
$e_0,\dots,e_{\varrho-1}$, $x_i$ $(i\in\kappa)$, $z$, $e^*_{\varrho-1},\dots,e^*_0$, the $e_a$ and $e^*_a$ in
hyperbolic pairs, the $x_i$ orthonormal and $h(z,z)=-\NS{S}$; the \emph{frame} is
$\fr{0},\dots,\fr{\varrho}=e_0,\dots,e_{\varrho-1},z$; $\UT{T}$ is the unipotent group of the flag of the $e_a$, and
$\ST{T}\supseteq\UT{T}$ the Bessel subgroup; $\Cl{S}$ is a finite set of classes, $\GRes{S}$ a finite group (a
product of general linear groups over the residue fields of the chosen places above $\gS{S}=\qS{S}{0}\qS{S}{1}$),
and $\clU{S}{c}$ the group of the rational $u\in\UT{\Trat{S}{E}}$ whose position at the class $c$ is congruent to
$1$ modulo $\gS{S}$.

\begin{setting}\label{set:su}
$E$ is a CM field with a CM type, $\fs$ is a shape and $S\in\Fill{E}{\fs}$ are filling data (Section~\ref{s:setting}).
Moreover $E$ is \emph{stably unramified}: it contains a primitive root of unity of order $5040$.
\end{setting}

\begin{setting}\label{set:adm}
$E$ is a CM field with a CM type that is stably unramified (it contains a primitive root of unity of order
$5040$), $\fs$ is a shape, and $S\in\Fill{E}{\fs}$ are filling data that are \emph{admissible}: the standing
hypotheses of Note~S hold for $S$; they are recalled in \ref{r:adm}.
\end{setting}

\statement{arith}{\phantomsection\label{t:A}\textbf{Theorem A (agreeing characters force congruences).}\enspace
In Setting~\ref{set:su}, let $k\in\KS{S}$ be such that $\Agr{S}{\redS{S}(k)}$, and let $x$ be a natural number. Let
$\Omega$ be a finite set with a bijection $\rho\colon\Omega\to\{1,2\}$, let $t$ be a matrix over $\AEf$ with rows
$\Ix{\fs}$ and columns $\Omega$ that is integral, and put $C=\fpair{k}^{*}\,t$; its rows are indexed by the pairs
$(i,a)$, $i\in\{1,2\}$, $0\le a\le\varrho$, the row $(1,a)$ belonging to $\fr{a}$ and the row $(2,a)$ to $k\fr{a}$.
Suppose that $C_{(i,a),l}=0$ for all $i\in\{1,2\}$, $a<\varrho$ and $l\in\Omega$, and let
\[\delta=\det\bigl(C_{(\rho(l),\varrho),\,l'}\bigr)_{l,l'\in\Omega}.\]
If $\delta$ is a unit of $\AEf$ and $\FixN{\redS{S}(k)}{x}$, then $x\le\gix{\delta}$.}

Here $\gix{\delta}$ is the generalized index (for $\delta\in\hO$ a unit of $\AEf$ it is $\#\hO/\delta\hO$), and
$\FixN{s}{x}$ says that $s$, or its inverse adjoint, fixes a non-zero vector of the span of the frame modulo
chosen places whose norms have product at least $x$. The columns of $t$ are two integral linear forms
$v\mapsto t_l^{*}v$ that vanish on $X=\spanE(e_a)$ and on $kX$; $\cj{\delta}$ is the determinant of their values at
$z$ and $kz$ (and $\gix{\cj{\delta}}=\gix{\delta}$). So wherever the frames of $1$ and of $k$ meet modulo a chosen place, that place divides $\delta$.

\statement{arith}{\phantomsection\label{t:B}\textbf{Theorem B (frame count).}\enspace
In Setting~\ref{set:adm}, let $c\in\Cl{S}$, $s\in\GRes{S}$, let $\sP$ be a set, $y,y'\in\prod_\tau\Gone{\fs}$ and
$x\in[0,\infty]$, and for every $p\in\sP$ let $d_p\in\Cl{S}$, $\gamma_p,\delta_p\in\GS{S}{E}$ and
$\sw_p\in[0,\infty]$ be given. Suppose that
\begin{enumerate}[label=(\roman*),itemsep=1pt]
\item for $p,q\in\sP$, if $d_p=d_q$ and $\delta_p=\delta_q$ then $p=q$;
\item for every $p$ with $\sw_p\ne0$: $\gamma_p\in\terms{S}{c}{s}{d_p}$,\ \ $\delta_p^{-1}\gamma_p\in\clU{S}{d_p}$,\ \
$\rinf{S}(\gamma_p)^{-1}y\in\bxp{S}$,\ \ and\ \ $\sw_p\le\cutoff{S}{d_p}\bigl(\rinf{S}(\delta_p)^{-1}y'\bigr)\cdot x$;
\item $V$ is a real subspace of the space of families $(v_{\tau,a})$ of vectors of $\CC^{\Ix{\fs}}$, indexed by
the infinite places $\tau$ of $E$ and $0\le a\le\varrho$, and $F_p-F_q\in V$ for all $p,q$ with $\sw_p\ne0\ne\sw_q$,
where $F_p=\bigl((\rinf{S}(\gamma_p)_\tau^{-1}y_\tau)^{-1}\fr{a}\bigr)_{\tau,a}$.
\end{enumerate}
Then\quad $\sum_{p\in\sP}\sw_p\ \le\ \bS{S}^{\,6\,\nn{\fs}^2\dim_{\RR}V}\cdot x$.}

\statement{arith}{\phantomsection\label{t:C}\textbf{Theorem C (the covolume).}\enspace
In Setting~\ref{set:adm}, for every $c\in\Cl{S}$ and every $y\in\prod_\tau\Gone{\fs}$,
\[\int_{\prod_\tau\UT{\Trg{\fs}{1}}}\cutoff{S}{c}\bigl((u_\tau y_\tau)_\tau\bigr)\,d\Bigl(\prod_\tau\muU{\fs}\Bigr)(u)\ =\ \ucov{S},\]
the integral over the product, over the infinite places $\tau$ of $E$, of the unipotent groups of $\Gone{\fs}$
with their Haar measures $\muU{\fs}$ of \ref{r:haar}; $\ucov{S}$ is the explicit number of \ref{r:const}.}

Since $\sum_{u\in\clU{S}{c}}\cutoff{S}{c}(\rinf{S}(u)y)=1$, Theorem~C says that the image of $\clU{S}{c}$ is a lattice
of covolume $\ucov{S}$ in $\prod_\tau\UT{\Trg{\fs}{1}}$, whatever the class $c$.

\statement{arith}{\phantomsection\label{t:D}\textbf{Theorem D (convergence; full against cut).}\enspace
In Setting~\ref{set:adm}, let $\sigma\in\{\full,\cut\}$ be a side, $\Xi\subseteq\GS{S}{E}$ a set, $c\in\Cl{S}$ and
$s\in\GRes{S}$, and put $X(d)=\Xi\cap\terms{S}{c}{s}{d}$ for $d\in\Cl{S}$. Then:
\begin{enumerate}[label=(\arabic*),itemsep=1pt]
\item for every $c'\in\Cl{S}$ there is a real number $M\ge0$ such that for all $y\in\prod_\tau\Gone{\fs}$
\[\sum_{\gamma\in X(c')}\cutoff{S}{c'}\bigl(\rinf{S}(\gamma)^{-1}y\bigr)\prod_\tau\bigl\|\Wsys{\sigma}{S}{\tau}\bigl((\rinf{S}(\gamma)^{-1}y)_\tau\bigr)\bigr\|\ \le\ M,\]
and moreover $\sum_{\gamma\in X(c')}\ampabs{\sigma}{S}{c}{c'}{\gamma}<\infty$;
\item with $f=\ninf{E}$ and $k=\kk{\fs}$,
\begin{multline*}\sum_{d\in\Cl{S}}\ \sum_{\gamma\in X(d)}\int\cutoff{S}{c}(y)\,\cutoff{S}{d}\bigl(\rinf{S}(\gamma)^{-1}y\bigr)\,
\bigl\|\besker{\sigma}{S}{\gamma}{y}-\besker{\cut}{S}{\gamma}{y}\bigr\|\,dy\\ \le\ \frac{\ucov{S}}{2\,\bS{S}^{\,k^3f}}.\end{multline*}
\end{enumerate}}

\statement{arith}{\phantomsection\label{t:E}\textbf{Theorem E (change of weight).}\enspace
In Setting~\ref{set:adm}, let $\sigma$ be a side, $\Xi\subseteq\GS{S}{E}$ a set with $\IsStab{S}{\Xi}$, let
$c,c'\in\Cl{S}$ and $s\in\GRes{S}$, and put $X=\Xi\cap\terms{S}{c}{s}{c'}$. Let
$\phi\colon\GS{S}{E}\times\prod_\tau\Gone{\fs}\to[0,\infty]$ be such that every $\phi(\gamma,\cdot)$ is measurable and
\[\sum_{(u,u')\in\clU{S}{c}\times\clU{S}{c'}}\phi\bigl(u\gamma u'^{-1},\ \rinf{S}(u)\,y\bigr)=1
\qquad\text{for all }\gamma\in X\text{ and all }y.\]
Then for every $\gamma\in X$ the function $y\mapsto\phi(\gamma,y)\,\besker{\sigma}{S}{\gamma}{y}$ is integrable, the
family of its integrals is summable with
\[\sum_{\gamma\in X}\int\phi(\gamma,y)\,\besker{\sigma}{S}{\gamma}{y}\,dy\ =\ \ampsum{\sigma}{S}{\Xi}{c}{s}{c'},\]
and\quad $\displaystyle\sum_{\gamma\in X}\int\phi(\gamma,y)\,\bigl\|\besker{\sigma}{S}{\gamma}{y}-\besker{\cut}{S}{\gamma}{y}\bigr\|\,dy
\ \le\ \frac{\ucov{S}}{2\,\bS{S}^{\,\kk{\fs}^3\ninf{E}}}$.}

\statement{arith}{\phantomsection\label{t:F}\textbf{Theorem F (characters that do not agree; the identity term).}\enspace
In Setting~\ref{set:adm}, let $c,c'\in\Cl{S}$.
\begin{enumerate}[label=(\roman*),itemsep=1pt]
\item Let $\sigma$ be a side, $\Xi\subseteq\GS{S}{E}$ a set with $\IsStab{S}{\Xi}$, and $s\in\GRes{S}$. If $s$
does not agree, that is, if $\Agr{S}{s}$ fails, then $\ampsum{\sigma}{S}{\Xi}{c}{s}{c'}=0$.
\item For $\Xi=\ST{\Trat{S}{E}}$, the group of rational points of the Bessel subgroup, for the full Bessel
functions, and for all $i,j\in\Bc{S}$, the coordinate $(i,j)$ of $\ampsum{\full}{S}{\ST{\Trat{S}{E}}}{c}{1}{c'}$ is
$\ucov{S}$ if $c=c'$ and $i=j$, and $0$ otherwise.
\end{enumerate}}

Theorems A to C are proved in Sections \ref{s:agree}, \ref{s:boxes} and \ref{s:covol}; Theorems D, E, F in
Section~\ref{s:ampl}. Sections \ref{s:links}, \ref{s:eichler} and \ref{s:count} hold the tools: a positive form on
links, integral Eichler elements, and covering lemmas, together with the definitions of the indices,
discriminants and weights that the later estimates of the unit use (Note~E, \emph{The energy of amplitude sums},
and Note~P, on the pairing of links).

\sect{The setting}{s:setting}
\begin{conventions}\phantomsection\label{c:conventions}
Throughout, $E$ is a CM field with a fixed CM type: a number field that is a totally imaginary quadratic extension
of a totally real field, with its complex conjugation $x\mapsto\cj{x}$ and, for every infinite place $\tau$, a chosen
embedding $x\mapsto x_\tau$ of $E$ into $\CC$; $\ninf{E}$ is the number of its infinite places, $d_E$ its
discriminant, $\AEf$ its ring of finite adeles with the involution induced by conjugation, and $\hO\subset\AEf$ the
closure of the integers. The letter $\fs$ denotes a shape, and $S$, $S'$, $S_1$ denote filling data in
$\Fill{E}{\fs}$ \ref{r:data}. $E'$ is a number field. $R$ and $A$ are commutative rings with an involution,
$T$ is an $\fs$-form over $R$, and $O$ a subring. Index sets written $\Omega$, $\Omega'$ are finite sets; $\sP$ is
an arbitrary set with a subset $\sA$; a matrix has its rows and columns indexed by the sets that are named, and
$M^{*}$ is its conjugate transpose. An element $k$, $\gamma$, $\gamma'$, $\gamma_2$ of a unitary group is
identified with its matrix. A side is an element $\sigma$ of $\{\full,\cut\}$. Indices of frames and levels are
counted from $0$, as in the formal text; $j$, $j'$ denote levels and $c$, $c'$, $d$ classes. A sum of numbers in
$[0,\infty]$ over an arbitrary set is the supremum of its finite partial sums. A matrix over a ring
containing $\hO$ is \emph{integral} if its entries lie in $\hO$ (or in the subring $O$ that is named). The letters
$S$, $s$, $\sigma$, $q$, $\theta$, $\Theta$, $\psi$, $\zeta$ have in the notes of this unit the meanings given here,
not those of the overview of the write-up (there: a set of columns, the complex variable, a field automorphism, a
rational number, a family of twists, a chain of lifts, a column, Riemann's function). For the filling data of
admissible data (Note~G), the primes $\qS{S}{0}$, $\qS{S}{1}$ are the primes that the overview calls $g_0$, $g_1$, and
$\gS{S}$ is their product.
\end{conventions}

\subhead{Recollections from Note~S.} The following objects are defined in Note~S (and, for the first three
items, in the background library); they are recalled without their formal names, which are in the concordance of
that note.
\begin{enumerate}[label=(R\arabic*),leftmargin=3.1em,itemsep=3pt]
\item\label{r:space} \emph{Shapes, forms, groups.} A shape $\fs=(\varrho,\kappa,k_0)$ is a natural number
$\varrho$, a finite set $\kappa$ and an element $k_0\in\kappa$; $\kk{\fs}=\#\kappa$, $\nn{\fs}=\varrho+1$. Its
index set $\Ix{\fs}$ labels the vectors $e_0,\dots,e_{\varrho-1}$; $x_i$ $(i\in\kappa)$; $z$;
$e^*_{\varrho-1},\dots,e^*_0$, so $\rr{\fs}=\#\Ix{\fs}=2\varrho+\kk{\fs}+1$; the middle index set
$\Mx{\fs}=\kappa\sqcup\{z\}$ labels the $x_i$ and $z$. An $\fs$-form over $R$ is a matrix $T$ on $\Mx{\fs}$; $\JT{T}$
is the matrix on $\Ix{\fs}$ that pairs $e_a$ with $e^*_a$ by $1$, is $T$ on the middle block and $0$ elsewhere;
$h_T(v,w)=v^{*}\JT{T}w$, and $\GT{T}$ is the group of the invertible $g$ with $g^{*}\JT{T}g=\JT{T}$. For $N\in R$,
$\Trg{\fs}{N}$ is the diagonal form with entries $1$ on $\kappa$ and $-N$ at $z$, and
$\Gone{\fs}=\GT{\Trg{\fs}{1}}$ over $\CC$. The \emph{level} of an index is $\lev(e_a)=a$, $\lev(x_i)=\lev(z)=\varrho$,
$\lev(e^*_a)=2\varrho-a$. The \emph{frame} is $\fr{a}=e_a$ for $a<\varrho$ and $\fr{\varrho}=z$; we write $\mathfrak f$ for
the matrix with columns $\fr{0},\dots,\fr{\varrho}$, so that $g\mathfrak f$ has the columns $g\fr{a}$ and
$\mathfrak f x=\sum_a x_a\fr{a}$, and $X=\spanE(e_0,\dots,e_{\varrho-1})$. The subgroups of $\GT{T}$ are: $\UT{T}$,
the $g$ with $g_{il}=\delta_{il}$ whenever $\lev(l)\le\lev(i)$; $\UoT{T}$, the $g\in\UT{T}$ with
$g\mathfrak f=\mathfrak f$; $\HT{T}$, the $g$ fixing all $e_a$, $e^*_a$ and $z$; and the Bessel subgroup $\ST{T}$,
the $g$ with $g(i,\fr{a})=\delta_{i,\fr{a}}$ whenever $\lev(i)\ge a$. Then $\UT{T}$ and $\HT{T}$ lie in $\ST{T}$, and
$\ST{T}$ normalizes $\UT{T}$. For $g\in\GT{T}$, $\sd{g}=\sum_{a<\varrho}g(\fr{a},\fr{a+1})$, the frame matrix is
$\fm{g}=(g(\fr{a},\fr{b}))_{a,b}$, and $\ZT{T}{g}=(h_T(\fr{a},g\fr{b}))_{a,b}$; $\lmin{l}{Z}$ is the leading
principal minor of order $l$ of a square matrix $Z$. On $\UT{T}$ the map $\fm{\cdot}$ is multiplicative with
upper unitriangular values, and $\fm{g}=1$ exactly for $g\in\UoT{T}$.
\item\label{r:real} \emph{At the complex numbers.} For $y\in\Gone{\fs}$, $\|y\|$ is the operator norm of its
matrix on the Euclidean space $\CC^{\Ix{\fs}}$, the vectors $y^{-1}\fr{a}$ are its \emph{frame vectors},
$\bsize{y}=\inf_{s\in\ST{\Trg{\fs}{1}}}(\|sy\|+1)$, and $\Zr{y}=\ZT{\Trg{\fs}{1}}{y}$. For real $b$,
$\InCut{y}{b}$ means $\bsize{y}\le b^{\,6\nn{\fs}+2}$, and for real $\rho$ the set $\frbox{\rho}$ consists of the
$y$ with $\|y^{-1}\fr{a}\|\le\rho$ for all $a$.
\item\label{r:haar} \emph{Haar measures.} $dy$ is the Haar measure of $\Gone{\fs}$ fixed by the background
library, and also its product on $\prod_\tau\Gone{\fs}$. On the unipotent group $\UT{\Trg{\fs}{1}}$ of $\Gone{\fs}$
the library fixes a measure $\muU{\fs}$ that is invariant under left and right translations and under
inversion, and $\muq{\fs}=\beta\,dy$ for a chosen measurable $\beta\ge0$ with
$\int\beta(uy)\,d\muU{\fs}(u)=1$ for all $y$. The map $(x,u)\mapsto ux$ carries $\muq{\fs}\times\muU{\fs}$ to $dy$; so
$\int F\,d\muq{\fs}$ is the integral over the quotient by the unipotent group of a function $F\ge0$ that is
invariant under it.
\item\label{r:data} \emph{Filling data.} A record $S\in\Fill{E}{\fs}$ has, among its components: a prime number
$\NS{S}$; two prime numbers $\qS{S}{0},\qS{S}{1}\ne\NS{S}$, with product $\gS{S}$; a finite idele $\piS{S}\in\hO$
with $\piS{S}\,\cj{\piS{S}}=\NS{S}$; an integer $\vpS{S}$ of $E$ that generates the different, with
$\cj{\vpS{S}}=-\vpS{S}$; a real number $\bS{S}$; for $i\in\{0,1\}$ a set $\ChS{S}{i}$ of $\ninf{E}$ finite places of
$E$ of norm $\qS{S}{i}$, the \emph{chosen places}, pairwise distinct, none the conjugate of another; and for
every infinite place $\tau$ a Bessel system \ref{r:bessel}. For $r\in\hO$, $\gS{S}$ divides $r$ if and only if
$r$ and $\cj{r}$ vanish modulo every chosen place. Among the conditions: $\varrho\ge10^5$;
$1000\,(2\varrho+2)^9\le\kk{\fs}\le\ninf{E}$; $\ninf{E}\le\bS{S}/1000$; $\varrho+1\le\bS{S}$; and
$\bS{S}^{10}\le\qS{S}{i}\le\bS{S}^{11}$.
\item\label{r:groups} \emph{Groups and coordinates.} $\Trat{S}{R}=\Trg{\fs}{\NS{S}}$ over $R$ and
$\GS{S}{R}=\GT{\Trat{S}{R}}$. The lattice basis $\latB{S}$ is the matrix over $\AEf$ that is the identity on the
$e_a$ and $e^*_a$ and has the columns $x_i$ $(i\ne k_0)$, $x_{k_0}/\cj{\piS{S}}+z/\NS{S}$ and $z$ on the middle
block; in this basis the form has the Gram matrix $\JT{\Tlat{S}{\AEf}}$, where $\Tlat{S}{R}$ has diagonal entries $1$
at $i\ne k_0$, $0$ at $k_0$, $-\NS{S}$ at $z$, and entries $-1$ at $(k_0,z)$ and $(z,k_0)$. $\GFin{S}=\GT{\Tlat{S}{\AEf}}$;
$\KS{S}$ is its subgroup of the integral elements with integral inverse; $\GRes{S}=\prod_{i}\prod_{w\in\ChS{S}{i}}\GL_{\Ix{\fs}}(\Fw{w})$,
with $\Fw{w}$ the residue field at $w$; $\redS{S}\colon\KS{S}\to\GRes{S}$ is the reduction of entries modulo the
chosen places, $\KgS{S}$ its kernel, and $\redim{S}{A}=\redS{S}(A\cap\KS{S})$ for a subgroup $A$ of $\GFin{S}$. The maps
$\rfin{S}\colon\GS{S}{E}\to\GFin{S}$, $\gamma\mapsto\latB{S}^{-1}\gamma\,\latB{S}$, and
$\rinf{S}\colon\GS{S}{E}\to\prod_\tau\Gone{\fs}$, $(\rinf{S}(\gamma)_\tau)_{il}=(\gamma_{il})_\tau\,\bscale{S}{\tau}{l}/\bscale{S}{\tau}{i}$,
are homomorphisms, with the scales of \ref{r:const}; $\rinf{S}$ maps $\UT{\Trat{S}{E}}$, $\UoT{\Trat{S}{E}}$,
$\ST{\Trat{S}{E}}$ into the groups of the same name of $\Gone{\fs}$. $\Cl{S}$ is the finite set of double cosets
of the finite adelic points of the unitary group of $X_0=\spanE(x_i)$ by its rational points and the level of the
record; every class $c$ has a fixed representative $\tilde c$ (a unitary matrix on $\AEf^{\kappa}$), and
$\clB{S}{c}=\operatorname{diag}(1,\tilde c,1,1)\,\latB{S}$ (blocks $e$, $\kappa$, $z$, $e^*$), $\clC{S}{c}=\clB{S}{c}^{-1}$,
$\clCX{S}{c}=\mbc{\fs}{\piS{S}}\,\tilde c^{-1}$, where $\mbc{\fs}{\piS{S}}$ is the matrix with rows $\Mx{\fs}$ and
columns $\kappa$ that expresses the $x_i$ in the middle part of the lattice basis; $\hc{S}{c}\in\GFin{S}$ is the
image of $\tilde c$. For $c,c'\in\Cl{S}$ and $\gamma\in\GS{S}{E}$ the \emph{position} is
$\pos{S}{c}{c'}{\gamma}=\hc{S}{c}^{-1}\rfin{S}(\gamma)\hc{S}{c'}=\clB{S}{c}^{-1}\gamma\,\clB{S}{c'}\in\GFin{S}$;
for $s\in\GRes{S}$, $\terms{S}{c}{s}{c'}$ is the set of the $\gamma$ with $\pos{S}{c}{c'}{\gamma}\in\KS{S}$ of
reduction $s$ (the \emph{terms over} $(c,s,c')$);
$\clU{S}{c}=\{u\in\UT{\Trat{S}{E}}:\pos{S}{c}{c}{u}\in\KgS{S}\}$ and $\clUo{S}{c}=\clU{S}{c}\cap\UoT{\Trat{S}{E}}$.
$\GS{S}{E}$ is countable.
\item\label{r:char} \emph{Characters and reductions.} $\gid{S}$ is the idele of $\gS{S}$; $\sdc{S}$ is the
additive character $x\mapsto\psi_f(-x/\vpS{S})$ of $\AEf$, $\psi_f$ the standard character; and
$\uchar{S}{u}=\sdc{S}\bigl(\gid{S}^{-1}\sd{u}\bigr)$ for $u\in\GFin{S}$. For a chosen place $w$ and
$\epsilon\in\{0,1\}$, $\rag{S}{(w,\epsilon)}\colon\hO\to\Fw{w}$ is the reduction modulo $w$, preceded by the
conjugation if $\epsilon=1$; $\ioia{S}{\epsilon}$ is the automorphism of $\GL_{\Ix{\fs}}(\Fw{w})$ that is the
identity if $\epsilon=0$ and $s\mapsto J^{-1}(s^{-1})^{\mathsf T}J$, $J=\JT{\Tlat{S}{\Fw{w}}}$, if $\epsilon=1$. For
$w\in\ChS{S}{i}$ and $\mathbf d=(\epsilon,v)$ with $0\ne v\in\Fw{w}^{\varrho+1}$, $\vfix{S}{w}{\mathbf d}$ is the subgroup
of the $s\in\GRes{S}$ such that $\ioia{S}{\epsilon}(s_{i,w})$ fixes $\mathfrak f v$. For a natural number $x$,
$\FixN{s}{x}$ means: there is a finite set $\sA$ of pairs $(w,\epsilon)$ such that $x$ is at most the product of
the norms of the places $w$ over $(w,\epsilon)\in\sA$, and for each of them some $v\ne0$ with
$s\in\vfix{S}{w}{(\epsilon,v)}$. An element $s\in\GRes{S}$ \emph{agrees}, $\Agr{S}{s}$, if
$\uchar{S}{u}=\uchar{S}{u'}$ for all $u,u'\in\KS{S}\cap\UT{\Tlat{S}{\AEf}}$ with $\redS{S}(u)\,s=s\,\redS{S}(u')$.
A set $\Xi\subseteq\GS{S}{E}$ is \emph{stable}, $\IsStab{S}{\Xi}$, if $u^{-1}\gamma u'\in\Xi\Leftrightarrow\gamma\in\Xi$
for all $u,u'\in\UT{\Trat{S}{E}}$ and all $\gamma$.
\item\label{r:links} \emph{Frames and links.} For $\gamma\in\GT{T}$: $\fpair{\gamma}$ is the matrix with the
columns $\fr{a}$ and $\gamma\fr{a}$ ($0\le a\le\varrho$, indexed by $(1,a)$ and $(2,a)$), $\ifpair{\gamma}$ its
submatrix of the columns with $a<\varrho$; $\lmat{\gamma}$ is the matrix of $t\mapsto(\Mid t,\Mid\gamma^{-1}t)$,
where $\Mid t\in R^{\Mx{\fs}}$ is the middle part of $t\in R^{\Ix{\fs}}$; the \emph{link} is
$\link{\gamma}=\lmat{\gamma}\{t:\ifpair{\gamma}^{*}\JT{T}t=0\}\subseteq R^{\Mx{\fs}\sqcup\Mx{\fs}}$; $\zfp{T}$ is the
matrix with two rows that sends $(p,p')$ to $((Tp)_z,(Tp')_z)$; and $\clink{\gamma}{\gamma'}\subseteq R^{\kappa}$
consists of the $\kappa$-parts of the first components of the $(p,p')\in\link{\gamma}\cap\link{\gamma'}$ with
$\zfp{T}(p,p')=0$. For filling data, $\ratG{S}=\JT{\Trat{S}{E}}$,
$\lform{S}=\operatorname{diag}(\Trat{S}{E},-\Trat{S}{E})$, $\Zf{S}{\gamma}=\ZT{\Trat{S}{E}}{\gamma}$,
$\tdisc{S}{\gamma}=\fdisc{\NS{S}}{\Zf{S}{\gamma}}$ with
$\fdisc{N}{Z}=\det Z\,\cj{\det Z}-N^2\,\lmin{\varrho}{Z}\,\cj{\lmin{\varrho}{Z}}$, and
$\tdnorm{S}{\gamma}=\prod_\tau|\tdisc{S}{\gamma}_\tau|$. For $j\le\varrho+1$, $\IsNF{j}{\gamma}$ (\emph{normal form
from~$j$}, as in Note~S) means: there is $c_0\in R$ with $c_0\cj{c_0}=1$ such that for every $a$ with $j\le a\le\varrho$ there
are $x,y\in R^{\varrho+1}$ with $\mathfrak f x=\gamma\mathfrak f y$, $x_a=c_0$, $y_a=1$ and $x_b=y_b=0$ for $b>a$; and
$\HasLev{S}{j}{\gamma}$ means $\IsNF{j}{\gamma}$, $\lmin{j}{\Zf{S}{\gamma}}\ne0$, and $j=\varrho+1$ exactly if
$\tdisc{S}{\gamma}\ne0$. $\Cong{S}{c}{c'}{\gamma}{\gamma'}$ means that $\gamma,\gamma'\in\terms{S}{c}{s}{c'}$ for one
$s$. For a matrix $M$ over $E$, $\fadm{M}$ is the same matrix over $\AEf$. For matrices $L,H$ over $\AEf$ with
$LH=1$ and a subspace $V$ of the space $E^{\Omega}$ on which $H$ acts, $\slb{H}{V}$ is a chosen matrix $\fadm{y}\,a$,
$y$ a basis matrix of $V$ and $a$ an invertible matrix over $\AEf$, such that $H\,\slb{H}{V}$ is integral with an
integral left inverse, and $\sld{H}{V}=\sqrt{\gix{\det(\slb{H}{V}^{*}\slb{H}{V})}}$. Finally
$\cw{S}{i}=\#\redim{S}{\ST{\Tlat{S}{\AEf}}}\cdot\#\Cl{S}\cdot\NS{S}^{-\kk{\fs}\ninf{E}}\,(2\,\tailS{S}^2)^{i}$.
\item\label{r:bessel} \emph{Bessel systems.} The Bessel system of $S$ at $\tau$ provides a finite non-empty set
$\Isys{S}{\tau}$, a real number $\deg_\tau(S)>0$ and continuous functions
$\Wsys{\full}{S}{\tau},\Wsys{\cut}{S}{\tau}\colon\Gone{\fs}\to\CC^{\Isys{S}{\tau}}$, with properties of which we
use those of Input~(S7). $\Bc{S}=\prod_\tau\Isys{S}{\tau}$; for $y\in\prod_\tau\Gone{\fs}$, $i\in\Bc{S}$ and a side
$\sigma$, $\besfn{\sigma}{S}{y}{i}=\prod_\tau\Wsys{\sigma}{S}{\tau}(y_\tau)_{i_\tau}$, and for $\gamma\in\GS{S}{E}$ the
\emph{kernel} $\besker{\sigma}{S}{\gamma}{y}\in\CC^{\Bc{S}\times\Bc{S}}$ has the coordinates
$\besfn{\sigma}{S}{y}{i}\,\cj{\besfn{\sigma}{S}{\rinf{S}(\gamma)^{-1}y}{j}}$, with the Euclidean norm.
\item\label{r:cutoff} \emph{Boxes, cutoffs, amplitudes.} $\bxs{S}=\bS{S}^{2\nn{\fs}^2}$; $\bxp{S}$ is the set of
the $y\in\prod_\tau\Gone{\fs}$ with $y_\tau\in\frbox{\bxs{S}}$ and $\InCut{y_\tau}{\bS{S}}$ for all $\tau$; and
$\IsNr{S}{\gamma}$ if some $y\in\bxp{S}$ has $\rinf{S}(\gamma)^{-1}y\in\bxp{S}$. A function
$\omega\colon\prod_\tau\Gone{\fs}\to[0,\infty]$ is a \emph{box partition} for $c\in\Cl{S}$, $\IsBP{S}{c}{\omega}$, if:
(i) $\omega$ is measurable; (ii) if $\omega(y)\ne0$ and $\InCut{y_\tau}{\bS{S}}$ then $y_\tau\in\frbox{\bxs{S}}$;
(iii) for every $y$ and every $\alpha\colon\clU{S}{c}\to[0,\infty]$ with
$\sum_{u\in\clUo{S}{c}}\alpha(u\gamma)=1$ for all $\gamma$, one has
$\sum_{\gamma\in\clU{S}{c}}\alpha(\gamma)\,\omega(\rinf{S}(\gamma)y)=1$; (iv) $\omega\le1$; (v) $\omega(sy)=\omega(y)$
whenever every $s_\tau$ lies in the group generated by $\UoT{\Trg{\fs}{1}}$ and $\HT{\Trg{\fs}{1}}$. The cutoff is
$\cutoff{S}{c}=\bfl{S}{c}\cdot\omS{S}{c}$, where $\bfl{S}{c}$ is a chosen measurable function with
$\sum_{u\in\clUo{S}{c}}\bfl{S}{c}(\rinf{S}(u)y)=1$ for all $y$ (Input~(S6)) and $\omS{S}{c}$ is a chosen box
partition for $c$ if there is one (and an arbitrary function otherwise). The amplitudes are
\[\amp{\sigma}{S}{c}{c'}{\gamma}=\int\cutoff{S}{c}(y)\,\cutoff{S}{c'}\bigl(\rinf{S}(\gamma)^{-1}y\bigr)\,\besker{\sigma}{S}{\gamma}{y}\,dy,\]
\[\ampsum{\sigma}{S}{\Xi}{c}{s}{c'}=\sum_{\gamma\in\Xi\cap\terms{S}{c}{s}{c'}}\amp{\sigma}{S}{c}{c'}{\gamma},\]
and $\ampabs{\sigma}{S}{c}{c'}{\gamma}\in[0,\infty]$ is the first integral with the kernel replaced by its norm. In the
first integral a value $\infty$ of the weight is read as $0$; by Lemma~\ref{l:sumone}(a) the weight is at most $1$.
\item\label{r:const} \emph{Constants.} $\kapS{S}{\tau}=\gS{S}\,\vpS{S}_\tau/(2\pi i)$, a non-zero real number.
The scales are $\bscale{S}{\tau}{e_a}=\kapS{S}{\tau}^{\varrho-a}/\sqrt{\NS{S}}$, $\bscale{S}{\tau}{x_i}=1$,
$\bscale{S}{\tau}{z}=1/\sqrt{\NS{S}}$, $\bscale{S}{\tau}{e^*_a}=\bscale{S}{\tau}{e_a}^{-1}$; so
$\bscale{S}{\tau}{\fr{a}}=\kapS{S}{\tau}^{\varrho-a}/\sqrt{\NS{S}}$ for $a\le\varrho$. Further
$\fscale{S}{\tau}{a}=\kapS{S}{\tau}^{2(\varrho-a)}/\NS{S}$,
$\fsprod{S}=\prod_\tau\prod_{a\le\varrho}\fscale{S}{\tau}{a}^{\kk{\fs}+2(\varrho-a)}$, $\tailS{S}=\bS{S}^{\kk{\fs}\ninf{E}/4}$, and,
with $k=\kk{\fs}$, $f=\ninf{E}$ and $\volO{E}=2^{f}/\sqrt{|d_E|}$,
\[\ucov{S}=\Bigl(\frac{\gS{S}^{\varrho(2\varrho+2k+1)}}{\NS{S}^{k}}\Bigr)^{f}\cdot\frac{|d_E|^{\varrho/4}}{\volO{E}^{\varrho(\varrho+k)}\ \fsprod{S}}\ >0.\]
\item\label{r:adm} \emph{Admissibility.} The filling data $S$ are admissible if a list of conditions of Note~S
holds, of which this note uses: $|\vpS{S}_\tau|\le\bS{S}$ and $1\le|\kapS{S}{\tau}|\le\bS{S}^{23}$ for all $\tau$;
$\NS{S}\le\bS{S}^{60\nn{\fs}}$; $\#\Bc{S}\le\bS{S}^{2\kk{\fs}\ninf{E}}$; and that the level of the record is the
preimage of $\KgS{S}$, which enters through Input~(S3).
\end{enumerate}

\statement{ana}{\textbf{Inputs from Note~S} (of this unit; \emph{The filling data: Bessel subgroups, links of
frames, Siegel's lemma and strong approximation}). $S\in\Fill{E}{\fs}$ throughout; $c,c',d\in\Cl{S}$.

\textbf{(S1)} \emph{Levels and characters.} $\KgS{S}$ is the set of the $k\in\KS{S}$ with $\gid{S}^{-1}(k-1)$
integral. $\JT{\Tlat{S}{\AEf}}$ is hermitian, integral with integral inverse. For $a\in\AEf$:
$\sdc{S}(ra)=1$ for all $r\in\hO$ if and only if $a\in\hO$. For $r\in\hO$: $\gid{S}^{-1}r\in\hO$ if and only if
$\rag{S}{(w,\epsilon)}(r)=0$ for all chosen $w$ and both $\epsilon$. Every family of residues
$(t_{w,\epsilon})$, $t_{w,\epsilon}\in\Fw{w}$, is the family of the $\rag{S}{(w,\epsilon)}(r)$ of some $r\in\hO$. For
$k\in\KS{S}$, a chosen place $w\in\ChS{S}{i}$ and $\epsilon\in\{0,1\}$: the entrywise image of $k$ under
$\rag{S}{(w,\epsilon)}$ is $\ioia{S}{\epsilon}(\redS{S}(k)_{i,w})$. The residue field of $w\in\ChS{S}{i}$ has
$\qS{S}{i}$ elements.

\textbf{(S2)} \emph{Eichler elements.} Let $A$ carry an element $\tau_0$ with $\tau_0+\cj{\tau_0}=1$, and let $T$
be an $\fs$-form over $A$ with $\JT{T}$ hermitian and invertible. Every linear form on $A^{\Ix{\fs}}$ is
$h_T(v,\cdot)$ for some $v$. Let $a\le\varrho$, $j\in A^{\varrho+1}$ with $j_b=0$ for $b\ge a$, $r\in A$, and let
$\mu=h_T(v,\cdot)$ vanish at $\fr{b}$ for all $b<a$. Then there is $u\in\UT{T}$ with
$\sd{u}=r\sum_{b=1}^{\varrho}\mu(\fr{b})\,j_{b-1}$ whose matrix is
$1+\xi\,v^{*}\JT{T}-(v+\varsigma\xi)\,\xi^{*}\JT{T}$, with $\xi=r\,\mathfrak f j$ and $\varsigma=\tau_0\,h_T(v,v)$.

\textbf{(S3)} \emph{Positions and terms.} $\clC{S}{c}^{*}\,\JT{\Tlat{S}{\AEf}}\,\clC{S}{c}=\fadm{\ratG{S}}$, and
$\pos{S}{c}{c'}{\gamma}\,\mathfrak f=\clC{S}{c}\,\fadm{\gamma\mathfrak f}$. If $\gamma\in\terms{S}{c}{s}{c'}$ and
$\gamma'\in\terms{S}{c'}{s'}{d}$ then $\gamma\gamma'\in\terms{S}{c}{ss'}{d}$ and
$\gamma^{-1}\in\terms{S}{c'}{s^{-1}}{c}$; for such $\gamma$ the matrix $\pos{S}{c}{c'}{\gamma}$ is integral and all
entries of $\Zf{S}{\gamma}$ are integral over $\mathbb Z$. Let $S$ be admissible. If $\gamma\in\ST{\Trat{S}{E}}$ and
$\pos{S}{c}{d}{\gamma}\in\KgS{S}$, then $c=d$ and $\gamma\in\clU{S}{c}$. If $\gamma\in\terms{S}{c}{s}{c'}$,
$\gamma'\in\terms{S}{c}{s}{d}$ and $\gamma\mathfrak f=\gamma'\mathfrak f$, then $c'=d$ and $\gamma^{-1}\gamma'\in\clUo{S}{c'}$.

\textbf{(S4)} \emph{Characters of the unipotent group.} For $\gamma\in\UT{\Trat{S}{E}}$, every side $\sigma$ and
all $y$, $i$: $\besfn{\sigma}{S}{\rinf{S}(\gamma)y}{i}=\uchar{S}{\rfin{S}(\gamma)}^{-1}\,\besfn{\sigma}{S}{y}{i}$. If
$u,v\in\UT{\Tlat{S}{\AEf}}$ and $u^{-1}\hc{S}{c}^{-1}v\,\hc{S}{c}\in\KgS{S}$, then $\uchar{S}{v}=\uchar{S}{u}$; this
holds with any element of the image of the finite adelic unitary group of $X_0$ in the place of $\hc{S}{c}$. For
every $u\in\UT{\Tlat{S}{\AEf}}$ there is $u_0\in\UT{\Trat{S}{E}}$ with $u^{-1}\pos{S}{c}{c}{u_0}\in\KgS{S}$ and
$\uchar{S}{\rfin{S}(u_0)}=\uchar{S}{u}$.}

\statement{ana}{\textbf{Inputs from Note~S, continued.}\enspace
\textbf{(S5)} \emph{Frames.} For $\gamma\in\GT{T}$ and $t\in R^{\Ix{\fs}}$: $\ifpair{\gamma}^{*}\JT{T}t=0$ if and only if
the $e^*$-coordinates of $t$ and of $\gamma^{-1}t$ vanish. For $\gamma\in\GS{S}{E}$ and $p\in E^{\Mx{\fs}\sqcup\Mx{\fs}}$:
$p\in\link{\gamma}$ if and only if $q^{*}\lform{S}\,p=0$ for all $q\in\link{\gamma}$. If $\IsNF{j}{\gamma}$ then
$\IsNF{j}{\gamma^{-1}}$, and for every $y\in R^{\varrho+1}$ there are $x,y'$ with $\mathfrak f x=\gamma\mathfrak f y'$ and
$y'_a=y_a$ for $a\ge j$. If $\gamma\in\GS{S}{E}$ and $\mathfrak f x=\gamma\mathfrak f y$, then
$\Zf{S}{\gamma}\,y=-\NS{S}\,x_\varrho\,\epsilon_\varrho$ with $\epsilon_\varrho$ the last unit vector. If
$v\in E^{\Mx{\fs}}$ has $v_z=0$ and $v^{*}\Trat{S}{E}\,v=0$, then $v=0$. For $s,s'\in\GT{T}$:
$(s^{-1}s')\mathfrak f=\mathfrak f$ if and only if $s\mathfrak f=s'\mathfrak f$. And
$\Zr{\rinf{S}(\gamma)_\tau}_{ab}=\kapS{S}{\tau}^{\varrho-a}\kapS{S}{\tau}^{\varrho-b}\,\NS{S}^{-1}\,(\Zf{S}{\gamma}_{ab})_\tau$.

\textbf{(S6)} \emph{Discreteness.} For a finite family of indices and a compact set $\mathcal C$ of families of
complex matrices, only finitely many points of $\mathcal C$ are of the form
$(\text{the submatrix of }\rinf{S}(u)_\tau\text{ on these indices})_\tau$ with $u\in\clU{S}{c}$. For every subgroup
$\Gamma$ of $\clU{S}{c}$ there is a measurable $\beta\colon\prod_\tau\Gone{\fs}\to[0,\infty]$ with
$\sum_{u\in\Gamma}\beta(\rinf{S}(u)y)=1$ for all $y$.

\textbf{(S7)} \emph{Bessel systems.} For both sides $t$, all $n\in\UT{\Trg{\fs}{1}}$ and $y\in\Gone{\fs}$:
$\Wsys{t}{S}{\tau}(ny)=e^{\sd{n}-\cj{\sd{n}}}\,\Wsys{t}{S}{\tau}(y)$. There are a measurable
$A_\tau\colon\Gone{\fs}\to[0,\infty]$, invariant under left translation by $\UT{\Trg{\fs}{1}}$, and a finite measure
$m_\tau$ on $\Gone{\fs}$ such that, with $\theta_\tau(y)=1$ if $\Wsys{\full}{S}{\tau}(y)=\Wsys{\cut}{S}{\tau}(y)$ and
$\theta_\tau(y)=\bS{S}^{3\rr{\fs}^3\ninf{E}}$ otherwise:
$\theta_\tau(y)\,\|\Wsys{t}{S}{\tau}(y)\|\le A_\tau(y)\le m_\tau\{g:\InCut{yg}{\bS{S}}\}$ for both sides $t$ and all $y$,
and $m_\tau(\Gone{\fs})\int A_\tau\,d\muq{\fs}\le\bS{S}^{2\rr{\fs}^3}\,\#\Isys{S}{\tau}$. Moreover
$\|\besker{\sigma}{S}{\gamma}{y}\|=\prod_\tau\|\Wsys{\sigma}{S}{\tau}(y_\tau)\|\cdot\prod_\tau\|\Wsys{\sigma}{S}{\tau}((\rinf{S}(\gamma)^{-1}y)_\tau)\|$,
$y\mapsto\besker{\sigma}{S}{\gamma}{y}$ is continuous, and
$\|\amp{\sigma}{S}{c}{c'}{\gamma}\|\le\ampabs{\sigma}{S}{c}{c'}{\gamma}$.}

\statement{ana}{\textbf{Inputs from other units.} \textbf{(J1)}, unit~132 (the real place). The maps $y\mapsto y^{-1}\fr{a}$ and $y\mapsto\bsize{y}$ are continuous on $\Gone{\fs}$. For $y\in\Gone{\fs}$:
$\|y^{-1}\fr{a}\|\le\|y\|$; if $s$ lies in the group generated by $\UoT{\Trg{\fs}{1}}$ and $\HT{\Trg{\fs}{1}}$ then
$(sy)^{-1}\fr{a}=y^{-1}\fr{a}$; if $s\in\UT{\Trg{\fs}{1}}$ then $y^{-1}\fr{a}=\sum_b\fm{s}_{ba}\,(sy)^{-1}\fr{b}$; if
$s\in\ST{\Trg{\fs}{1}}$ then $\bsize{sy}=\bsize{y}$; there is $u\in\UT{\Trg{\fs}{1}}$ with $\|uy\|\le\bsize{y}$; for
complex numbers $c_a$, an index $i\in\Ix{\fs}$ and $s\in\Gone{\fs}$,
$\bigl|\sum_a c_a\,s(i,\fr{a})\bigr|\le\|sy\|\cdot\bigl\|\sum_a c_a\,y^{-1}\fr{a}\bigr\|$, and likewise for the sum
over $a<\varrho$ only; $|\langle v,\JT{\Trg{\fs}{1}}w\rangle|\le\|v\|\,\|w\|$; and for $x,y'\in\Gone{\fs}$,
$\Zr{x}_{ab}=\langle y'^{-1}\fr{a},\,\JT{\Trg{\fs}{1}}\,(x^{-1}y')^{-1}\fr{b}\rangle$. The coordinates of
$\Wsys{\full}{S}{\tau}$ are orthonormal: $\int\Wsys{\full}{S}{\tau}(y)_i\,\cj{\Wsys{\full}{S}{\tau}(y)_j}\,d\muq{\fs}(y)=\delta_{ij}$.

\textbf{(J2)}, units 113 and 57 (generalized index). $\gix{\cdot}$ is multiplicative on $\AEf$, zero on
non-units, positive on units and invariant under conjugation; $\gix{a}=\#\hO/a\hO$ for a unit $a\in\hO$;
$\gix{x}=|\NmE(x)|=\prod_\tau|x_\tau|^2$ for $x\in E$; in particular $\gix{n}=n^{2\ninf{E}}$ for a natural number
$n$. For an invertible integral matrix $M$ over $\AEf$ on a finite set $\Omega$, the index of $M\hO^{\Omega}$ in
$\hO^{\Omega}$ is $\gix{\det M}$ (unit~57).

\textbf{(J3)}, unit~114 (packing). In a real normed space $V$ of finite dimension, a set of points that are
pairwise at distance at least $m>0$ and of norm at most $r\ge0$ has at most $(1+2r/m)^{\dim V}$ elements.

\textbf{(J4)}, unit~121. An algebraic integer $x$ of a number field with $|x_\tau|<1$ at all infinite places
$\tau$ is $0$.

\textbf{(J5)}, unit~112 (duality). Over a field $F$, let $P,P'$ be matrices with the same finite set of rows
and the same finite set of columns, $\mathfrak a,\mathfrak a'$ sets of columns and $X,Y$ vectors. If every linear form $m$ that
kills the columns of $P$ in $\mathfrak a$ and the columns of $P'$ in $\mathfrak a'$ satisfies $m(PX)=m(P'Y)$, then
$Px=P'y$ for some $x$ equal to $X$ outside $\mathfrak a$ and some $y$ equal to $Y$ outside $\mathfrak a'$.

\textbf{(J6)}, unit~123. If $E$ is stably unramified, the integers of $E$ contain $\omega_0$ with
$\omega_0+\cj{\omega_0}=1$.

\noindent The catalogue of the write-up classes (J2) to (J6) as elementary and gives no reference for them.}

\subhead{Definitions of this note.} They are stated where they are first needed: the link metric and the link
discriminant (Definition~\ref{d:linkmetric}), integral lifts (Definition~\ref{d:lifts}), and the indices, gauges
and weights of the later estimates (Definitions \ref{d:pad} to \ref{d:weights}).

\sect{A positive form on links; class coordinates}{s:links}
The link of $\gamma$ is a Lagrangian subspace for the indefinite form $\lform{S}$ (Input~(S5)). To measure
lattices in links one needs a positive form; it is obtained by changing the sign of $\lform{S}$ on the second
copy and correcting the two $z$-directions.

\begin{definition}[link metric, link discriminant]\label{d:linkmetric}
Let $S\in\Fill{E}{\fs}$ and write $T=\Trat{S}{E}$. The \emph{link metric} is the matrix over $E$ on
$\Mx{\fs}\sqcup\Mx{\fs}$
\[\lmet{S}=\tfrac12\operatorname{diag}(T,T)+\NS{S}^{-1}\,\zfp{T}^{*}\,\zfp{T}.\]
For a finite set $\Omega$, a matrix $v$ over $E$ with rows $\Mx{\fs}\sqcup\Mx{\fs}$ and columns $\Omega$, and an
invertible matrix $e$ over $\AEf$ on $\Omega$, the \emph{link discriminant} is
\[\ldsq{S}{v}{e}=\bigl|\NmE\det\bigl(v^{*}\,\lmet{S}\,v\bigr)\bigr|\cdot\gix{\det e}^{2},\]
the squared discriminant, for the link metric, of the lattice with the adelic basis $\fadm{v}\,e$.
\end{definition}
\leanline{\lean{linkMetric}\sep\lean{linkDiscSq}}

\begin{definition}[integral lifts]\label{d:lifts}
Let $S\in\Fill{E}{\fs}$, $c\in\Cl{S}$, $\gamma\in\GS{S}{E}$, let $\Omega'$ be a set and $m$ a matrix over $\AEf$ with
rows $\Mx{\fs}\sqcup\Mx{\fs}$ and columns $\Omega'$. We write $\HIL{c}{\gamma}{m}$, and say that $m$ \emph{has
integral lifts} for $c$ and $\gamma$, if there are a finite set $\Omega$, a matrix~$t$ over $E$ with rows $\Ix{\fs}$
and columns $\Omega$, and a matrix $w$ over $\AEf$ with rows $\Omega$ and columns $\Omega'$, such that
\[\ifpair{\gamma}^{*}\,\ratG{S}\,t=0,\qquad\clC{S}{c}\,\fadm{t}\,w\ \text{ is integral},\qquad\fadm{\lmat{\gamma}\,t}\,w=m.\]
So the columns of $m$ are the images in the link of vectors of the lattice of the class $c$ that are adelic
combinations of rational vectors orthogonal to $X$ and to $\gamma X$.
\end{definition}
\leanline{\lean{HasIntegralLifts}}

\begin{lemma}[class coordinates and frames]\label{l:formulas}
Let $S\in\Fill{E}{\fs}$, $c,c'\in\Cl{S}$ and $\gamma\in\GS{S}{E}$.
\begin{enumerate}[label=(\alph*),itemsep=1pt]
\item For $k\in\GT{T}$, the two rows $(1,\varrho)$ and $(2,\varrho)$ of $\fpair{k}^{*}\JT{T}$ form the matrix
$\zfp{T}\,\lmat{k}$.
\item $\clC{S}{c}$ is block diagonal for the blocks $e$, middle, $e^*$: the identity on the outer blocks, and on
the middle block the matrix whose columns at $\kappa$ are those of $\clCX{S}{c}$ and whose column at $z$ is the
unit vector at $z$.
\item $\clCX{S}{c}^{*}\ \Tlat{S}{\AEf}\ \clCX{S}{c}=1$.
\item For matrices $t$, $t'$ over $E$ with rows $\Ix{\fs}$ and any sets of columns,
\[(\clC{S}{c}\fadm{t})^{*}\,\JT{\Tlat{S}{\AEf}}\,(\clC{S}{c}\fadm{t'})=\fadm{t^{*}\,\ratG{S}\,t'}.\]
\item $\clC{S}{c}\,\fadm{\fpair{\gamma}}=\fpair{k}$ and $\clC{S}{c}\,\fadm{\ifpair{\gamma}}=\ifpair{k}$ for
$k=\pos{S}{c}{c'}{\gamma}$.
\item If $k\in\GT{T}$ is integral over the subring $O$ of $R$, then so is $\fpair{k}$.
\item If a matrix $v$ over a field, with a finite set $\Omega'$ of rows and a finite set $\Omega$ of columns,
has a left inverse, then $\#\Omega$ is the dimension of the span of its columns.
\end{enumerate}
\end{lemma}
\begin{proof}
(a) The row of $\fpair{k}^{*}\JT{T}$ at $(1,\varrho)$ is the linear form $t\mapsto h_T(z,t)=(T\Mid t)_z$, and the
row at $(2,\varrho)$ is $t\mapsto h_T(kz,t)=h_T(z,k^{-1}t)=(T\Mid k^{-1}t)_z$, because $k$ is unitary. These are the
two rows of $\zfp{T}\lmat{k}$. (b) By \ref{r:groups}, $\clB{S}{c}$ is the product of two matrices that are the
identity on the outer blocks, with middle blocks $\operatorname{diag}(\tilde c,1)$ and the middle block of
$\latB{S}$. The inverse of the latter has the columns $\mbc{\fs}{\piS{S}}$ at $\kappa$ and the unit vector at $z$,
since $z$ belongs to the lattice basis. Hence the middle block of $\clC{S}{c}=\clB{S}{c}^{-1}$ has the columns
$\mbc{\fs}{\piS{S}}\tilde c^{-1}=\clCX{S}{c}$ at $\kappa$ and the unit vector at $z$. (c) Restrict the first
identity of Input~(S3) to the rows and columns of $\kappa$ and use (b): the block of $\ratG{S}$ on $\kappa$ is $1$.
(d) is the first identity of (S3) multiplied by $\fadm{t}^{*}$ and $\fadm{t'}$. (e) By (S3),
$k\mathfrak f=\clC{S}{c}\fadm{\gamma\mathfrak f}$, and $\clC{S}{c}\mathfrak f=\mathfrak f$ by (b); omitting the columns of
$z$ and $kz$ gives the second identity. (f) The entries of $\fpair{k}$ are entries of $1$ and of $k$. (g) The
matrix is injective, so its rank is $\#\Omega$.
\end{proof}
\leanline{\lean{framePairlasteqzFormPairmul}\sep\lean{classCoordeqfromBlocks}\sep\lean{gramclassCoordX}\sep
\lean{classCoordgram}\sep\lean{classCoordmulframePair}\sep\lean{classCoordmulisotropicFramePair}\sep
\lean{isIntMatrixframePair}\sep\lean{cardeqfinrankcolSpan}}

\begin{lemma}[the link is Lagrangian]\label{l:lagr}
Let $S\in\Fill{E}{\fs}$ and $\gamma\in\GS{S}{E}$.
\begin{enumerate}[label=(\alph*),itemsep=1pt]
\item $\dim_E\link{\gamma}=\kk{\fs}+1$.
\item If $v$ and $u$ are matrices over $E$ with rows $\Mx{\fs}\sqcup\Mx{\fs}$ and finite sets of columns
$\Omega'$ and $\Omega$, whose columns all lie in $\link{\gamma}$, then $v^{*}\,\lform{S}\,u=0$.
\end{enumerate}
\end{lemma}
\begin{proof}
By Input~(S5) the link is its own orthogonal for the hermitian form $\lform{S}$ on $E^{\Mx{\fs}\sqcup\Mx{\fs}}$;
this is (b). The form is non-degenerate, because $\Trat{S}{E}=\operatorname{diag}(1,\dots,1,-\NS{S})$ is invertible.
So the orthogonal of the link has dimension $2(\kk{\fs}+1)-\dim\link{\gamma}$, and (a) follows. (The formal proof
applies the rank theorem to $p\mapsto v^{*}\lform{S}p$ for a basis matrix $v$ of the link.)
\end{proof}
\leanline{\lean{finranklink}\sep\lean{gramlinkFormeqzerooflelink}}

\begin{lemma}[positivity of the link metric]\label{l:posdef}
Let $S\in\Fill{E}{\fs}$.
\begin{enumerate}[label=(\alph*),itemsep=1pt]
\item For every infinite place $\tau$ of $E$ the image of $\lmet{S}$ under $x\mapsto x_\tau$ is positive
definite.
\item If a matrix $v$ over $E$ with rows $\Mx{\fs}\sqcup\Mx{\fs}$ and a finite set $\Omega'$ of columns has a left
inverse, then $\det(v^{*}\,\lmet{S}\,v)\ne0$.
\end{enumerate}
\end{lemma}
\begin{proof}
(a) With $T=\Trat{S}{E}$, the matrix $\zfp{T}$ has the entry $-\NS{S}$ at the two positions (first row, first
$z$) and (second row, second $z$) and zeros elsewhere, so $\zfp{T}^{*}\zfp{T}$ is diagonal with entries
$\NS{S}^2$ at the two $z$ and $0$ elsewhere. Hence $\lmet{S}$ is the diagonal matrix with entries $\tfrac12$ at
the indices of $\kappa$ and $-\tfrac12\NS{S}+\NS{S}=\tfrac12\NS{S}$ at the two $z$: rational and positive. (b)
Under $x\mapsto x_\tau$ the matrix $v$ stays injective, so $v_\tau^{*}(\lmet{S})_\tau v_\tau$ is positive definite
by (a), and its determinant $\det(v^{*}\lmet{S}v)_\tau$ is positive.
\end{proof}
\leanline{\lean{linkMetricposDef}\sep\lean{detgramlinkMetricnezero}}

\begin{lemma}[products in two links]\label{l:gramlink}
Let $S\in\Fill{E}{\fs}$, $\gamma,\gamma'\in\GS{S}{E}$, and let $t$, $t'$ be matrices over $E$ with rows $\Ix{\fs}$
and sets of columns $\Omega$, $\Omega'$ such that $\ifpair{\gamma}^{*}\,\ratG{S}\,t=0$ and
$\ifpair{\gamma'}^{*}\,\ratG{S}\,t'=0$. Then for all $\alpha,\alpha'\in E$, with $T=\Trat{S}{E}$,
\[(\lmat{\gamma}t)^{*}\,\operatorname{diag}(\alpha T,-\alpha'T)\,(\lmat{\gamma'}t')
=\alpha\,t^{*}\,\ratG{S}\,t'-\alpha'\,t^{*}\,\ratG{S}\,(\gamma\gamma'^{-1})\,t'.\]
\end{lemma}
\begin{proof}
Let $t_i$ and $t'_l$ be columns of $t$ and $t'$. By Input~(S5) the $e^*$-coordinates of $t_i$, $\gamma^{-1}t_i$,
$t'_l$ and $\gamma'^{-1}t'_l$ vanish, and for two vectors $v,w$ without $e^*$-coordinates
$h_T(v,w)=(\Mid v)^{*}T\,\Mid w$. Since $\lmat{\gamma}t_i=(\Mid t_i,\Mid\gamma^{-1}t_i)$, the entry $(i,l)$ of the
left side is $\alpha\,h_T(t_i,t'_l)-\alpha'\,h_T(\gamma^{-1}t_i,\gamma'^{-1}t'_l)$, and
$h_T(\gamma^{-1}t_i,\gamma'^{-1}t'_l)=h_T(t_i,\gamma\gamma'^{-1}t'_l)$ because $\gamma$ is unitary.
\end{proof}
\leanline{\lean{gramlinkMatrixmul}}

\begin{lemma}[a Liouville inequality for Gram determinants]\label{l:liouville}
Let $E$ be stably unramified, let $\Omega$, $\Omega'$ be finite sets, $Q$ a matrix over $E$ on $\Omega'$ whose
image $Q_\tau$ is positive definite for every infinite place $\tau$, let $a,b$ be matrices over $E$ with rows
$\Omega'$ and columns $\Omega$ such that $\det(a^{*}Qb)\ne0$, and $v,v'$ invertible matrices over $\AEf$ on
$\Omega$. Let $x\ge0$ be a real number such that $x\le\gix{\delta}$ whenever
$\delta=\det\bigl(v^{*}\,\fadm{a^{*}Qb}\,v'\bigr)$ is a unit of $\AEf$. Then
\[x^2\ \le\ \gix{\det v}^2\cdot\bigl|\NmE\det(a^{*}Qa)\bigr|\cdot\bigl|\NmE\det(b^{*}Qb)\bigr|\cdot\gix{\det v'}^2.\]
\end{lemma}
\begin{proof}
$\delta=\cj{\det v}\cdot\det(a^{*}Qb)\cdot\det v'$ is a product of units, so $x\le\gix{\delta}$, and by (J2)
$\gix{\delta}=\gix{\det v}\,|\NmE\det(a^{*}Qb)|\,\gix{\det v'}$ with $|\NmE\det(a^{*}Qb)|=\prod_\tau|\det(a^{*}Qb)_\tau|^2$.
At every $\tau$ the block matrix with blocks $a_\tau^{*}Q_\tau a_\tau$, $a_\tau^{*}Q_\tau b_\tau$, $b_\tau^{*}Q_\tau a_\tau$,
$b_\tau^{*}Q_\tau b_\tau$ is positive semidefinite, being $(a_\tau\ b_\tau)^{*}Q_\tau(a_\tau\ b_\tau)$. For such a matrix
with square blocks $P_1$, $A_1$, $A_1^{*}$, $P_2$ one has $|\det A_1|^2\le\det P_1\cdot\det P_2$ (an elementary
input from unit~112; the artifact calls it the Cauchy--Schwarz inequality for determinants). Hence
$|\det(a^{*}Qb)_\tau|^{4}\le|\det(a^{*}Qa)_\tau|^{2}\,|\det(b^{*}Qb)_\tau|^{2}$, and the product over $\tau$ gives
$|\NmE\det(a^{*}Qb)|^2\le|\NmE\det(a^{*}Qa)|\cdot|\NmE\det(b^{*}Qb)|$. Squaring the bound for $x$ finishes the proof.
\end{proof}
\leanline{\lean{liouvillegram}}

\sect{Agreeing characters: proof of Theorem A}{s:agree}
\begin{proof}[Proof of Theorem A]
\emph{Step 1: reduction to congruences.} By \ref{r:char}, the hypothesis $\FixN{\redS{S}(k)}{x}$ gives a finite
set $\sA$ of pairs $p=(w_p,\epsilon_p)$ of a chosen place and $\epsilon_p\in\{0,1\}$, with
$x\le\prod_{p\in\sA}\#\Fw{w_p}$ (the norm of a chosen place is the number of elements of its residue field,
(S1)), and for every $p\in\sA$ a vector $v_p\ne0$ in $\Fw{w_p}^{\varrho+1}$ such that
$\ioia{S}{\epsilon_p}(\redS{S}(k)_{w_p})$ fixes $\mathfrak f v_p$.

Put $Q=t^{*}\fpair{k}=C^{*}$. Its entries lie in $\hO$, because $t$ is integral, $\hO$ is stable under the
conjugation, and $k\in\KS{S}$ is integral. Its columns are indexed by the pairs $(i,a)$; those with $a<\varrho$
vanish. Let $Q^{\circ}$ be the matrix on $\Omega$ whose column $l$ is the column $(\rho(l),\varrho)$ of $Q$, and
$\delta'=\det Q^{\circ}$. Then $\delta'=\cj{\delta}$, so $\delta'\in\hO$ is a unit of $\AEf$ and, by (J2),
$\gix{\delta}=\gix{\delta'}=\#\hO/\delta'\hO$. By (S1) the ring homomorphism
$\hO\to\prod_{p\in\sA}\Fw{w_p}$, $r\mapsto(\rag{S}{p}(r))_{p}$, is surjective; let $\mathfrak a$ be its kernel, so
that $\#\hO/\mathfrak a=\prod_{p\in\sA}\#\Fw{w_p}\ge x$. If $\rag{S}{p}(\delta')=0$ for every $p\in\sA$, then
$\delta'\hO\subseteq\mathfrak a$ and $\gix{\delta}=\#\hO/\delta'\hO\ge\#\hO/\mathfrak a\ge x$. So it suffices to
prove that $\rag{S}{p}(\delta')=0$ for every $p\in\sA$.

\emph{Step 2: the frames do not meet where $\delta'$ is a unit.} Fix $p=(w,\epsilon)\in\sA$, write
$\mathbb F=\Fw{w}$, and let a subscript $\mathbb F$ denote the image under $\rag{S}{p}$, entry by entry: for a matrix
$M$ over $\hO$, $\rF{M}$ is its entrywise image, and $\rF{(MM')}=\rF{M}\,\rF{M'}$. The field $\mathbb F$ has no
involution, so a conjugate transpose is always taken over $\hO$, before the image. Suppose that
$\rF{\delta'}=\det\rF{Q^{\circ}}\ne0$ in $\mathbb F$. By (S1), $\rF{k}\in\GL_{\Ix{\fs}}(\mathbb F)$ equals
$\ioia{S}{\epsilon}(\redS{S}(k)_{w})$, so $\rF{k}\,\mathfrak f v_p=\mathfrak f v_p$. We claim:
\[\text{if }\xi,\eta\in\mathbb F^{\varrho+1}\text{ satisfy }\mathfrak f\xi=\rF{k}\,\mathfrak f\eta,\text{ then }\xi=0.\tag{$*$}\]
For $\xi=\eta=v_p$ this contradicts $v_p\ne0$, and the theorem follows. We prove by descending induction on
$a\in\{\varrho,\dots,0\}$ that $\xi_a=0$ for all pairs $(\xi,\eta)$ as in $(*)$.

\emph{The case $a=\varrho$.} For such a pair, $\rF{Q}\,(\xi,-\eta)=\rF{(t^{*})}(\mathfrak f\xi-\rF{k}\,\mathfrak f\eta)=0$,
where $(\xi,-\eta)$ has the coordinate $\xi_a$ at $(1,a)$ and $-\eta_a$ at $(2,a)$. The columns $(i,a)$ of
$\rF{Q}$ with $a<\varrho$ vanish, so $\xi_\varrho$ times the column $(1,\varrho)$ equals $\eta_\varrho$ times the
column $(2,\varrho)$. These two columns are those of $\rF{Q^{\circ}}$, which is invertible. Hence
$\xi_\varrho=\eta_\varrho=0$.

\emph{The step.} Let $a<\varrho$, suppose the assertion holds for all indices above $a$, and let
$\mathfrak f\xi=\rF{k}\,\mathfrak f\eta$. Then $\xi_b=0$ for $b>a$ and $\eta_\varrho=0$. For
$\zeta\in\mathbb F^{\varrho+1}$ write $\zeta^{+}=(0,\zeta_0,\dots,\zeta_{\varrho-1})$. We shall prove that
\[m(\mathfrak f\xi^{+})=m(\rF{k}\,\mathfrak f\eta^{+})\tag{$**$}\]
for every linear form $m$ on $\mathbb F^{\Ix{\fs}}$ with $m(\fr{b})=0$ for $b\le a$ and $m(\rF{k}\fr{b})=0$ for
$b<\varrho$.
By the duality (J5), applied to the matrices $\mathfrak f$ and $\rF{k}\,\mathfrak f$, the sets of columns
$\{b\le a\}$ and $\{b<\varrho\}$ and the vectors $\xi^{+}$, $\eta^{+}$, there are then $\xi',\eta'$ with
$\mathfrak f\xi'=\rF{k}\,\mathfrak f\eta'$ and $\xi'_b=\xi^{+}_b$ for $b>a$. By the induction hypothesis at $a+1$,
$\xi_a=\xi^{+}_{a+1}=\xi'_{a+1}=0$.

\emph{Proof of $(**)$: integral lifts.} By (S1) there is a map $o\colon\mathbb F\to\hO$ with $o(0)=0$,
$\rag{S}{p}(o(\theta))=\theta$, and $\rag{S}{p'}(o(\theta))=0$ for every pair $p'\ne p$. Then, for every chosen
place $w'$, $\rag{S}{(w',\epsilon)}\bigl(\cj{o(\theta)}\bigr)=\rag{S}{(w',1-\epsilon)}(o(\theta))=0$. For a matrix $M$
over $\mathbb F$ on $\Ix{\fs}$ and $\zeta\in\mathbb F^{\varrho+1}$ put
\[\mathsf A(M,\zeta)=\sum_{b=1}^{\varrho}o\bigl(m(M\fr{b})\bigr)\,o(\zeta_{b-1})\in\hO,\quad\text{so that}\quad
\rag{S}{p}\bigl(\mathsf A(M,\zeta)\bigr)=m(M\mathfrak f\zeta^{+}).\]
Thus $(**)$ says that $\mathsf A(1,\xi)-\mathsf A(\rF{k},\eta)$ vanishes under $\rag{S}{p}$. By the two
equivalences of (S1) it suffices to show that
\[\sdc{S}\Bigl(r\,\gid{S}^{-1}\bigl(\mathsf A(1,\xi)-\mathsf A(\rF{k},\eta)\bigr)\Bigr)=1\qquad\text{for every }r\in\hO.\]

\emph{Eichler elements over $\hO$.} Fix $r\in\hO$. Let $a_0\le\varrho$, let $\zeta\in\mathbb F^{\varrho+1}$ have
$\zeta_b=0$ for $b\ge a_0$, and let $M$ be such that $m(M\fr{b})=0$ for $b<a_0$. We claim that there is
$u\in\KS{S}\cap\UT{\Tlat{S}{\AEf}}$ with $\sd{u}=r\,\mathsf A(M,\zeta)$ whose entrywise image under
$\rag{S}{(w',\epsilon)}$ is $1$ for every chosen place $w'\ne w$, and is $1+\rF{r}\,(\mathfrak f\zeta)(mM)$ for
$w'=w$; here $(\mathfrak f\zeta)(mM)$ is the product of a column and a row. Indeed, apply (S2) to the ring $\hO$
with its conjugation and the form $\Tlat{S}{\hO}$: its Gram matrix $J$ is hermitian and invertible over $\hO$ by
(S1), and $\hO$ contains $\omega_0$ with $\omega_0+\cj{\omega_0}=1$ by (J6). The linear form
$\mu=o\circ(mM)$ on $\hO^{\Ix{\fs}}$ is $v^{*}J$ for some $v$, it vanishes at $\fr{b}$ for $b<a_0$ since $o(0)=0$,
and $j=o\circ\zeta$ vanishes at $b\ge a_0$. So (S2) gives $u$ unitary over $\hO$, in the unipotent group, with
$\sd{u}=r\sum_{b\ge1}\mu(\fr{b})\,j_{b-1}=r\,\mathsf A(M,\zeta)$ and with matrix
$1+\xi_0\,\mu-(v+\varsigma\xi_0)\,\xi_0^{*}J$, where $\xi_0=r\,\mathfrak f j$. It is integral with integral inverse,
so it lies in $\KS{S}$. Under $\rag{S}{(w',\epsilon)}$ the entries $\cj{r}\,\cj{o(\zeta_b)}$ of $\xi_0^{*}$ go to $0$,
while $\xi_0$ goes to $\rF{r}\,\mathfrak f\zeta$ and $\mu$ to $mM$ if $w'=w$, and $\xi_0$ goes to $0$ if $w'\ne w$.

\emph{Conclusion.} Let $u$ be the element of the claim for $(a_0,\zeta,M)=(a+1,\xi,1)$ and $u'$ the one for
$(\varrho,\eta,\rF{k})$; the hypotheses hold by the properties of $\xi$, $\eta$ and $m$. We compare
$\redS{S}(u)\,\redS{S}(k)$ and $\redS{S}(k)\,\redS{S}(u')$ in $\GRes{S}$. At the component of a chosen place $w'$,
after the automorphism $\ioia{S}{\epsilon}$, these are by (S1) the entrywise images of $uk$ and $ku'$ under
$\rag{S}{(w',\epsilon)}$. For $w'\ne w$ both are the image of $k$. For $w'=w$ they are
\begin{gather*}\bigl(1+\rF{r}\,(\mathfrak f\xi)\,m\bigr)\rF{k}=\rF{k}+\rF{r}\,(\mathfrak f\xi)(m\rF{k})\quad\text{and}\\
\rF{k}\bigl(1+\rF{r}\,(\mathfrak f\eta)(m\rF{k})\bigr)=\rF{k}+\rF{r}\,(\rF{k}\,\mathfrak f\eta)(m\rF{k}),\end{gather*}
which agree because $\mathfrak f\xi=\rF{k}\,\mathfrak f\eta$. Hence $\redS{S}(u)\,\redS{S}(k)=\redS{S}(k)\,\redS{S}(u')$,
and since $\redS{S}(k)$ agrees, $\uchar{S}{u}=\uchar{S}{u'}$, that is
$\sdc{S}(\gid{S}^{-1}r\,\mathsf A(1,\xi))=\sdc{S}(\gid{S}^{-1}r\,\mathsf A(\rF{k},\eta))$. As $\sdc{S}$ is an
additive character, this is the required identity.
\end{proof}
\leanline{\lean{legenIndexdetoffixesNormAtLeast}}

\sect{Integral Eichler elements of the arithmetic unipotent group}{s:eichler}
Rational elements of $\clU{S}{c}$ with prescribed entries are produced by Eichler transformations. A position
$(a,l)$, $a<\varrho$, $l\in\Ix{\fs}$, is called \emph{free} if $l=e_b$ or $l=e^*_b$ with $b>a$, or if $l$ is a
middle index (this is a notion of unit~118); and $[P]$ denotes $1$ if the statement $P$ holds and $0$ otherwise.

\begin{lemma}[Eichler elements with several columns]\label{l:eichler}
In Setting~\ref{set:su}, let $c\in\Cl{S}$, let $\Omega$ be a finite set, $p\colon\Omega\to\{0,\dots,\varrho-1\}$ a
map, and $x$ the matrix with rows $\Ix{\fs}$ and columns $\Omega$ whose column $b$ is the unit vector $e_{p(b)}$.
Let $u$ be a matrix over $E$ with rows $\Ix{\fs}$ and columns $\Omega$, and $s$ a matrix over $E$ on $\Omega$, such
that
\begin{enumerate}[label=(\roman*),itemsep=1pt]
\item $u(e^*_a,b)=0$ for all $a$, $b$, $b'$ with $a\le p(b')$;
\item $s^{*}=-s$;
\item the matrices $\clC{S}{c}\,\fadm{u}$ and $\fadm{s}$ are integral.
\end{enumerate}
Then there are $\gamma\in\clU{S}{c}$ and $\lambda\in E$ with
\[\gamma=1+\gS{S}\Bigl(x\,u^{*}\ratG{S}-\bigl(u+x\,(\lambda\,u^{*}\ratG{S}\,u-s)\bigr)\,x^{*}\ratG{S}\Bigr),\]
and if $(u^{*}\ratG{S})(b,\fr{a})=0$ for all $b\in\Omega$ and all $a\le\varrho$, then $\gamma\in\clUo{S}{c}$.
\end{lemma}
\begin{proof}
Write $J=\ratG{S}$ and $g=\gS{S}$, take $\omega_0$ from (J6), and put $\lambda=g\,\omega_0$,
$B=\lambda\,u^{*}Ju-s$ and $N=x\,u^{*}J-(u+xB)\,x^{*}J$, so that the matrix in question is $1+gN$. The columns of
$x$ are isotropic and mutually orthogonal, and by (i) they are orthogonal to the columns of $u$:
$x^{*}Jx=0$ and $x^{*}Ju=0$. Moreover $x^{*}J$ vanishes at all frame vectors, and $\clC{S}{c}\,\fadm{x}=\fadm{x}$ by
Lemma~\ref{l:formulas}(b).

\emph{Unitarity.} $(1+gN)^{*}J(1+gN)=J$ means $N^{*}J+JN+g\,N^{*}JN=0$. Now
$JN+N^{*}J=-Jx\,(B+B^{*})\,x^{*}J$, the terms $Jxu^{*}J$ and $Jux^{*}J$ cancelling; and in $N^{*}JN$ every product
contains $x^{*}Jx$, $x^{*}Ju$ or $u^{*}Jx$, except $Jx\,(u^{*}Ju)\,x^{*}J$. Since $s^{*}=-s$ and $u^{*}Ju$ is hermitian,
$B+B^{*}=(\lambda+\cj{\lambda})\,u^{*}Ju=g\,u^{*}Ju$. So the sum vanishes, and $\gamma=1+gN\in\GS{S}{E}$.

\emph{$\gamma$ lies in the unipotent group.} The row $b$ of $u^{*}J$ has the entry $\cj{u(e^*_a,b)}$ at $e_a$,
which by (i) vanishes unless $a>p(b')$ for all $b'$. So the entries of $x\,u^{*}J$ sit in rows $e_{p(b')}$ and in
columns of level greater than $p(b')$. The matrix $x^{*}J$ has its entries in the columns $e^*_{p(b)}$, of level
$2\varrho-p(b)\ge\varrho+1$, and the rows of $u+xB$ that do not vanish are rows $e_a$, middle rows, or rows
$e^*_a$ with $a>p(b)$ for all $b$, all of level less than $2\varrho-p(b)$. Hence $N$ is strictly upper
triangular for the levels and $\gamma\in\UT{\Trat{S}{E}}$.

\emph{The position.} By \ref{r:groups}, $\pos{S}{c}{c}{\gamma}=1+\gid{S}\,\clC{S}{c}\fadm{N}\clC{S}{c}^{-1}$, and by
(S3), $\fadm{J}\,\clC{S}{c}^{-1}=\clC{S}{c}^{*}\,J'$ with $J'=\JT{\Tlat{S}{\AEf}}$. Therefore, writing
$\hat x=\clC{S}{c}\fadm{x}$ and $\hat u=\clC{S}{c}\fadm{u}$,
\[\clC{S}{c}\,\fadm{N}\,\clC{S}{c}^{-1}=\bigl(\hat x\,\hat u^{*}-(\hat u+\hat x\,\fadm{B})\,\hat x^{*}\bigr)\,J',
\qquad\fadm{B}=\fadm{\lambda}\,\hat u^{*}J'\hat u-\fadm{s},\]
the last by Lemma~\ref{l:formulas}(d). All matrices on the right are integral, by (iii) and (S1). So
$\gid{S}^{-1}(\pos{S}{c}{c}{\gamma}-1)$ is integral; then $\pos{S}{c}{c}{\gamma}$ is integral, and so is its inverse
$J'^{-1}\pos{S}{c}{c}{\gamma}^{*}J'$. By (S1) the position lies in $\KgS{S}$, that is $\gamma\in\clU{S}{c}$.

\emph{The frame.} If $u^{*}J$ vanishes at the frame vectors, then $N\fr{a}=0$ for all $a$, so
$\gamma\mathfrak f=\mathfrak f$ and $\gamma\in\clUo{S}{c}$.
\end{proof}
\leanline{\lean{existsmemclassUnip}}

\begin{lemma}[Eichler elements in one row]\label{l:eichlerrow}
In Setting~\ref{set:su}, let $c\in\Cl{S}$ and $a<\varrho$, let $l$ be a row vector over $E$ indexed by $\Ix{\fs}$
with $l_j=0$ unless $(a,j)$ is a free position, such that the row vector $\fadm{l}\,\clB{S}{c}$ is integral, and
let $s$ be an integer of $E$ with $\cj{s}=-s$. Then there is $\gamma\in\clU{S}{c}$ such that for all $b<\varrho$
and all $j\in\Ix{\fs}$ with $j\ne e^*_a$ or $l=0$,
\[\gamma(e_b,j)=\delta_{e_b,j}+\gS{S}\,[b=a]\,\bigl(l_j+[j=e^*_a]\,s\bigr).\]
\end{lemma}
\begin{proof}
Since $J=\ratG{S}$ is invertible there is a column vector $v$ with $v^{*}J=l$. We apply
Lemma~\ref{l:eichler} with $\Omega$ a single point, $p=a$, $x=e_a$, the column $u=v$ and the matrix $(s)$.
Condition (i): $\cj{v(e^*_{a'})}=l_{e_{a'}}=0$ for $a'\le a$, because $(a,e_{a'})$ is not free. Condition (ii) is
$\cj{s}=-s$. Condition (iii): $v=J^{-1}l^{*}$, and from (S3), $\clC{S}{c}\fadm{J}^{-1}=J'^{-1}\clB{S}{c}^{*}$ with
$J'=\JT{\Tlat{S}{\AEf}}$; so $\clC{S}{c}\fadm{v}=J'^{-1}(\fadm{l}\,\clB{S}{c})^{*}$, which is integral by hypothesis
and (S1). The lemma gives $\gamma\in\clU{S}{c}$ and $\lambda$ with
$\gamma=1+\gS{S}\bigl(e_a\,l-(v+e_a(\lambda\,v^{*}Jv-s))\,e_a^{*}J\bigr)$. The row vector $e_a^{*}J$ is the unit row
at $e^*_a$. So for $j\ne e^*_a$ the entry $\gamma(e_b,j)$ is $\delta_{e_b,j}+\gS{S}[b=a]\,l_j$; and if $l=0$ then
$v=0$ and $\gamma(e_b,e^*_a)=\gS{S}[b=a]\,s$.
\end{proof}
\leanline{\lean{existseichlerclassUnip}}

\sect{Covering lemmas; indices, gauges and weights}{s:count}
The first four lemmas are elementary counting tools; of them this note uses only the first. The definitions that
follow, and the two lemmas after them, prepare the later estimates of the unit.

\begin{lemma}[colouring]\label{l:colour}
Let $\sP$ and $\sC$ be sets, $\sA\subseteq\sP$, let $\bowtie$ be a relation on $\sP$, $m,x\in[0,\infty]$, and
$\sw\colon\sP\to[0,\infty]$. Let $\col\colon\sA\to\sC$ be a map whose image has at most $m$ elements and such that
$\col(p)=\col(q)$ implies $p\bowtie q$. Suppose that $\sum_{p\in\sB}\sw_p\le x$ for every non-empty
$\sB\subseteq\sA$ with $p\bowtie q$ for all $p,q\in\sB$. Then $\sum_{p\in\sA}\sw_p\le m\cdot x$.
\end{lemma}
\begin{proof}
The sum over $\sA$ is the sum, over the colours in the image of $\col$, of the sums over the fibres. Each fibre
is non-empty and consists of pairwise related elements, so its sum is at most $x$; and there are at most $m$
colours.
\end{proof}
\leanline{\lean{tsumleofencardrangele}}

\begin{lemma}[colouring by cosets]\label{l:colourcoset}
Let $\sP$ be a set, $\sA\subseteq\sP$, $\Omega$ a finite set, $\mathsf M$ an abelian group with subgroups
$\mathsf L$, $\mathsf L'$ such that the index of $\mathsf L\cap\mathsf L'$ in $\mathsf L'$ is finite, and
$f\colon\sP\to\mathsf M^{\Omega}$ a map with $f(p)_i\in\mathsf L'$ for all $p\in\sA$ and $i\in\Omega$. Let
$\sw\colon\sP\to[0,\infty]$ and $x\in[0,\infty]$, and suppose that $\sum_{p\in\sB}\sw_p\le x$ for every non-empty
$\sB\subseteq\sA$ such that $f(p)_i-f(q)_i\in\mathsf L$ for all $p,q\in\sB$ and $i\in\Omega$. Then
\[\sum_{p\in\sA}\sw_p\ \le\ [\mathsf L':\mathsf L\cap\mathsf L']^{\#\Omega}\cdot x.\]
\end{lemma}
\begin{proof}
Colour $p\in\sA$ by the family of the classes of the $f(p)_i$ in $\mathsf L'/(\mathsf L\cap\mathsf L')$. There are
at most $[\mathsf L':\mathsf L\cap\mathsf L']^{\#\Omega}$ colours, and two elements of one colour have
$f(p)_i-f(q)_i\in\mathsf L$ for all $i$. Lemma~\ref{l:colour} applies.
\end{proof}
\leanline{\lean{tsumleofcolouring}}

\begin{lemma}[bounded integral vectors]\label{l:intcount}
Let $E'$ be a number field with a chosen complex embedding $x\mapsto x_\tau$ at each of its infinite places
$\tau$, let $\sP$ be a set, $\sA\subseteq\sP$, $\Omega$ a finite set, $r\ge0$ a real number, and
$F\colon\sP\to E'^{\,\Omega}$ a map that is injective on $\sA$ and such that, for all $p\in\sA$ and $i\in\Omega$,
$F(p)_i$ is integral over $\mathbb Z$ and $|(F(p)_i)_\tau|\le r$ for all $\tau$. Then
$\#\sA\le(1+2r)^{\,2\,\#\Omega\,\ninf{E'}}$, where $\ninf{E'}$ is the number of infinite places of $E'$.
\end{lemma}
\begin{proof}
Map $p\in\sA$ to the point $((F(p)_i)_\tau)_{i,\tau}$ of the real vector space $\CC^{\Omega\times\{\tau\}}$, of
dimension $2\,\#\Omega\,\ninf{E'}$, with the supremum norm. All these points have norm at most $r$. Two of them at
distance less than $1$ belong to $p,q$ for which every $F(p)_i-F(q)_i$ is an algebraic integer with all
absolute values less than $1$, hence $0$ by (J4), so $p=q$ by injectivity. The packing bound (J3) with $m=1$
gives the claim.
\end{proof}
\leanline{\lean{encardintegralle}}

\begin{lemma}[index of a scaled lattice]\label{l:scaled}
Let $E'$ be a CM field that is stably unramified, $\Omega$ a finite set, and $c$ a finite idele of $E'$ that lies
in the closure $\hat O_{E'}$ of the integers. Then the index of $c\,\hat O_{E'}^{\Omega}$ in $\hat O_{E'}^{\Omega}$
is $\gix{c}^{\#\Omega}$; in particular it is finite and not zero.
\end{lemma}
\begin{proof}
This is (J2) for the invertible integral matrix $c\cdot1$, of determinant $c^{\#\Omega}$; and $\gix{c}>0$.
\end{proof}
\leanline{\lean{submoduleIndexsmulintVectors}\sep\lean{submoduleIndexsmulnezero}}

\subhead{Indices, gauges and weights.} In the following definitions an index $[\mathsf L':\mathsf L\cap\mathsf L']$
of subgroups of a group is read as $0$ when it is infinite, as in the formal text, and $\colL{M}$ denotes the
$\hO$-module generated by the columns of a matrix $M$ over $\AEf$.

\begin{definition}[rows, padded matrices, middle parts]\label{d:pad}
Let $j\le\varrho+1$ be a level. $\IdxB{j}$ is the set of the $a\in\{0,\dots,\varrho-1\}$ with $a<j$. For a matrix
$Z$ with rows $0,\dots,\varrho$ and $\varrho+1$ columns, $\xrows{Z}$ is its submatrix of the rows
$0,\dots,\varrho-1$. For a square matrix $Z$ of size $\varrho+1$ over a commutative ring, $\padl{j}{Z}$ is the
matrix whose rows $a<j$ are those of $Z$ and whose other rows are those of the identity. For a square matrix
$\gamma$ on $\Ix{\fs}$, $\fmid{\gamma}$ is its submatrix of the rows in $\Mx{\fs}$ and the columns $\fr{0},\dots,\fr{\varrho}$:
the middle parts of the vectors $\gamma\fr{a}$.
\end{definition}
\leanline{\lean{IndexBelow}\sep\lean{xRows}\sep\lean{paddedLeading}\sep\lean{frameMiddle}}

\begin{lemma}\label{l:paddet}
For a level $j\le\varrho+1$ and a square matrix $Z$ of size $\varrho+1$ over a commutative ring $R$:
$\det\padl{j}{Z}=\lmin{j}{Z}$.
\end{lemma}
\begin{proof}
Expanding along the rows $a\ge j$, which are rows of the identity, leaves the determinant of the upper left
block of size $j$.
\end{proof}
\leanline{\lean{detpaddedLeading}}

\begin{definition}[lattice discriminants and indices]\label{d:indices}
Let $E$ be stably unramified, $S\in\Fill{E}{\fs}$ and $c\in\Cl{S}$; write $T'=\Tlat{S}{\AEf}$ and, for
$k\in\GFin{S}$, $Z'(k)=\ZT{T'}{k}$.
\begin{enumerate}[label=(\alph*),itemsep=1pt]
\item $\latd{S}{c}$ is the function on the subspaces $V$ of $E^{\kappa}$ given by
$\latd{S}{c}(V)=\sld{\clCX{S}{c}}{V}$, the discriminant of \ref{r:links} for the matrix $H=\clCX{S}{c}$, which
has the left inverse $\clCX{S}{c}^{*}T'$ by Lemma~\ref{l:formulas}(c). It is the discriminant of the lattice cut
out by $V$ in the $X_0$-part of the lattice of the class $c$.
\item For $k\in\GFin{S}$, a finite idele $a$ and a subspace $V$ of $E^{\kappa}$,
\[\fvi{S}{c}{k}{a}{V}=[\mathsf L_2:a\,\mathsf L_1\cap\mathsf L_2],\]
where\quad $\mathsf L_1=\colL{(\xrows{Z'(k)})^{*}}$\quad and\quad
$\mathsf L_2=\colL{{\fmid{k}}^{*}\,T'\,\clCX{S}{c}\,\slb{\clCX{S}{c}}{V}}$.
\item For a level $j\le\varrho+1$ and $k\in\GFin{S}$, $\xri{S}{j}{k}=[\mathsf L_1:\mathsf L_1^{(j)}\cap\mathsf L_1]$,
where $\mathsf L_1$ is as in (b) and $\mathsf L_1^{(j)}$ is generated by the columns of the conjugate transpose
of the submatrix of $Z'(k)$ of the rows $a\in\IdxB{j}$.
\item For $k\in\GFin{S}$, $\sati{S}{k}$ is the index of $\colL{\ifpair{k}}$ in the module of the vectors of
$\hO^{\Ix{\fs}}$ that lie in the $\AEf$-span of the columns of $\ifpair{k}$.
\end{enumerate}
\end{definition}
\leanline{\lean{latticeDisc}\sep\lean{frameValueIndex}\sep\lean{xRowsIndex}\sep\lean{saturationIndex}}

In (b), $\mathsf L_1$ is spanned by the vectors $(\cj{h(e_i,k\fr{l})})_l$, $i<\varrho$, and $\mathsf L_2$ by the
vectors $(h(k\fr{l},y))_l$ for $y$ in the basis of the lattice cut out by $V$; in (d), the columns of $\ifpair{k}$
are the $e_a$ and $ke_a$.

\begin{definition}[invariants of pairs, gauges, near representatives]\label{d:gauge}
Let $E$ be stably unramified, $S\in\Fill{E}{\fs}$, $c,c'\in\Cl{S}$, $\gamma,\gamma'\in\GS{S}{E}$ and $j\le\varrho+1$
a level.
\begin{enumerate}[label=(\alph*),itemsep=1pt]
\item For a pair $p=(c',\gamma)$: $\telem{S}{c}{p}=\pos{S}{c}{c'}{\gamma}$.
\item $\pinv{S}{c}{c'}{\gamma}{\gamma'}=\latd{S}{c}(V)\cdot\fvi{S}{c}{k}{1}{V}\cdot\sati{S}{k}$ with
$V=\clink{\gamma}{\gamma'}$ and $k=\pos{S}{c}{c'}{\gamma}$.
\item $\mnorm{S}{j}{\gamma}=\bigl|\NmE\lmin{j}{\Zf{S}{\gamma}}\bigr|\cdot\NS{S}^{-2\,[j=\varrho+1]\,\ninf{E}}$.
\item $\cgauge{S}{j}{c}{c'}{\gamma}=\max\bigl(\tdnorm{S}{\gamma},\ \NS{S}^{\ninf{E}}\,\mnorm{S}{j}{\gamma}\bigr)\big/
\bigl(\xri{S}{j}{k}\cdot\xri{S}{j}{k^{-1}}\bigr)$ with $k=\pos{S}{c}{c'}{\gamma}$.
\item $\IsNRep{S}{c}{c'}{\gamma}$ means that $\IsNr{S}{\gamma}$ and that $\gamma$ is the chosen representative of
its double coset $\clUo{S}{c}\,\gamma\,\clUo{S}{c'}$ (one representative is chosen in every double coset).
\end{enumerate}
\end{definition}
\leanline{\lean{termElem}\sep\lean{pairInvariant}\sep\lean{minorNorm}\sep\lean{classGauge}\sep\lean{IsNearRep}}

\begin{definition}[weights and labelled terms]\label{d:weights}
Let $E$ be stably unramified, $S\in\Fill{E}{\fs}$, $c\in\Cl{S}$, $j\le\varrho+1$ a level, and write
$k=\kk{\fs}$.
\begin{enumerate}[label=(\alph*),itemsep=1pt]
\item For a function $n\colon\GFin{S}\to\mathbb N$ and a pair $p=(c',\gamma)\in\Cl{S}\times\GS{S}{E}$:
$\ipw{S}{c}{n}{p}=n\bigl(\telem{S}{c}{p}\bigr)^{-k}\in[0,\infty)$, read as $0$ when $n(\telem{S}{c}{p})=0$.
\item For a natural number $i$, pairs $p,p'\in\Cl{S}\times\GS{S}{E}$ and a subspace $V$ of $E^{\kappa}$:
\begin{multline*}\termw{S}{j}{i}{c}{p}{p'}{V}=\cw{S}{i}^{-1}\cdot\latd{S}{c}(V)^{-(k+1)}\\
\cdot\ipw{S}{c}{n_{j,V}}{p}\cdot\ipw{S}{c}{\sati{S}{\cdot}}{p'},\end{multline*}
where $n_{j,V}(k')=[\hO^{\IdxB{j}}:\gid{S}\hO^{\IdxB{j}}]\cdot\fvi{S}{c}{k'}{1}{V}$ for $k'\in\GFin{S}$, and an inverse
of $0$ is read as $0$.
\item For $c'\in\Cl{S}$ and $\gamma,\gamma'\in\GS{S}{E}$, with $V=\clink{\gamma}{\gamma'}$:
\[\pairw{S}{j}{c}{c'}{\gamma}{\gamma'}=\termw{S}{j}{k-\dim V}{c}{(c',\gamma)}{(c',\gamma')}{V}\]
if $\Cong{S}{c}{c'}{\gamma}{\gamma'}$, $\IsNRep{S}{c}{c'}{\gamma}$, $\HasLev{S}{j}{\gamma}$, $\IsNRep{S}{c}{c'}{\gamma'}$
and $\HasLev{S}{j}{\gamma'}$ all hold (the difference $k-\dim V$ being $0$ if $\dim V>k$), and
$\pairw{S}{j}{c}{c'}{\gamma}{\gamma'}=0$ otherwise.
\item For $t\in\redim{S}{\ST{\Tlat{S}{\AEf}}}$ and a pair $p=(c',\gamma)$: $\IsLab{S}{j}{c}{t}{p}$ means that
$\IsNRep{S}{c}{c'}{\gamma}$, $\HasLev{S}{j}{\gamma}$, and $\gamma\in\terms{S}{c}{s}{c'}$ for some $s\in\GRes{S}$ whose
component in the subgroup $\redim{S}{\ST{\Tlat{S}{\AEf}}}$, for a fixed bijection of $\GRes{S}$ with the product of
the set of cosets of this subgroup and the subgroup, is $t$.
\end{enumerate}
\end{definition}
\leanline{\lean{invPowWeight}\sep\lean{termWeight}\sep\lean{pairWeight}\sep\lean{IsLabelledTerm}}

By Lemma~\ref{l:scaled} and (J2), $[\hO^{\IdxB{j}}:\gid{S}\hO^{\IdxB{j}}]=\gS{S}^{2\ninf{E}\min(j,\varrho)}$.

\begin{lemma}[terms of level $j$]\label{l:levels}
Let $S\in\Fill{E}{\fs}$, let $j\le\varrho+1$ be a level and $\gamma\in\GS{S}{E}$ with $\HasLev{S}{j}{\gamma}$. Then:
\begin{enumerate}[label=(\alph*),itemsep=1pt]
\item $\HasLev{S}{j}{\gamma^{-1}}$;
\item $\mnorm{S}{j}{\gamma}>0$;
\item $\padl{j}{\Zf{S}{\gamma}}$ is invertible: $\padl{j}{\Zf{S}{\gamma}}\cdot\padl{j}{\Zf{S}{\gamma}}^{-1}=1$;
\item for every frame index $b$ with $j\le b$ there is $w\in E^{\varrho+1}$ with
\[\mathfrak f w=\gamma\mathfrak f\cdot\bigl(\text{column $b$ of }\padl{j}{\Zf{S}{\gamma}}^{-1}\bigr):\]
this combination of the vectors $\gamma\fr{a}$ lies in the span of the frame.
\end{enumerate}
\end{lemma}
\begin{proof}
(b) and (c) hold because $\lmin{j}{\Zf{S}{\gamma}}\ne0$ and by Lemma~\ref{l:paddet}. (a) The normal form passes
to $\gamma^{-1}$ by (S5). Since $\gamma$ is unitary, $\Zf{S}{\gamma^{-1}}=\Zf{S}{\gamma}^{*}$, so the leading minors of
$\Zf{S}{\gamma^{-1}}$ are the conjugates of those of $\Zf{S}{\gamma}$, and
$\tdisc{S}{\gamma^{-1}}=\tdisc{S}{\gamma}$. (d) By (S5), applied to the unit vector $\epsilon_b$, there are $w$ and
$y'$ with $\mathfrak f w=\gamma\mathfrak f y'$ and $y'_a=\delta_{ab}$ for $a\ge j$. By (S5) again,
$\Zf{S}{\gamma}y'$ is a multiple of the last unit vector; as $b\le\varrho$ forces $j\le\varrho$, the rows $a<j$
of $\Zf{S}{\gamma}y'$ vanish. So $\padl{j}{\Zf{S}{\gamma}}\,y'=\epsilon_b$, and $y'$ is the column $b$ of the inverse.
\end{proof}
\leanline{\lean{inv}\sep\lean{minorNormpos}\sep\lean{paddedLeadingmulinv}\sep\lean{existssharedframevector}}

\begin{lemma}\label{l:xrows}
Let $k\in\GT{T}$ and let $n$ be a matrix over $R$ with rows $\Ix{\fs}$ and a set $\Omega$ of columns. If
$\ifpair{k}^{*}\,\JT{T}\,n=0$ and $\zfp{T}\,\lmat{k}\,n=0$, then
\[(\xrows{\ZT{T}{k}})^{*}\cdot(-n_e)={\fmid{k}}^{*}\,T\,n_{\Mid},\]
where $n_e$ and $n_{\Mid}$ are the submatrices of $n$ of the rows $e_0,\dots,e_{\varrho-1}$ and of the middle rows.
\end{lemma}
\begin{proof}
Let $n_l$ be a column of $n$ and $b\le\varrho$. The entries $(b,l)$ of the two sides are
\[-\sum_{a<\varrho}h_T(k\fr{b},e_a)\,n_l(e_a)\qquad\text{and}\qquad(\Mid k\fr{b})^{*}\,T\,\Mid n_l.\]
By (S5) the $e^*$-coordinates of $n_l$ vanish, so
\[h_T(k\fr{b},n_l)=\sum_{a<\varrho}h_T(k\fr{b},e_a)\,n_l(e_a)+(\Mid k\fr{b})^{*}T\,\Mid n_l.\]
Finally
$h_T(k\fr{b},n_l)=0$: for $b<\varrho$ by the first hypothesis, and for $b=\varrho$ by the second and
Lemma~\ref{l:formulas}(a).
\end{proof}
\leanline{\lean{xRowsconjTransposemuleq}}

\sect{Box partitions and the frame count}{s:boxes}
Throughout this section $\mathcal U=\UT{\Trg{\fs}{1}}$ is the unipotent group of $\Gone{\fs}$, and in Setting~\ref{set:adm}
we abbreviate $b=\bS{S}$, $n=\nn{\fs}$, $N=\NS{S}$, $f=\ninf{E}$; these are abbreviations of numbers, not of
functions.

\begin{lemma}[frame vectors of two elements]\label{l:framediff}
Let $S\in\Fill{E}{\fs}$ and let $\gamma,\gamma'\in\GS{S}{E}$ be such that $\gamma\mathfrak f$ and $\gamma'\mathfrak f$ have the
same rows outside the rows $e_0,\dots,e_{\varrho-1}$. Then for every $y\in\Gone{\fs}$, every infinite place $\tau$ of
$E$ and every frame index $l$,
\[\bigl(\rinf{S}(\gamma)_\tau^{-1}y\bigr)^{-1}\fr{l}-\bigl(\rinf{S}(\gamma')_\tau^{-1}y\bigr)^{-1}\fr{l}
=\sum_{a<\varrho}\bigl((\gamma-\gamma')(e_a,\fr{l})\bigr)_\tau\,\frac{\bscale{S}{\tau}{\fr{l}}}{\bscale{S}{\tau}{e_a}}\ y^{-1}\fr{a}.\]
\end{lemma}
\begin{proof}
$(\rinf{S}(\gamma)_\tau^{-1}y)^{-1}\fr{l}=y^{-1}\rinf{S}(\gamma)_\tau\fr{l}=\sum_{i\in\Ix{\fs}}\rinf{S}(\gamma)_\tau(i,\fr{l})\,y^{-1}e_i$,
with $e_i$ the unit vectors; by \ref{r:groups} the coefficient is $(\gamma(i,\fr{l}))_\tau\,\bscale{S}{\tau}{\fr{l}}/\bscale{S}{\tau}{i}$.
In the difference only the rows $i=e_a$ remain, and $y^{-1}e_a=y^{-1}\fr{a}$.
\end{proof}
\leanline{\lean{frameVectorsubeqsum}}

\statement{ana}{\textbf{Input.} \textbf{(J7)}, unit~99 (classed as elementary by the catalogue, without a
reference). Let $d\ne0$ lie in the different of $E$. For every family $(z_\tau)$ of complex numbers there is an
integer $x$ of $E$ with $|z_\tau-x_\tau|\le(2\,\ninf{E})^2\,|d_\tau|$ for all infinite places $\tau$. (The formal
statement also prescribes $x$ modulo $\hO$ by a finite adele; it is used with the adele $0$.)}

\begin{proposition}[box partitions exist]\label{p:boxexists}
In Setting~\ref{set:adm}, for every $c\in\Cl{S}$ the chosen function $\omS{S}{c}$ of \ref{r:cutoff} is a box
partition for $c$: $\IsBP{S}{c}{\omS{S}{c}}$.
\end{proposition}
\begin{proof}
It suffices to construct one box partition for $c$; the choice then picks one. By \ref{r:groups},
$\rinf{S}(u)_\tau\in\mathcal U$ for $u\in\clU{S}{c}$.

\emph{Step 1: reduction of frame matrices.} We claim: for every family $(t_\tau)$ of elements of $\mathcal U$ there
is $u\in\clU{S}{c}$ such that all entries of the frame matrices $\fm{t_\tau\rinf{S}(u)_\tau}$ have absolute value at
most $b^{60\varrho+91}$. The entries $(k,l)$ with $k\ge l$ are $0$ or $1$ \ref{r:space}. We treat the pairs
$(k,l)$, $k<l$, one after the other, the rows $k$ in decreasing order. Suppose that $u_0\in\clU{S}{c}$ satisfies
the bound at the pairs already treated, and let $(k,l)$ be the next pair; write
$F_\tau=\fm{t_\tau\rinf{S}(u_0)_\tau}$ and
\[r_\tau=\gS{S}\,N\,\bscale{S}{\tau}{\fr{l}}/\bscale{S}{\tau}{e_k}=\gS{S}\,N\,\kapS{S}{\tau}^{\,k-l}.\]
By (J7), with $d=\vpS{S}$, there is an integer $x$ of $E$ with
$|F_{\tau,kl}/r_\tau-x_\tau|\le(2f)^2|\vpS{S}_\tau|$ for all $\tau$. Let $\lambda$ be the row vector with the entry $-xN$
at $\fr{l}$ and $0$ elsewhere. The position $(k,\fr{l})$ is free because $l>k$, and $\fadm{\lambda}\,\clB{S}{c}$ is
integral: the row $e_l$ of $\clB{S}{c}$ is a unit row, and its row $z$ is that of $\latB{S}$, with the entries
$1/N$ at $k_0$ and $1$ at $z$ \ref{r:groups}. Lemma~\ref{l:eichlerrow}, with $s=0$, gives $u_1\in\clU{S}{c}$ with
$u_1(e_{k'},i)=\delta_{e_{k'},i}+\gS{S}[k'=k]\,\lambda_i$ for $i\ne e^*_k$. By the formula for $\rinf{S}$, and because the
last row of the frame matrix of an element of $\mathcal U$ is a unit row, $\fm{\rinf{S}(u_1)_\tau}=1-r_\tau x_\tau\,E_{kl}$
with $E_{kl}$ the elementary matrix. As $\fm{\cdot}$ is multiplicative on $\mathcal U$,
\[\fm{t_\tau\rinf{S}(u_0u_1)_\tau}=F_\tau\,(1-r_\tau x_\tau E_{kl}):\]
the column $l$ of $F_\tau$ changes by $-r_\tau x_\tau$ times its column $k$, which is $0$ below the row $k$ and $1$ in
the row $k$. So the rows below $k$ and the other entries of the row $k$ are unchanged, and the entry $(k,l)$ becomes
$r_\tau(F_{\tau,kl}/r_\tau-x_\tau)$, of absolute value at most
\[|r_\tau|\,(2f)^2\,|\vpS{S}_\tau|\ \le\ b^{22}\cdot b^{60n}\cdot b^{2}\cdot b\ \le\ b^{60\varrho+91},\]
by $|\kapS{S}{\tau}|\ge1$ and $l>k$, $\gS{S}\le b^{22}$, $N\le b^{60n}$, $2f\le b$ and $|\vpS{S}_\tau|\le b$
\ref{r:data} and \ref{r:adm}. After the last pair the claim is proved.

\emph{Step 2: the function.} The map $\theta\colon u\mapsto(\fm{\rinf{S}(u)_\tau})_\tau$ is a homomorphism on $\clU{S}{c}$
with kernel $\clUo{S}{c}$ \ref{r:space}. For $y\in\prod_\tau\Gone{\fs}$ and a coset $q=u\,\clUo{S}{c}$ let
$\mathsf I(y,q)=1$ if
\[\Bigl\|\sum_{a}\theta(u)_{\tau,al}\ y_\tau^{-1}\fr{a}\Bigr\|\ \le\ b^{60\varrho+92}\,\bsize{y_\tau}\qquad\text{for all $\tau$ and all $l\le\varrho$,}\]
and $\mathsf I(y,q)=0$ otherwise. By (J1) the vector in the norm is $(\rinf{S}(u)_\tau^{-1}y_\tau)^{-1}\fr{l}$. Put
$\Sigma(y)=\sum_q\mathsf I(y,q)$ and $\omega(y)=\mathsf I(y,1)/\Sigma(y)$. For $u\in\clU{S}{c}$, (J1) gives
$\bsize{\rinf{S}(u)_\tau y_\tau}=\bsize{y_\tau}$, hence $\mathsf I(\rinf{S}(u)y,q)=\mathsf I(y,u^{-1}q)$ and
$\Sigma(\rinf{S}(u)y)=\Sigma(y)$.

\emph{$\Sigma(y)\ge1$.} By (J1) choose $w_\tau\in\mathcal U$ with $\|w_\tau y_\tau\|\le\bsize{y_\tau}$, and let $u$ be as
in Step~1 for $t=w$. By (J1) for $w_\tau\rinf{S}(u)_\tau\in\mathcal U$,
$(\rinf{S}(u)_\tau^{-1}y_\tau)^{-1}\fr{l}=\sum_a\fm{w_\tau\rinf{S}(u)_\tau}_{al}\,(w_\tau y_\tau)^{-1}\fr{a}$, of norm at most
$(\varrho+1)\,b^{60\varrho+91}\,\|w_\tau y_\tau\|\le b^{60\varrho+92}\,\bsize{y_\tau}$. So $\mathsf I(y,u\,\clUo{S}{c})=1$.

\emph{$\Sigma(y)<\infty$.} If $\mathsf I(y,u\,\clUo{S}{c})=1$ then, by the last inequality of (J1) with $s=1$ and the
index $\fr{a}$, $|\theta(u)_{\tau,al}|\le\|y_\tau\|\,b^{60\varrho+92}\,\bsize{y_\tau}$ for all $a,l$. So $\theta(u)$ lies in
a compact set that depends only on $y$, and by (S6) only finitely many values $\theta(u)$ occur there; $\theta$
is injective on cosets.

\emph{The five properties.} (i) Every $\mathsf I(\cdot,q)$ is the indicator of a set defined by inequalities
between continuous functions, and there are countably many $q$. (ii) If $\omega(y)\ne0$ then $\mathsf I(y,1)=1$, so
$\|y_\tau^{-1}\fr{l}\|\le b^{60\varrho+92}\bsize{y_\tau}$; if also $\InCut{y_\tau}{b}$ this is at most
$b^{60\varrho+92+6n+2}\le b^{2n^2}=\bxs{S}$, as $\varrho\ge10^5$. (iii) Let $\alpha$ be as in \ref{r:cutoff}. Then
\[\sum_{\gamma\in\clU{S}{c}}\alpha(\gamma)\,\omega(\rinf{S}(\gamma)y)=\frac{1}{\Sigma(y)}\sum_{q}\mathsf I(y,q^{-1})\sum_{\gamma\in q}\alpha(\gamma)
=\frac{1}{\Sigma(y)}\sum_q\mathsf I(y,q^{-1})=1,\]
because the sum of $\alpha$ over a coset of $\clUo{S}{c}$ is $1$ and $0<\Sigma(y)<\infty$. (iv)
$\mathsf I(y,1)\le\Sigma(y)$. (v) If every $s_\tau$ lies in the group generated by $\UoT{\Trg{\fs}{1}}$ and
$\HT{\Trg{\fs}{1}}$, which is contained in $\ST{\Trg{\fs}{1}}$, then by (J1) the frame vectors and the numbers
$\bsize{\cdot}$ of $sy$ and $y$ agree, so $\mathsf I(sy,q)=\mathsf I(y,q)$ for all $q$.
\end{proof}
\leanline{\lean{isBoxPartitionepsilon}}

\begin{lemma}[the cutoffs]\label{l:sumone}
In Setting~\ref{set:adm}, let $c\in\Cl{S}$.
\begin{enumerate}[label=(\alph*),itemsep=1pt]
\item $\sum_{u\in\clU{S}{c}}\cutoff{S}{c}\bigl(\rinf{S}(u)\,y\bigr)=1$ for every $y\in\prod_\tau\Gone{\fs}$.
\item $\cutoff{S}{c}$ is measurable.
\item If $y\in\prod_\tau\Gone{\fs}$ satisfies $\InCut{y_\tau}{\bS{S}}$ for all $\tau$, then there is $u\in\clU{S}{c}$
with $\rinf{S}(u)\,y\in\bxp{S}$.
\end{enumerate}
\end{lemma}
\begin{proof}
(a) Apply property (iii) of the box partition $\omS{S}{c}$ (Proposition~\ref{p:boxexists}) to
$\alpha(\gamma)=\bfl{S}{c}(\rinf{S}(\gamma)y)$, whose sums over the cosets $\clUo{S}{c}\gamma$ are $1$ by the choice of
$\bfl{S}{c}$. (b) Both factors are measurable. (c) By (a) some $u\in\clU{S}{c}$ has
$\cutoff{S}{c}(\rinf{S}(u)y)\ne0$, hence $\omS{S}{c}(\rinf{S}(u)y)\ne0$. Since $\rinf{S}(u)_\tau\in\mathcal U$, (J1)
gives $\InCut{\rinf{S}(u)_\tau y_\tau}{\bS{S}}$, and property (ii) puts $\rinf{S}(u)_\tau y_\tau$ in $\frbox{\bxs{S}}$.
\end{proof}
\leanline{\lean{tsumclassCutoffeqone}\sep\lean{measurableclassCutoff}\sep\lean{existsclassUnipmulmemboxedPoints}}

\begin{proof}[Proof of Theorem B]
Let $\sA=\{p:\sw_p\ne0\}$, and give the space of families $(v_{\tau,a})$ the norm $\sup_{\tau,a}\|v_{\tau,a}\|$. Since
$\rinf{S}(\gamma_p)^{-1}y\in\bxp{S}$, we have $\|F_p\|\le\bxs{S}$ for $p\in\sA$.

\emph{Step 1: close frames are equal.} Let $p,q\in\sA$ with $\varepsilon=\|F_p-F_q\|<(N\,\bxs{S})^{-1}$. We show
$\gamma_p\mathfrak f=\gamma_q\mathfrak f$. Put $g=\gamma_p^{-1}\gamma_q$ and $y^{\circ}=\rinf{S}(\gamma_p)^{-1}y\in\bxp{S}$. By
(S3), $g\in\terms{S}{d_p}{1}{d_q}$, and $1\in\terms{S}{d_p}{1}{d_p}$. For $u\in\GS{S}{E}$ let $F'(u)$ be the family
of the vectors $(\rinf{S}(u)_\tau^{-1}y^{\circ}_\tau)^{-1}\fr{a}$; then $F'(1)=F_p$ and $F'(g)=F_q$.

(a) $\Zf{S}{g}=\Zf{S}{1}$. By (J1), $\Zr{\rinf{S}(u)_\tau}_{la}$ is the hermitian product of
$(y^{\circ}_\tau)^{-1}\fr{l}$ with the image under $\JT{\Trg{\fs}{1}}$ of the component $(\tau,a)$ of $F'(u)$. So, by (J1) again,
$|\Zr{\rinf{S}(g)_\tau}_{la}-\Zr{1}_{la}|\le\bxs{S}\,\varepsilon$, and by the last formula of (S5) and $|\kapS{S}{\tau}|\ge1$,
$|(\Zf{S}{g}_{la}-\Zf{S}{1}_{la})_\tau|\le N\,\bxs{S}\,\varepsilon<1$ for all $\tau$. The entries of $\Zf{S}{g}$ and
$\Zf{S}{1}$ are algebraic integers (S3); by (J4) they agree.

(b) \emph{The rows of $g\mathfrak f$ and $\mathfrak f$ outside the rows $e_a$ agree.} $\Zf{S}{1}$ is the diagonal matrix
with the entries $0,\dots,0,-N$. The entry $(a,l)$ of $\Zf{S}{g}$ is $g(e^*_a,\fr{l})$ for $a<\varrho$ and
$-N\,g(z,\fr{l})$ for $a=\varrho$. So the rows $e^*_a$ and $z$ of $g\mathfrak f$ are those of $\mathfrak f$. Let $v_l$ be the
middle part of $(g-1)\fr{l}$; its $z$-coordinate is $0$. The vectors $g\fr{l}$ and $\fr{l}$ have no
$e^*$-coordinates, so $h(g\fr{l},g\fr{l})=h(\fr{l},\fr{l})$ reads
$(v_l+\Mid\fr{l})^{*}T(v_l+\Mid\fr{l})=(\Mid\fr{l})^{*}T\,\Mid\fr{l}$ with $T=\Trat{S}{E}$; as $T$ is diagonal and
$\Mid\fr{l}$ is $0$ or the unit vector at $z$, this gives $v_l^{*}Tv_l=0$, and $v_l=0$ by (S5).

(c) \emph{The rows $e_a$.} We show $(g-1)(e_a,\fr{l})=0$ for all $l$, by descending induction on $a<\varrho$. By
(b) and Lemma~\ref{l:framediff} (for $g$ and $1$), the component $(\tau,l)$ of $F'(g)-F'(1)$ is
$\sum_{a'<\varrho}c_{a'}\,(y^{\circ}_\tau)^{-1}\fr{a'}$ with
$c_{a'}=((g-1)(e_{a'},\fr{l}))_\tau\,\bscale{S}{\tau}{\fr{l}}/\bscale{S}{\tau}{e_{a'}}$, and $c_{a'}=0$ for $a'>a$ by the
induction hypothesis. By (J1) there is $w\in\mathcal U$ with $\|wy^{\circ}_\tau\|\le\bsize{y^{\circ}_\tau}\le b^{6n+2}$, and
\[\Bigl|\sum_{a'<\varrho}c_{a'}\,w(e_a,e_{a'})\Bigr|\ \le\ \|wy^{\circ}_\tau\|\cdot\Bigl\|\sum_{a'<\varrho}c_{a'}\,(y^{\circ}_\tau)^{-1}\fr{a'}\Bigr\|\ \le\ b^{6n+2}\,\varepsilon.\]
As $w(e_a,e_{a'})$ is $0$ for $a'<a$ and $1$ for $a'=a$, the left side is $|c_a|$. Hence
\[\bigl|((g-1)(e_a,\fr{l}))_\tau\bigr|\ \le\ |\kapS{S}{\tau}|^{\,l-a}\,b^{6n+2}\,\varepsilon\ \le\ b^{23\varrho+6n+2}\,\varepsilon\ \le\ \bxs{S}\,\varepsilon<1,\]
using $1\le|\kapS{S}{\tau}|\le b^{23}$. The numbers $g(e_a,\fr{l})$ are algebraic integers: by \ref{r:groups} the
matrices $\clB{S}{\cdot}$ are the identity on the rows $e_a$ and fix the frame vectors, so $g(e_a,\fr{l})$ is an
entry of the position $\pos{S}{d_p}{d_q}{g}$, which is integral by (S3). By (J4), $(g-1)(e_a,\fr{l})=0$.

So $g\mathfrak f=\mathfrak f$, that is $\gamma_p\mathfrak f=\gamma_q\mathfrak f$ by (S5).

\emph{Step 2: the count.} Let $\sA_0\subseteq\sA$ contain exactly one $p$ for every matrix $\gamma_p\mathfrak f$, $p\in\sA$.
By Step~1 the points $F_p$, $p\in\sA_0$, are pairwise at distance at least $(N\bxs{S})^{-1}$. Fix $p_0\in\sA_0$; the
points $F_p-F_{p_0}$ lie in $V$, by (iii), and have norm at most $2\,\bxs{S}$. By (J3),
\[\#\sA_0\ \le\ \bigl(1+4N\,\bxs{S}^2\bigr)^{\dim_{\RR}V}\ \le\ b^{\,6n^2\dim_{\RR}V},\]
since $1+4N\bxs{S}^2\le b^{3+60n+4n^2}\le b^{6n^2}$.

Colour $p\in\sA$ by the matrix $\gamma_p\mathfrak f$; there are $\#\sA_0$ colours. Call $p,q$ related if $d_p=d_q$ and
$\gamma_p^{-1}\gamma_q\in\clUo{S}{d_p}$. Two elements of one colour are related, by (S3) and (ii). Let
$\sB\subseteq\sA$ be a set of pairwise related elements and $p\in\sB$; write $d=d_p$. For $q\in\sB$ put
$u_q=\delta_q^{-1}\gamma_p=(\delta_q^{-1}\gamma_q)(\gamma_q^{-1}\gamma_p)\in\clU{S}{d}$; by (i) the map $q\mapsto u_q$ is
injective on $\sB$. By (ii) and Lemma~\ref{l:sumone}(a),
\[\sum_{q\in\sB}\sw_q\ \le\ x\sum_{q\in\sB}\cutoff{S}{d}\bigl(\rinf{S}(u_q)\,\rinf{S}(\gamma_p)^{-1}y'\bigr)\ \le\ x.\]
Lemma~\ref{l:colour} gives $\sum_{p\in\sA}\sw_p\le\#\sA_0\cdot x$, which is the assertion.
\end{proof}
\leanline{\lean{tsumclassCutoffboxle}}

\sect{The covolume of the arithmetic unipotent group}{s:covol}
As before $\mathcal U=\UT{\Trg{\fs}{1}}$. For $a<\varrho$ let $\mathrm{Fr}(a)$ be the set of the $l\in\Ix{\fs}$ such
that $(a,l)$ is a free position (Section~\ref{s:eichler}): the $e_b$ and $e^*_b$ with $b>a$ and the middle indices;
it has $2(\varrho-1-a)+\kk{\fs}+1$ elements.

\statement{ana}{\textbf{Inputs.} \textbf{(J8)}, unit~118 (coordinates of the unipotent group). The
\emph{coordinates} of $u\in\mathcal U$ are the complex numbers $u(e_a,l)$, $a<\varrho$, $l\in\mathrm{Fr}(a)$, and the
real numbers $\ImPt u(e_a,e^*_a)$, $a<\varrho$. The coordinate map is injective on $\mathcal U$ and carries the measure
$\muU{\fs}$ to Lebesgue measure.

\textbf{(J9)}, unit~106 (fundamental domains). (i) Let $\Omega$ be a finite set, $(t_{l,\tau})$ non-zero complex
numbers indexed by $l\in\Omega$ and the infinite places $\tau$ of $E$, and $M$ an invertible matrix over $\AEf$ on
$\Omega$. The group of the families $(t_{l,\tau}\,x_{l,\tau})_{l,\tau}$, for $x\in E^{\Omega}$ with $M\fadm{x}$ integral,
has a measurable set $F\subseteq\CC^{\Omega\times\{\tau\}}$ of volume
$\prod_\tau\bigl|\prod_lt_{l,\tau}\bigr|^2\big/\bigl(\gix{\det M}\,\volO{E}^{\#\Omega}\bigr)$ that contains exactly one
translate of every point by the group. (ii) Let $E$ be stably unramified and $(t_\tau)$ non-zero real numbers.
The group of the families $(t_\tau\,\ImPt s_\tau)_\tau$, for integers $s$ of $E$ with $\cj{s}=-s$, has a measurable set
$F\subseteq\RR^{\{\tau\}}$ of volume $\prod_\tau|t_\tau|\cdot|d_E|^{1/4}$ with the same property. (iii) Let $G$ be a
group, $\Gamma$ a subgroup, $(V_i)_{i\in\iota}$ a finite family of abelian groups with a degree
$\deg\colon\iota\to\mathbb N$, $\mathsf c=(\mathsf c_i)\colon G\to\prod_iV_i$ an injective map, and for every $i$ a
submonoid $L_i\subseteq V_i$ and a set $F_i\subseteq V_i$ that contains exactly one translate of every point
by $L_i$. Suppose that (1) for every $i$ and every $l\in L_i$ there is $\gamma\in\Gamma$ with $\mathsf c_i(\gamma)=l$ and
$\mathsf c_{i'}(\gamma)=\mathsf c_{i'}(1)$ for all $i'\ne i$ with $\deg i'\le\deg i$; and (2) for every $i$ and every
$\gamma\in\Gamma$ with $\mathsf c_{i'}(\gamma)=\mathsf c_{i'}(1)$ for all $i'$ with $\deg i'<\deg i$, one has
$\mathsf c_i(\gamma)\in L_i$ and $\mathsf c_i(x\gamma)=\mathsf c_i(\gamma)+\mathsf c_i(x)$ for all $x\in G$. Then for every
$x\in G$ there is exactly one $\gamma\in\Gamma$ with $\mathsf c(x\gamma)\in\prod_iF_i$.

The catalogue classes these inputs as elementary and gives no reference.}

\begin{lemma}[additivity of the coordinates]\label{l:additive}
Let $u\in\mathcal U$ and $a<\varrho$. Suppose that the coordinates of $u$ in the rows $e_b$, $b>a$, vanish: the
entries $u(e_b,l)$, $l\in\mathrm{Fr}(b)$, and the numbers $\ImPt u(e_b,e^*_b)$.
\begin{enumerate}[label=(\alph*),itemsep=1pt]
\item Then $(xu)(e_a,l)=x(e_a,l)+u(e_a,l)$ for all $x\in\mathcal U$ and all $l\in\mathrm{Fr}(a)$.
\item If moreover $u(e_a,l)=0$ for all $l\in\mathrm{Fr}(a)$, then $\RePt u(e_a,e^*_a)=0$ and
$(xu)(e_a,e^*_a)=x(e_a,e^*_a)+u(e_a,e^*_a)$ for all $x\in\mathcal U$.
\end{enumerate}
\end{lemma}
\begin{proof}
Since $(xu)(e_a,l)=x(e_a,l)+u(e_a,l)+\sum_i(x-1)(e_a,i)\,(u-1)(i,l)$ and $(x-1)(e_a,i)=0$ unless $\lev(i)>a$, it
suffices to show that $(u-1)(i,l)=0$ for all $i$ with $\lev(i)>a$, that is for $i=e_b$ with $b>a$, $i$ a middle
index, or $i=e^*_b$; and for $l\in\mathrm{Fr}(a)$ in (a), $l=e^*_a$ in (b).

Let $J=\JT{\Trg{\fs}{1}}$; then $J^2=1$, so $uJu^{*}=J$: writing $\rho_i$ for the row $i$ of $u$ and
$\langle\rho,\rho'\rangle=\rho J\rho'^{*}$, we have $\langle\rho_i,\rho_{i'}\rangle=J(i,i')$. Let $c>a$, or $c=a$ in
case (b). Outside the columns $e^*_d$ the row $\rho_{e_c}$ is the unit row at $e_c$, because $u\in\mathcal U$ and
its free entries vanish. Hence $\langle\rho,\rho_{e_c}\rangle=\rho(e^*_c)+\sum_d\rho(e_d)\,\cj{u(e_c,e^*_d)}$ for every
row vector $\rho$.

For $i$ a middle index or $i=e^*_b$ we have $\rho_i(e_d)=0$, so $u(i,e^*_c)=J(i,e_c)=\delta_{i,e^*_c}$. Now let $i=e_b$ with $b>a$, or $b=a$ in case (b). Then $\rho_{e_b}(e_d)=\delta_{bd}$, so
$u(e_b,e^*_c)+\cj{u(e_c,e^*_b)}=0$. If $b<c$ the first term is a free entry of the row $e_b$, hence $0$, and so is
the second; by symmetry both terms vanish whenever $b\ne c$. If $b=c$ the relation says $\RePt u(e_b,e^*_b)=0$;
for $b>a$ also $\ImPt u(e_b,e^*_b)=0$ by hypothesis, and for $b=a$ this is the first assertion of (b). This proves
$(u-1)(i,e^*_c)=0$ for all $i$ of level greater than $a$ and all such $c$.

It remains to treat, in (a), the columns $l=e_c$ with $c>a$ and the middle columns $l$. There $u(i,l)=\delta_{il}$
if $\lev(i)\ge\lev(l)$, and otherwise $i=e_b$ with $a<b$, and $(b,l)$ is a free position.
\end{proof}

\begin{proof}[Proof of Theorem C]
Write $g=\gS{S}$, $N=\NS{S}$, $k=\kk{\fs}$, $f=\ninf{E}$. Let $G=\prod_\tau\mathcal U$ with the measure
$\mu=\prod_\tau\muU{\fs}$, let $\varphi\colon\clU{S}{c}\to G$, $u\mapsto(\rinf{S}(u)_\tau)_\tau$, an injective
homomorphism \ref{r:groups}, and $\Gamma$ its image.

\emph{The class basis.} By \ref{r:groups}, $\clB{S}{c}$ is block diagonal: the identity on the $e_a$ and $e^*_a$,
and a matrix $m_c$ on $\Mx{\fs}$, the product of $\operatorname{diag}(\tilde c,1)$ and the middle block of $\latB{S}$,
whose determinant is $\cj{\piS{S}}^{-1}$. As $\tilde c$ is unitary, $\gix{\det\tilde c}^2=1$; and
$\gix{\piS{S}}^2=\gix{N}=N^{2f}$ by (J2). Hence $\gix{\det m_c}=N^{-f}$. For $a<\varrho$ let $m(a)$ be the transpose
of the block of $\clB{S}{c}$ on $\mathrm{Fr}(a)$: for a row vector $x$ supported on $\mathrm{Fr}(a)$,
$m(a)x$ is the restriction of $x\,\clB{S}{c}$ to $\mathrm{Fr}(a)$, and $x\,\clB{S}{c}$ vanishes outside
$\mathrm{Fr}(a)$. Its determinant is $\det m_c$.

\emph{Graded coordinates.} The index set $\iota$ consists of the pairs $(a,\mathrm{fr})$ and $(a,\mathrm{sk})$,
$a<\varrho$, with $\deg(a,\mathrm{fr})=2(\varrho-1-a)$ and $\deg(a,\mathrm{sk})=2(\varrho-1-a)+1$. For
$u=(u_\tau)\in G$ put $\mathsf c_{a,\mathrm{fr}}(u)=(u_\tau(e_a,l))_{l\in\mathrm{Fr}(a),\tau}$ and
$\mathsf c_{a,\mathrm{sk}}(u)=(\ImPt u_\tau(e_a,e^*_a))_\tau$. By (J8), $\mathsf c$ is injective and carries $\mu$ to
Lebesgue measure. Let
\[t_{l,\tau}=g\,\bscale{S}{\tau}{l}/\bscale{S}{\tau}{e_a}\quad(l\in\mathrm{Fr}(a)),\qquad t'_\tau=g/\fscale{S}{\tau}{a},\]
let $L_{a,\mathrm{fr}}$ be the group of the families $(t_{l,\tau}x_{l,\tau})$, $x\in E^{\mathrm{Fr}(a)}$ with
$m(a)\fadm{x}$ integral, and $L_{a,\mathrm{sk}}$ the group of the families $(t'_\tau\ImPt s_\tau)$, $s$ an integer of $E$
with $\cj{s}=-s$. Let $F_{a,\mathrm{fr}}$, $F_{a,\mathrm{sk}}$ be the sets given by (J9)(i),(ii).

\emph{Condition (1) of (J9)(iii).} Let $(t_{l,\tau}x_{l,\tau})\in L_{a,\mathrm{fr}}$, and extend $x$ by $0$ to a row
vector on $\Ix{\fs}$. Then $\fadm{x}\,\clB{S}{c}$ is integral, and Lemma~\ref{l:eichlerrow} with $s=0$ gives
$u\in\clU{S}{c}$ with $u(e_b,l)=\delta_{e_b,l}+g\,[b=a]\,x_l$ for all $b<\varrho$ and $l\ne e^*_a$. By the formula for
$\rinf{S}$, $\mathsf c_{a,\mathrm{fr}}(\varphi(u))=(t_{l,\tau}x_{l,\tau})$, and all other coordinates of $\varphi(u)$
vanish, except possibly $\mathsf c_{a,\mathrm{sk}}$, whose degree is larger. Let $(t'_\tau\ImPt s_\tau)\in L_{a,\mathrm{sk}}$.
Lemma~\ref{l:eichlerrow} with $l=0$ gives $u\in\clU{S}{c}$ with $u(e_b,j)=\delta_{e_b,j}+g\,[b=a]\,[j=e^*_a]\,s$ for
all $b$ and $j$. Since $\bscale{S}{\tau}{e^*_a}/\bscale{S}{\tau}{e_a}=1/\fscale{S}{\tau}{a}$ is real,
$\mathsf c_{a,\mathrm{sk}}(\varphi(u))=(t'_\tau\ImPt s_\tau)$, and all other coordinates vanish.

\emph{Condition (2).} Let $u\in\clU{S}{c}$, let $i=(a,\mathrm{fr})$ or $i=(a,\mathrm{sk})$, and suppose that the
coordinates of $\varphi(u)$ of degree less than $\deg i$ vanish. The additivity
$\mathsf c_i(x\varphi(u))=\mathsf c_i(\varphi(u))+\mathsf c_i(x)$ is Lemma~\ref{l:additive}, applied at every $\tau$. For
the membership in $L_i$ note first that the row $e_a$ of $\gid{S}^{-1}\fadm{u-1}\,\clB{S}{c}$ is integral: this matrix
equals $\clB{S}{c}\,\gid{S}^{-1}(\pos{S}{c}{c}{u}-1)$, the row $e_a$ of $\clB{S}{c}$ is a unit row, and
$\gid{S}^{-1}(\pos{S}{c}{c}{u}-1)$ is integral by (S1). If $i=(a,\mathrm{fr})$, put $x_l=g^{-1}u(e_a,l)$ for
$l\in\mathrm{Fr}(a)$; then $\mathsf c_i(\varphi(u))=(t_{l,\tau}x_{l,\tau})$, and $m(a)\fadm{x}$ is the restriction to
$\mathrm{Fr}(a)$ of that integral row, because the entries of the row $e_a$ of $u-1$ outside $\mathrm{Fr}(a)$ sit
in columns $e^*_d$, $d\le a$, on which $\clB{S}{c}$ is the identity. If $i=(a,\mathrm{sk})$, put
$s=g^{-1}u(e_a,e^*_a)$. It is integral at all finite places, hence an integer of $E$; and
$\RePt\rinf{S}(u)_\tau(e_a,e^*_a)=0$ by Lemma~\ref{l:additive}(b), that is $s_\tau+\cj{s_\tau}=0$ for all $\tau$, so
$\cj{s}=-s$. Then $\mathsf c_i(\varphi(u))=(t'_\tau\ImPt s_\tau)\in L_{a,\mathrm{sk}}$.

\emph{Unfolding.} Let $\mathfrak B=\prod_iF_i$ and $D=\mathsf c^{-1}(\mathfrak B)\subseteq G$. By (J9)(iii), for every
$x\in G$ there is exactly one $\gamma\in\Gamma$ with $x\gamma\in D$; applied to $x^{-1}$: exactly one $\gamma\in\Gamma$
with $\gamma x\in D^{-1}$. So $G$ is the disjoint union of the sets $\gamma^{-1}D^{-1}$, and by the left invariance of
$\mu$ and Lemma~\ref{l:sumone}(a),
\[\int_G\cutoff{S}{c}(uy)\,d\mu(u)=\int_{D^{-1}}\sum_{\gamma\in\clU{S}{c}}\cutoff{S}{c}\bigl(\rinf{S}(\gamma)\,uy\bigr)\,d\mu(u)=\mu(D^{-1})=\mu(D),\]
the last by the invariance of $\mu$ under inversion \ref{r:haar}. By (J8), $\mu(D)$ is the Lebesgue measure of
$\mathfrak B$, the product of the volumes of the $F_i$.

\emph{The volume.} The scales of the $e_b$ and $e^*_b$ are inverse to each other, so
\[\prod_{l\in\mathrm{Fr}(a)}|\bscale{S}{\tau}{l}|^2=|\bscale{S}{\tau}{z}|^2=N^{-1}\qquad\text{and}\qquad
|\bscale{S}{\tau}{e_a}|^{-2}=\fscale{S}{\tau}{a}^{-1}.\] With $n_a=\#\mathrm{Fr}(a)$, (J9) and
$\gix{\det m(a)}=N^{-f}$ give
\[\vol(F_{a,\mathrm{fr}})=\frac{g^{2n_af}}{\volO{E}^{\,n_a}}\prod_\tau\fscale{S}{\tau}{a}^{-n_a},\qquad
\vol(F_{a,\mathrm{sk}})=g^{f}\,|d_E|^{1/4}\prod_\tau\fscale{S}{\tau}{a}^{-1}.\]
Now $n_a+1=k+2(\varrho-a)$ and $\sum_{a<\varrho}n_a=\varrho(\varrho+k)$, so the product over $a<\varrho$ of these volumes is
\[g^{\varrho(2\varrho+2k+1)f}\ \volO{E}^{-\varrho(\varrho+k)}\ |d_E|^{\varrho/4}\ \prod_\tau\prod_{a<\varrho}\fscale{S}{\tau}{a}^{-(k+2(\varrho-a))}.\]
This is $\ucov{S}$ \ref{r:const}: the factor $a=\varrho$ of $\fsprod{S}$ is $\prod_\tau\fscale{S}{\tau}{\varrho}^{k}=N^{-kf}$,
which cancels the factor $N^{-kf}$ of $\ucov{S}$.
\end{proof}
\leanline{\lean{lintegralclassCutoffequnipCovolume}}

\begin{lemma}[unfolding in products]\label{l:unfold}
Let $A$ be a subgroup of $\Gone{\fs}$ and $\mu_A$ a measure on $A$ that is a unimodular Haar measure with a
cutoff: a $\sigma$-finite Haar measure of $A$, invariant under right translations and inversion, for which
there is a measurable $\beta\ge0$ on $\Gone{\fs}$ with $\int_A\beta(ay)\,d\mu_A(a)=1$ for all $y$; let
$\mu_{A\backslash}=\beta\,dy$ for the chosen such $\beta$. Then for every finite set $\Omega$ the map
$\bigl((x_p)_p,(a_p)_p\bigr)\mapsto(a_px_p)_p$ carries the measure
$\bigl(\prod_{p\in\Omega}\mu_{A\backslash}\bigr)\times\bigl(\prod_{p\in\Omega}\mu_A\bigr)$ to the Haar measure
$\prod_{p\in\Omega}dy$ of $\Gone{\fs}^{\Omega}$.
\end{lemma}
\begin{proof}
For one factor this is the unfolding of the background library of unit~114, recalled in \ref{r:haar}. A product of
measure preserving maps preserves the product measures, and regrouping the factors does.
\end{proof}
\leanline{\lean{measurePreservingpiorbitMeasure}}

The formal text records here that $\mu_{A\backslash}$ is $\sigma$-finite, the chosen $\beta$ being finite almost
everywhere, and later that $\GRes{S}$ is finite; both are unnamed instances.

\begin{lemma}[integrals against a cutoff]\label{l:cutprod}
In Setting~\ref{set:adm}, let $c\in\Cl{S}$ and, for every infinite place $\tau$ of $E$, let
$F_\tau\colon\Gone{\fs}\to[0,\infty]$ be measurable and invariant under left translation by $\mathcal U$. Then
\[\int\cutoff{S}{c}(y)\prod_\tau F_\tau(y_\tau)\,dy=\ucov{S}\cdot\prod_\tau\int F_\tau\,d\muq{\fs}.\]
\end{lemma}
\begin{proof}
By Lemma~\ref{l:unfold} for $A=\mathcal U$ and the set of the infinite places, the left side is
$\int\!\!\int\cutoff{S}{c}\bigl((u_\tau x_\tau)_\tau\bigr)\prod_\tau F_\tau(u_\tau x_\tau)\,d\mu(u)\,d\bigl(\prod_\tau\muq{\fs}\bigr)(x)$. The
factors $F_\tau(u_\tau x_\tau)=F_\tau(x_\tau)$ do not depend on $u$, the integral of the cutoff over $u$ is $\ucov{S}$ by
Theorem~C, and the remaining integral over $x$ is the product of the integrals of the $F_\tau$.
\end{proof}
\leanline{\lean{lintegralclassCutoffmulprod}}

\sect{Amplitude sums}{s:ampl}
In this section $\mathcal U=\UT{\Trg{\fs}{1}}$, and for $\gamma\in\GS{S}{E}$ and $y\in\prod_\tau\Gone{\fs}$ we write
$\gamma^{-1}y$ for $\rinf{S}(\gamma)^{-1}y$ and $\gamma y$ for $\rinf{S}(\gamma)y$. All integrals over $y$ are over
$\prod_\tau\Gone{\fs}$ for the Haar measure $dy$, which is invariant under these translations.

\begin{lemma}\label{l:stableS}
For $S\in\Fill{E}{\fs}$, the group $\ST{\Trat{S}{E}}$ of the rational points of the Bessel subgroup is stable:
$\IsStab{S}{\ST{\Trat{S}{E}}}$.
\end{lemma}
\begin{proof}
It is a group that contains $\UT{\Trat{S}{E}}$ \ref{r:space}.
\end{proof}
\leanline{\lean{isStablebesselS}}

\begin{lemma}[coordinates of a sum of integrals]\label{l:sumint}
Let $(\mathsf Y,\mu)$ be a measure space, $\sP$ a countable set, $I$ a finite set, and for $p\in\sP$ let
$g_p\colon\mathsf Y\to\CC^{I}$ be measurable, with $\sum_{p\in\sP}\int\|g_p\|\,d\mu<\infty$. Then for every $i\in I$,
\[\Bigl(\sum_{p\in\sP}\int g_p\,d\mu\Bigr)_i=\int\sum_{p\in\sP}g_p(x)_i\,d\mu(x).\]
\end{lemma}
\begin{proof}
Each $g_p$ is integrable and its integral has the coordinates $\int(g_p)_i\,d\mu$; the sum of the integrals
converges absolutely, so its coordinate $i$ is $\sum_p\int(g_p)_i\,d\mu$; and since
$\sum_p\int|(g_p)_i|\,d\mu<\infty$, sum and integral may be exchanged.
\end{proof}
\leanline{\lean{tsumintegralapply}}

\begin{lemma}[invariance of the Bessel functions and of the kernel]\label{l:Winv}
Let $S\in\Fill{E}{\fs}$ and let $\sigma$ be a side.
\begin{enumerate}[label=(\alph*),itemsep=1pt]
\item For $c\in\Cl{S}$, $u\in\clU{S}{c}$, $y\in\prod_\tau\Gone{\fs}$ and $i\in\Bc{S}$:
$\besfn{\sigma}{S}{\rinf{S}(u)\,y}{i}=\besfn{\sigma}{S}{y}{i}$.
\item For $u,u'\in\UT{\Trat{S}{E}}$, $\gamma\in\GS{S}{E}$ and $y\in\prod_\tau\Gone{\fs}$:
\[\besker{\sigma}{S}{u\gamma u'^{-1}}{\rinf{S}(u)\,y}=\uchar{S}{\rfin{S}(u)}^{-1}\,\uchar{S}{\rfin{S}(u')}\cdot\besker{\sigma}{S}{\gamma}{y}.\]
\end{enumerate}
\end{lemma}
\begin{proof}
(a) By (S4), $\besfn{\sigma}{S}{\rinf{S}(u)y}{i}=\uchar{S}{\rfin{S}(u)}^{-1}\besfn{\sigma}{S}{y}{i}$. The position of
$u$, which is $\pos{S}{c}{c}{u}=\hc{S}{c}^{-1}\rfin{S}(u)\hc{S}{c}$, lies in $\KgS{S}$, so the second statement of (S4), for the pair
$(1,\rfin{S}(u))$, gives $\uchar{S}{\rfin{S}(u)}=\uchar{S}{1}=1$. (b) The coordinate $(i,j)$ of the left side is
$\besfn{\sigma}{S}{\rinf{S}(u)y}{i}$ times the conjugate of
$\besfn{\sigma}{S}{\rinf{S}(u')\rinf{S}(\gamma)^{-1}y}{j}$. By (S4) these are
$\uchar{S}{\rfin{S}(u)}^{-1}\besfn{\sigma}{S}{y}{i}$ and the conjugate of
$\uchar{S}{\rfin{S}(u')}^{-1}\besfn{\sigma}{S}{\rinf{S}(\gamma)^{-1}y}{j}$, and the values of $\uchar{S}{\cdot}$ have absolute
value $1$.
\end{proof}
\leanline{\lean{besselFnmulofmemclassUnip}\sep\lean{besselKernelbesselUmul}}

\begin{proof}[Proof of Theorem D]
Write $b=\bS{S}$, $n=\nn{\fs}$, $r=\rr{\fs}$, $k=\kk{\fs}$, $f=\ninf{E}$. Take $A_\tau$, $m_\tau$ and $\theta_\tau$ from
(S7) and put, for $y\in\prod_\tau\Gone{\fs}$ and a side $t$,
\[\mathbf A(y)=\prod_\tau A_\tau(y_\tau),\qquad\mathbf W^{t}(y)=\prod_\tau\|\Wsys{t}{S}{\tau}(y_\tau)\|,\qquad
\Theta(y)=\prod_\tau\theta_\tau(y_\tau)\ \ge1,\]
so that $\Theta(y)\,\mathbf W^{t}(y)\le\mathbf A(y)$, and $\|\besker{t}{S}{\gamma}{y}\|=\mathbf W^{t}(y)\,\mathbf W^{t}(\gamma^{-1}y)$.

\emph{Step 1: the frame count.} We claim that for every $y$, with $C_0=b^{12n^3rf}$,
\[\sum_{d\in\Cl{S}}\ \sum_{\gamma\in X(d)}\cutoff{S}{d}(\gamma^{-1}y)\,\mathbf A(\gamma^{-1}y)\ \le\ C_0\prod_\tau m_\tau(\Gone{\fs}).\tag{1}\]
Let $m=\prod_\tau m_\tau$, a finite measure on $\prod_\tau\Gone{\fs}$, and for a point $v$ let $Q(v)$ be the set of the
$g$ with $\InCut{v_\tau g_\tau}{b}$ for all $\tau$. By (S7), $\mathbf A(v)\le m(Q(v))$. So the left side of (1) is at
most the integral over $g$, for $m$, of
\[\sum_{d}\ \sum_{\gamma\in X(d)}\cutoff{S}{d}(\gamma^{-1}y)\,1_{Q(\gamma^{-1}y)}(g),\]
and it suffices to bound this sum by $C_0$ for every $g$. Fix $g$ and apply Theorem~B to the following data: the
index set $\sP$ of the pairs $p=(d,\gamma)$ with $\gamma\in X(d)$; $d_p=d$, $\delta_p=\gamma$; the weight
$\sw_p=\cutoff{S}{d}(\gamma^{-1}y)\,1_{Q(\gamma^{-1}y)}(g)$; $x=1$; the point $y$ in the place of $y'$ and the point
$yg=(y_\tau g_\tau)_\tau$ in the place of $y$; and $V$ the whole space of families, of real dimension $2rnf$. For $p$
with $\sw_p\ne0$ we have $g\in Q(\gamma^{-1}y)$, so by Lemma~\ref{l:sumone}(c) there is $u_p\in\clU{S}{d}$ with
$\rinf{S}(u_p)(\gamma^{-1}y)g\in\bxp{S}$; put $\gamma_p=\gamma u_p^{-1}$. Then hypothesis (i) of Theorem~B is clear;
for (ii), $\gamma_p\in\terms{S}{c}{s}{d}$ by (S3), $\delta_p^{-1}\gamma_p=u_p^{-1}\in\clU{S}{d}$,
$\gamma_p^{-1}(yg)=\rinf{S}(u_p)(\gamma^{-1}y)g\in\bxp{S}$, and $\sw_p\le\cutoff{S}{d}(\delta_p^{-1}y)$; and (iii) is
empty. Theorem~B bounds the sum by $b^{6n^2\cdot2rnf}=C_0$.

\emph{Step 2: integration.} Multiply (1) by $\cutoff{S}{c}(y)\,\mathbf A(y)$ and integrate over $y$. By
Lemma~\ref{l:cutprod} and (S7),
$\int\cutoff{S}{c}(y)\mathbf A(y)\,dy=\ucov{S}\prod_\tau\int A_\tau\,d\muq{\fs}$ and
$\prod_\tau m_\tau(\Gone{\fs})\int A_\tau\,d\muq{\fs}\le b^{2r^3f}\,\#\Bc{S}$. Hence
\[\sum_{d\in\Cl{S}}\ \sum_{\gamma\in X(d)}\int\cutoff{S}{c}(y)\,\cutoff{S}{d}(\gamma^{-1}y)\,\mathbf A(y)\,\mathbf A(\gamma^{-1}y)\,dy
\ \le\ C_0\,b^{2r^3f}\,\#\Bc{S}\cdot\ucov{S}.\tag{2}\]

\emph{Assertion (1) of the theorem.} Since $\mathbf W^{\sigma}\le\mathbf A$, the first sum is at most the left side
of (1), so $M=C_0\prod_\tau m_\tau(\Gone{\fs})$ will do. By (S7) and the definition,
$\ampabs{\sigma}{S}{c}{c'}{\gamma}=\int\cutoff{S}{c}(y)\cutoff{S}{c'}(\gamma^{-1}y)\mathbf W^{\sigma}(y)\mathbf W^{\sigma}(\gamma^{-1}y)\,dy$, and
the sum of these over $\gamma\in X(c')$ is at most the left side of (2), which is finite.

\emph{Assertion (2).} We claim that for all $\gamma$ and $y$
\[\bigl\|\besker{\sigma}{S}{\gamma}{y}-\besker{\cut}{S}{\gamma}{y}\bigr\|\ \le\ 2\,b^{-3r^3f}\,\mathbf A(y)\,\mathbf A(\gamma^{-1}y).\]
If $\Theta(y)=\Theta(\gamma^{-1}y)=1$, the full and the cut Bessel functions agree at all $y_\tau$ and all
$(\gamma^{-1}y)_\tau$, and the left side is $0$. Otherwise $\Theta(y)\,\Theta(\gamma^{-1}y)\ge b^{3r^3f}$, and for both
sides $t$, $\|\besker{t}{S}{\gamma}{y}\|=\mathbf W^{t}(y)\mathbf W^{t}(\gamma^{-1}y)\le b^{-3r^3f}\mathbf A(y)\mathbf A(\gamma^{-1}y)$;
the claim follows by the triangle inequality. With (2) and $\#\Bc{S}\le b^{2kf}$ \ref{r:adm}, the left side of
assertion (2) is at most
\[2\,b^{-3r^3f}\cdot C_0\,b^{2r^3f}\,b^{2kf}\cdot\ucov{S}=2\,b^{(12n^3r+2k-r^3)f}\,\ucov{S}.\]
Now $r=k+2n-1$ with $2n-1\le k$ \ref{r:data}, so $r\le2k$ and $r^3-k^3\ge3k^2$, while
$12n^3r+2k\le24n^3k+2k\le2k^2$ because $k\ge1000(2n)^9$. Hence the exponent is at most $-(k^3+k^2)f$, and
$2\,b^{-k^2f}\le\tfrac12$.
\end{proof}
\leanline{\lean{amplitudeAbslttopandcutle}}

\begin{corollary}\label{c:summable}
In Setting~\ref{set:adm}, for every side $\sigma$, every set $\Xi\subseteq\GS{S}{E}$, all $c,c'\in\Cl{S}$ and
$s\in\GRes{S}$, the family of the amplitudes $\amp{\sigma}{S}{c}{c'}{\gamma}$, $\gamma\in\Xi\cap\terms{S}{c}{s}{c'}$, is
absolutely summable.
\end{corollary}
\begin{proof}
$\|\amp{\sigma}{S}{c}{c'}{\gamma}\|\le\ampabs{\sigma}{S}{c}{c'}{\gamma}$ by (S7), and Theorem~D(1).
\end{proof}
\leanline{\lean{summableamplitude}}

\begin{proof}[Proof of Theorem E]
The countable group $\Delta=\clU{S}{c}\times\clU{S}{c'}$ acts on $\mathsf P=X\times\prod_\tau\Gone{\fs}$ by
$(u,u')\cdot(\gamma,y)=(u\gamma u'^{-1},\rinf{S}(u)y)$: indeed $u\gamma u'^{-1}\in\Xi$ because $\Xi$ is stable, and
$u\gamma u'^{-1}\in\terms{S}{c}{s}{c'}$ by (S3), as $u\in\terms{S}{c}{1}{c}$ and $u'^{-1}\in\terms{S}{c'}{1}{c'}$. The action
preserves the measure $\mu_{\mathsf P}$ on $\mathsf P$ that is the product of the counting measure on $X$ and $dy$. By
Lemma~\ref{l:Winv}(a) the functions $\mathsf K^{t}(\gamma,y)=\besker{t}{S}{\gamma}{y}$ are invariant under $\Delta$,
for both sides $t$. Put $\beta_0(\gamma,y)=\cutoff{S}{c}(y)\,\cutoff{S}{c'}(\gamma^{-1}y)$. By
Lemma~\ref{l:sumone}(a),
\[\sum_{(u,u')\in\Delta}\beta_0\bigl((u,u')\cdot(\gamma,y)\bigr)=\sum_{u}\cutoff{S}{c}(\rinf{S}(u)y)\cdot\sum_{u'}\cutoff{S}{c'}\bigl(\rinf{S}(u')\gamma^{-1}y\bigr)=1,\]
and $\phi$ has the same property by hypothesis; in particular $\phi\le1$ on $\mathsf P$. For a measurable
$\Delta$-invariant $F\colon\mathsf P\to[0,\infty]$ we get, inserting the sum of $\beta_0$ and then moving $\delta$ to
the other side by the invariance of $\mu_{\mathsf P}$ and of $F$,
\begin{align*}\int\phi F\,d\mu_{\mathsf P}&=\sum_{\delta\in\Delta}\int\phi(p)F(p)\beta_0(\delta p)\,d\mu_{\mathsf P}(p)\\
&=\sum_{\delta\in\Delta}\int\phi(\delta^{-1}p)F(p)\beta_0(p)\,d\mu_{\mathsf P}(p)=\int\beta_0F\,d\mu_{\mathsf P}.\tag{3}\end{align*}
For $F=\|\mathsf K^{\sigma}\|$ the right side is $\sum_{\gamma\in X}\ampabs{\sigma}{S}{c}{c'}{\gamma}$, finite by Theorem~D(1).
So every $y\mapsto\phi(\gamma,y)\besker{\sigma}{S}{\gamma}{y}$ is integrable, the sums below converge absolutely, and
the computation (3), now with the vector valued $\mathsf K^{\sigma}$ in the place of $F$ and Fubini's theorem,
gives
\[\sum_{\gamma\in X}\int\phi(\gamma,y)\besker{\sigma}{S}{\gamma}{y}\,dy=\sum_{\gamma\in X}\int\beta_0(\gamma,y)\besker{\sigma}{S}{\gamma}{y}\,dy
=\ampsum{\sigma}{S}{\Xi}{c}{s}{c'}.\]
Finally (3) for $F=\|\mathsf K^{\sigma}-\mathsf K^{\cut}\|$ bounds the last sum of the theorem by the term $d=c'$ of
Theorem~D(2).
\end{proof}
\leanline{\lean{hasSumamplitudeSumchangeofweight}}

\begin{proposition}[moves]\label{p:moves}
In Setting~\ref{set:adm}, let $\sigma$ be a side and $\Xi\subseteq\GS{S}{E}$ a set with $\IsStab{S}{\Xi}$.
\begin{enumerate}[label=(\alph*),itemsep=1pt]
\item Let $c,c_1,c',c_1'\in\Cl{S}$, $t,t'\in\GRes{S}$, and let $\eta,\eta'\in\ST{\Trat{S}{E}}$ be such that
$\eta\in\terms{S}{c}{t}{c_1}$ and $\eta'\in\terms{S}{c'}{t'}{c_1'}$, and that
$\gamma\in\Xi\Leftrightarrow\eta\gamma\eta'^{-1}\in\Xi$ for all $\gamma\in\GS{S}{E}$. Let $L$ be a continuous $\CC$-linear
map of $\CC^{\Bc{S}\times\Bc{S}}$ to itself such that
$\besker{\sigma}{S}{\eta\gamma\eta'^{-1}}{\rinf{S}(\eta)\,y}=L\bigl(\besker{\sigma}{S}{\gamma}{y}\bigr)$ for all $\gamma$ and $y$.
Then for every $s\in\GRes{S}$,
\[\ampsum{\sigma}{S}{\Xi}{c}{t\,s\,t'^{-1}}{c'}=L\bigl(\ampsum{\sigma}{S}{\Xi}{c_1}{s}{c_1'}\bigr).\]
\item Let $u,u'\in\KS{S}\cap\UT{\Tlat{S}{\AEf}}$, $c,c'\in\Cl{S}$ and $s\in\GRes{S}$. Then
\[\ampsum{\sigma}{S}{\Xi}{c}{\redS{S}(u)\,s\,\redS{S}(u')^{-1}}{c'}=\uchar{S}{u}^{-1}\,\uchar{S}{u'}\cdot\ampsum{\sigma}{S}{\Xi}{c}{s}{c'}.\]
\end{enumerate}
\end{proposition}
\begin{proof}
(a) By (S3), $\gamma\in\terms{S}{c_1}{s}{c_1'}$ if and only if $\eta\gamma\eta'^{-1}\in\terms{S}{c}{t\,s\,t'^{-1}}{c'}$. For
$\eta'=\eta$ and $s=1$ this shows, together with the fact that $\eta$ normalizes $\UT{\Trat{S}{E}}$, that
$u\mapsto\eta u\eta^{-1}$ is a bijection of $\clU{S}{c_1}$ onto $\clU{S}{c}$; likewise for $\eta'$. Put
$\phi(\gamma,y)=\cutoff{S}{c}(\eta y)\,\cutoff{S}{c'}(\eta'\gamma^{-1}y)$. For $\gamma\in X_1=\Xi\cap\terms{S}{c_1}{s}{c_1'}$,
\begin{multline*}\sum_{(u,u')\in\clU{S}{c_1}\times\clU{S}{c_1'}}\phi\bigl(u\gamma u'^{-1},uy\bigr)\\
=\sum_{u}\cutoff{S}{c}\bigl((\eta u\eta^{-1})\,\eta y\bigr)\cdot\sum_{u'}\cutoff{S}{c'}\bigl((\eta'u'\eta'^{-1})\,\eta'\gamma^{-1}y\bigr)=1\end{multline*}
by the two bijections and Lemma~\ref{l:sumone}(a). Theorem~E, for $(c_1,s,c_1')$, gives
$\ampsum{\sigma}{S}{\Xi}{c_1}{s}{c_1'}=\sum_{\gamma\in X_1}\int\phi(\gamma,y)\besker{\sigma}{S}{\gamma}{y}\,dy$. A continuous linear map
commutes with these sums and integrals, so by the hypothesis on $L$ and the substitution $y\mapsto\eta^{-1}y$,
\begin{align*}L\bigl(\ampsum{\sigma}{S}{\Xi}{c_1}{s}{c_1'}\bigr)&=\sum_{\gamma\in X_1}\int\cutoff{S}{c}(y)\,\cutoff{S}{c'}\bigl((\eta\gamma\eta'^{-1})^{-1}y\bigr)\,
\besker{\sigma}{S}{\eta\gamma\eta'^{-1}}{y}\,dy\\&=\sum_{\gamma\in X_1}\amp{\sigma}{S}{c}{c'}{\eta\gamma\eta'^{-1}}.\end{align*}
As $\gamma\mapsto\eta\gamma\eta'^{-1}$ maps $X_1$ bijectively onto $\Xi\cap\terms{S}{c}{t\,s\,t'^{-1}}{c'}$, this is the left
side of the assertion.

(b) By (S4) there is $x\in\UT{\Trat{S}{E}}$ with $u^{-1}\pos{S}{c}{c}{x}\in\KgS{S}$ and
$\uchar{S}{\rfin{S}(x)}=\uchar{S}{u}$; then $\pos{S}{c}{c}{x}\in u\,\KgS{S}$, so $x\in\terms{S}{c}{\redS{S}(u)}{c}$.
Let $x'$ be the same for $u'$ and $c'$. Apply (a) with $\eta=x$, $\eta'=x'$, $c_1=c$, $c_1'=c'$,
$t=\redS{S}(u)$, $t'=\redS{S}(u')$: the condition on $\Xi$ holds by stability, and by Lemma~\ref{l:Winv}(b) the
multiplication by $\uchar{S}{u}^{-1}\uchar{S}{u'}$ can be taken for $L$.
\end{proof}
\leanline{\lean{amplitudeSummove}\sep\lean{amplitudeSumbesselUmul}}

\begin{proof}[Proof of Theorem F]
(i) Suppose $\ampsum{\sigma}{S}{\Xi}{c}{s}{c'}\ne0$, and let $u,u'\in\KS{S}\cap\UT{\Tlat{S}{\AEf}}$ with
$\redS{S}(u)\,s=s\,\redS{S}(u')$. Then $\redS{S}(u)\,s\,\redS{S}(u')^{-1}=s$, and Proposition~\ref{p:moves}(b) gives
$\ampsum{\sigma}{S}{\Xi}{c}{s}{c'}=\uchar{S}{u}^{-1}\uchar{S}{u'}\,\ampsum{\sigma}{S}{\Xi}{c}{s}{c'}$, hence
$\uchar{S}{u}=\uchar{S}{u'}$. So $s$ agrees.

(ii) Write $\Xi=\ST{\Trat{S}{E}}$; it is stable by Lemma~\ref{l:stableS}. A term $\gamma\in\Xi$ over $(c,1,c')$ has
$\pos{S}{c}{c'}{\gamma}\in\KgS{S}$, so by (S3) $c=c'$ and $\gamma\in\clU{S}{c}$. Thus for $c\ne c'$ the sum is empty.
Let $c=c'$, and let $\phi(\gamma,y)=\cutoff{S}{c}(y)$ if $\gamma=1$ and $\phi(\gamma,y)=0$ otherwise. For
$\gamma\in\Xi\cap\terms{S}{c}{1}{c}\subseteq\clU{S}{c}$ the condition $u\gamma u'^{-1}=1$ determines $u'=u\gamma$, so
$\sum_{(u,u')}\phi(u\gamma u'^{-1},uy)=\sum_{u\in\clU{S}{c}}\cutoff{S}{c}(uy)=1$. By Theorem~E, with
$\besker{\full}{S}{1}{y}$ of coordinates $\besfn{\full}{S}{y}{i}\,\cj{\besfn{\full}{S}{y}{j}}$,
\[\ampsum{\full}{S}{\Xi}{c}{1}{c}_{(i,j)}=\int\cutoff{S}{c}(y)\,\besfn{\full}{S}{y}{i}\,\cj{\besfn{\full}{S}{y}{j}}\,dy.\]
By (S7), $\Wsys{\full}{S}{\tau}(ny)_{i_\tau}\cj{\Wsys{\full}{S}{\tau}(ny)_{j_\tau}}=\Wsys{\full}{S}{\tau}(y)_{i_\tau}\cj{\Wsys{\full}{S}{\tau}(y)_{j_\tau}}$
for $n\in\mathcal U$, the factor $e^{\sd{n}-\cj{\sd{n}}}$ having absolute value $1$. So we may unfold as in
Lemma~\ref{l:cutprod} (for integrable complex functions, by Lemma~\ref{l:unfold} and Fubini's theorem): the
integral over $\prod_\tau\mathcal U$ of the cutoff is $\ucov{S}$ by Theorem~C, and there remains
\[\ucov{S}\cdot\prod_\tau\int\Wsys{\full}{S}{\tau}(y)_{i_\tau}\,\cj{\Wsys{\full}{S}{\tau}(y)_{j_\tau}}\,d\muq{\fs}(y)=\ucov{S}\prod_\tau\delta_{i_\tau j_\tau}\]
by the orthonormality (J1).
\end{proof}
\leanline{\lean{amplitudeSumeqzeroofnotagrees}\sep\lean{amplitudeSumbesselSone}}

\sect{Notes on the formal text}{s:notes}
\subhead{Arguments that the census reads differently from Lean.} The notations of this note show the explicit
arguments of the Lean definitions. For four definitions the census of the artifact, which reads the arguments
from the text, lists other arguments than Lean takes, and the notation follows Lean: integral lifts
(Definition~\ref{d:lifts}; Lean takes the class, $\gamma$ and $m$, the census adds the bound matrix $t$), and the
three weights of Definition~\ref{d:weights}(a)--(c) (the census adds a variable $k$ of the file, because the
definitions contain $\kk{\fs}$). Likewise the census lists a matrix variable of the file (and, for the first of them, instances of the file's
field) among the arguments of
Lemma~\ref{l:intcount}, Lemma~\ref{l:sumone}(c) and Theorem~B, and a type among those of Lemma~\ref{l:eichler}, where
the statements only bind variables of the same name; and it stops at the first \texttt{let} of the statements
of Theorems A, D and E and of Lemma~\ref{l:liouville}, whose remaining arguments were read in the Lean text. The
notation $\latd{S}{c}$ is a function of the subspace, as in Lean. The discriminant $\sld{H}{V}$ of Note~S takes in
Lean the proof that $H$ has the left inverse $L$; we show $H$.

\subhead{The side.} In Lean the side is an instance argument with the default value ``full''. Theorems D and E
compare the kernel of the side $\sigma$ with the cut kernel; for $\sigma=\cut$ their second assertions are empty.
Theorem~F(ii) is stated for the full side only, as in Lean.

\subhead{Where the formal statements are more general or more special.} Theorem~A needs no admissibility, and
its number $x$ is a natural number. In Lemma~\ref{l:liouville} the hypothesis on $x$ is conditional on $\delta$
being a unit, although under the other hypotheses $\delta$ always is one. In Theorem~B the element $s$ only
serves to make the $\gamma_p$ terms over a common $s$. Proposition~\ref{p:boxexists} is stated in Lean for the
function selected by a choice operator, as here. Lemma~\ref{l:formulas}(g) holds over any field. In
Lemma~\ref{l:framediff} the second element enters only through its matrix. Lemma~\ref{l:levels}(c) is stated in
Lean as the identity $\padl{j}{Z}\padl{j}{Z}^{-1}=1$ for the matrix inverse of the library, which is $0$ for
singular matrices.

\subhead{What the three files do not use themselves.} Lemmas \ref{l:gramlink}, \ref{l:liouville},
\ref{l:colourcoset}, \ref{l:intcount}, \ref{l:scaled}, \ref{l:levels}, \ref{l:xrows}, \ref{l:sumint},
Lemma~\ref{l:lagr}, Lemma~\ref{l:posdef} and the definitions of Sections \ref{s:links} and~\ref{s:count} are used
only by later files of the unit (the pairing of links and the energy estimates). Theorem~A is used in the
estimate of pairs of terms, Theorems B to F throughout the second moment.

\subhead{Proofs that differ from the formal ones.} Lemma~\ref{l:lagr}(a): the formal proof applies the rank
theorem to a basis of the link, we use that the link is its own orthogonal. Theorem~B: the formal proof takes a
subset on which the colouring is injective with the same image, and measures distances from a chosen element;
this is Step~2. Theorem~C: Lemma~\ref{l:additive} is inside the formal proof, where the vanishing of the entries
is found by a decision procedure from the same relations between rows. Theorem~D: the comparison of the
exponents at the end is done by a decision procedure in Lean and by hand here. Theorem~E: the formal proof uses
the library's measure on a quotient and the fact that the distribution of an invariant function does not
depend on the cutoff (unit~114); identity~(3) is the same fact, proved directly. In Proposition~\ref{p:boxexists}
the formal induction adds the pairs with the row index decreasing, as here.

\subhead{Related work named by the authors.} The comments at the head of the first file name, as related
work on the sorting of double cosets by unipotent characters and on second moments, papers of Friedberg (1987), Y.~Liu
(2014), Furusawa and Shalika (2003), Michel, Ramakrishnan and Yang (2023), Kowalski, Saha and Tsimerman (2012)
and Blomer (2023). These are the authors' pointers; they were not opened for this note, and the formal proofs
do not depend on them.

\subhead{How this note was written.} From the dossier of the unit: the texts of the members of the three files
and their structured proofs. I read the three Lean files completely. The proofs of Theorems A, B, C and D and of
Proposition~\ref{p:boxexists}, which the readers of the dossier marked as their least certain, I read line by
line in the Lean text, and the proofs given here follow those lines; the full statements of Theorems A, B, D, E
and of Lemma~\ref{l:liouville} were read in the source, beyond the point where the census cuts them. Every
statement of this note was compared with its Lean statement, argument by argument. Of the Inputs I read the
Lean statements (for Note~S in its files, for the other units in theirs), and not the proofs; the description of
the Inputs (J2) to (J9) as elementary is the catalogue's. Two remarks on the dossier: its reading of the
induction in Proposition~\ref{p:boxexists} says that the pairs are added with increasing row index, while the Lean
text adds them with decreasing row index; and the statement of Theorem~D is complete in its text although the
census lists only the side among its arguments. Nothing was compiled for this note. The program check of the
unit's dossier covered \UnitSteps{} leaf steps of \UnitProofs{} proofs.

\head{Appendix A. Concordance with the Lean artifact}
The formal names of the objects: \texttt{S.b}, \texttt{S.N}, \texttt{S.g}, \texttt{S.piBar}, \texttt{S.ϖ} are
$\bS{S}$, $\NS{S}$, $\gS{S}$, $\piS{S}$, $\vpS{S}$; \texttt{𝔰.n}, \texttt{𝔰.k}, \texttt{𝔰.r} are $\nn{\fs}$, $\kk{\fs}$,
$\rr{\fs}$; an infinite place \texttt{v} is $\tau$ and \texttt{Φ v x} is $x_\tau$; \texttt{numInfPlaces E} is
$\ninf{E}$; \texttt{besselU}, \texttt{besselU0}, \texttt{besselH}, \texttt{besselS} are $\UT{T}$, $\UoT{T}$, $\HT{T}$,
$\ST{T}$; \texttt{classUnip}, \texttt{classUnip₀}, \texttt{classCutoff}, \texttt{unipCovolume} are $\clU{S}{c}$,
$\clUo{S}{c}$, $\cutoff{S}{c}$, $\ucov{S}$; \texttt{termsOver}, \texttt{termPosition}, \texttt{ratToInf},
\texttt{ratToFin} are $\terms{S}{c}{s}{c'}$, $\pos{S}{c}{c'}{\gamma}$, $\rinf{S}$, $\rfin{S}$; \texttt{genIndex} is
$\gix{\cdot}$; \texttt{frameVector y a} is $y^{-1}\fr{a}$ and \texttt{frameImage} $\gamma$ is $\gamma\mathfrak f$; a set
\texttt{T} of rational elements is $\Xi$; the elements \texttt{σ}, \texttt{σ'} of Proposition~\ref{p:moves} are
$\eta$, $\eta'$; the types \texttt{σ}, \texttt{τ}, \texttt{τ'}, \texttt{ι}, \texttt{m} of indices are $\Omega$,
$\Omega'$; the type \texttt{α} is $\sP$. Three statements, as the artifact has them:
\quoted{tsumclassCutoffboxle}
\quoted{lintegralclassCutoffequnipCovolume}
\quoted{amplitudeSumbesselSone}
{\fontsize{8.6}{11.6}\selectfont\renewcommand{\leanfont}{\ttfamily\fontsize{7.7}{10}\selectfont\color{leancol}}
\setlength{\LTleft}{0pt}\setlength{\LTright}{\fill}\setlength{\tabcolsep}{0pt}
\begin{longtable}{@{}p{27mm}p{60mm}p{42mm}r@{}}
\textbf{In this note} & \textbf{Lean declaration} & \textbf{File} & \textbf{Lines}\\\midrule\endhead
\ConcordanceRows
\end{longtable}}

\head{Notation}
\begingroup\fontsize{8.6}{10.9}\selectfont\raggedright
\setlength{\columnsep}{9mm}\setlength{\parskip}{0pt}
\begin{multicols}{2}
\nt{$E$, $x\mapsto\cj{x}$, $x_\tau$}{a CM field with a CM type, its conjugation, the embedding at the infinite place $\tau$}
\nt{$\ninf{E}$, $d_E$, $\volO{E}$}{number of infinite places; discriminant; $2^{\ninf{E}}/\sqrt{|d_E|}$}
\nt{$\AEf$, $\hO$, $\Fw{w}$}{finite adeles; closure of the integers; residue field of a finite place $w$}
\nt{$\fs=(\varrho,\kappa,k_0)$}{a shape; $\kk{\fs}$, $\nn{\fs}$, $\rr{\fs}$, $\Ix{\fs}$, $\Mx{\fs}$ as in \ref{r:space}}
\nt{$e_a$, $x_i$, $z$, $e^*_a$; $X$, $X_0$}{the basis of the space; the spans of the $e_a$ and of the $x_i$}
\nt{$\fr{a}$, $\mathfrak f$, $\lev$}{the frame, its matrix, the level of an index}
\nt{$T$, $\JT{T}$, $h_T$, $\GT{T}$}{an $\fs$-form over a ring $R$ (or $A$), its Gram matrix, the form, the unitary group}
\nt{$\UT{T}$, $\UoT{T}$, $\HT{T}$, $\ST{T}$}{unipotent group, its part fixing the frame, the group of $X_0$, the Bessel subgroup}
\nt{$\Trg{\fs}{N}$, $\Gone{\fs}$, $\mathcal U$}{the diagonal form; the group at $\CC$; its unipotent group}
\nt{$\sd{g}$, $\fm{g}$, $\ZT{T}{g}$, $\Zr{y}$, $\lmin{l}{Z}$}{sum of the superdiagonal entries on the frame; frame matrix; frame Gram matrices; leading minor}
\nt{$\bsize{y}$, $\InCut{y}{b}$, $\frbox{\rho}$}{size modulo the Bessel subgroup, the cut condition, the frame box}
\nt{$dy$, $\muU{\fs}$, $\muq{\fs}$, $\mu$}{Haar measures; the measure for the quotient by $\mathcal U$; in Section~\ref{s:covol} the product measure}
\nt{$S$, $S'$, $S_1$}{filling data; $\NS{S}$, $\qS{S}{i}$, $\gS{S}$, $\piS{S}$, $\vpS{S}$, $\bS{S}$, $\ChS{S}{i}$ as in \ref{r:data}}
\nt{$\Trat{S}{R}$, $\Tlat{S}{R}$, $\GS{S}{R}$, $\GFin{S}$}{forms and groups of \ref{r:groups}}
\nt{$\KS{S}$, $\KgS{S}$, $\GRes{S}$, $\redS{S}$, $\redim{S}{A}$}{integral subgroup, principal level, finite group, reduction, reduced image}
\nt{$\latB{S}$, $\rfin{S}$, $\rinf{S}$, $\bscale{S}{\tau}{i}$}{lattice basis; finite and real coordinates; the scales}
\nt{$\Cl{S}$, $c$, $c'$, $d$, $\tilde c$, $\hc{S}{c}$}{classes, their representatives}
\nt{$\clB{S}{c}$, $\clC{S}{c}$, $\clCX{S}{c}$, $\mbc{\fs}{p}$, $m_c$, $m(a)$}{class basis, class coordinates, their $X_0$-part; blocks of the class basis}
\nt{$\pos{S}{c}{c'}{\gamma}$, $\terms{S}{c}{s}{c'}$}{position; terms over $(c,s,c')$}
\nt{$\clU{S}{c}$, $\clUo{S}{c}$}{arithmetic unipotent group of level $\gS{S}$, its part fixing the frame}
\nt{$\gid{S}$, $\sdc{S}$, $\psi_f$, $\uchar{S}{u}$}{idele of $\gS{S}$; additive characters; the unipotent character}
\nt{$\rag{S}{(w,\epsilon)}$, $\ioia{S}{\epsilon}$, $\vfix{S}{w}{\mathbf d}$}{reductions, inverse adjoint, fixer of a vector of the frame}
\nt{$\FixN{s}{x}$, $\Agr{S}{s}$, $\IsStab{S}{\Xi}$}{$s$ fixes vectors at places of norm product $\ge x$; $s$ agrees; $\Xi$ is stable}
\nt{$\fpair{\gamma}$, $\ifpair{\gamma}$, $\lmat{\gamma}$, $\link{\gamma}$, $\zfp{T}$, $\clink{\gamma}{\gamma'}$}{frames, link matrix, link, the $z$-rows, common link}
\nt{$\ratG{S}$, $\lform{S}$, $\Zf{S}{\gamma}$, $\tdisc{S}{\gamma}$, $\fdisc{N}{Z}$, $\tdnorm{S}{\gamma}$}{Gram matrix, link form, frame products, discriminants}
\nt{$\IsNF{j}{\gamma}$, $\HasLev{S}{j}{\gamma}$, $\Cong{S}{c}{c'}{\gamma}{\gamma'}$}{normal form, level $j$, congruent terms}
\nt{$\fadm{M}$, $\slb{H}{V}$, $\sld{H}{V}$, $\gix{a}$}{a matrix over the adeles; lattice basis and discriminant of a subspace; generalized index}
\nt{$\Wsys{\sigma}{S}{\tau}$, $\Isys{S}{\tau}$, $\Bc{S}$, $\deg$}{Bessel functions of the system at $\tau$, their index sets, tuples of indices, the constant of the system}
\nt{$\besfn{\sigma}{S}{y}{i}$, $\besker{\sigma}{S}{\gamma}{y}$}{product Bessel function, kernel}
\nt{$\bxs{S}$, $\bxp{S}$, $\IsNr{S}{\gamma}$, $\IsBP{S}{c}{\omega}$}{box size, boxed points, near elements, box partitions}
\nt{$\cutoff{S}{c}$, $\bfl{S}{c}$, $\omS{S}{c}$}{the cutoff of the class $c$ and its two factors}
\nt{$\amp{\sigma}{S}{c}{c'}{\gamma}$, $\ampabs{\sigma}{S}{c}{c'}{\gamma}$, $\ampsum{\sigma}{S}{\Xi}{c}{s}{c'}$}{amplitude, its absolute version, amplitude sum}
\nt{$\kapS{S}{\tau}$, $\fscale{S}{\tau}{a}$, $\fsprod{S}$, $\tailS{S}$, $\ucov{S}$, $\cw{S}{i}$}{constants of \ref{r:const} and \ref{r:links}}
\nt{$\lmet{S}$, $\ldsq{S}{v}{e}$, $\HIL{c}{\gamma}{m}$}{link metric, link discriminant, integral lifts (Section~\ref{s:links})}
\nt{$\IdxB{j}$, $\xrows{Z}$, $\padl{j}{Z}$, $\fmid{\gamma}$}{Definition~\ref{d:pad}}
\nt{$\latd{S}{c}$, $\fvi{S}{c}{k}{a}{V}$, $\xri{S}{j}{k}$, $\sati{S}{k}$}{Definition~\ref{d:indices}; $\mathsf L_1$, $\mathsf L_2$, $T'$, $Z'$ there}
\nt{$\telem{S}{c}{p}$, $\pinv{S}{c}{c'}{\gamma}{\gamma'}$, $\mnorm{S}{j}{\gamma}$, $\cgauge{S}{j}{c}{c'}{\gamma}$, $\IsNRep{S}{c}{c'}{\gamma}$}{Definition~\ref{d:gauge}}
\nt{$\ipw{S}{c}{n}{p}$, $\termw{S}{j}{i}{c}{p}{p'}{V}$, $\pairw{S}{j}{c}{c'}{\gamma}{\gamma'}$, $\IsLab{S}{j}{c}{t}{p}$}{Definition~\ref{d:weights}}
\nt{$\Omega$, $\Omega'$, $\iota$; $\sP$, $\sA$, $\sB$, $\sC$, $\sw$, $\col$}{finite index sets; sets, colours and weights of the counting lemmas}
\nt{$\sigma$; $\Xi$; $\gamma,\gamma',\gamma_2,\gamma_p,\delta_p$; $u,u'$; $s,t,t'$}{a side; a set of rational elements; rational elements; unipotent elements; elements of $\GRes{S}$ (also matrices where said)}
\nt{$y,y',y^{\circ}$; $g$; $k$; $\eta,\eta'$}{points at infinity; a group element, or $\gS{S}$ where said; an element of $\KS{S}$ or $\GFin{S}$, or $\kk{\fs}$ where said; elements of the Bessel subgroup, or vectors in Section~\ref{s:agree}}
\nt{$a,b,l,i,j,j'$; $p,q$; $r$; $n$, $N$, $f$, $M$, $m$, $x$}{indices and levels; elements of $\sP$ or pairs, $p$ also a map in Lemma~\ref{l:eichler}; a real number or an adele, $\rr{\fs}$ where said; numbers, with the abbreviations said in each proof}
\nt{$C$, $Q$, $Q^{\circ}$, $\delta$, $\delta'$, $\rho$, $\mathfrak a$, $\mathbb F$, $o$, $\mathsf A$, $\xi,\zeta,\theta$, $\xi_0$, $\mu$, $v$, $\varsigma$, $\omega_0$, $\tau_0$, $\epsilon$}{objects of Theorem~A and its proof (Section~\ref{s:agree}) and of Input~(S2)}
\nt{$B$, $\lambda$, $\hat x$, $\hat u$, $J$, $J'$, $P$}{matrices and scalars of Section~\ref{s:eichler}; Gram matrices; a statement in $[P]$}
\nt{$\alpha$, $\alpha'$, $\beta$, $\beta_0$, $\phi$, $\varepsilon$}{scalars; cutoffs and weights; a distance}
\nt{$\theta$, $\mathsf I$, $\Sigma$, $F_\tau$, $E_{kl}$, $r_\tau$, $w_\tau$, $w$}{objects of Section~\ref{s:boxes}; $w$ also a finite place}
\nt{$F_p$, $F'$, $V$, $v_l$, $c_a$, $D$, $\mathfrak B$, $G$, $\Gamma$, $\varphi$, $\mathsf c$, $L_i$, $F_i$, $\mathrm{Fr}(a)$, $t_{l,\tau}$}{frames and spaces of Theorem~B; objects of Section~\ref{s:covol}}
\nt{$A_\tau$, $m_\tau$, $\theta_\tau$, $\mathbf A$, $\mathbf W$, $\Theta$, $C_0$, $X(d)$, $X_1$, $\Delta$, $\mathsf K$, $\mathsf Y$, $g_p$, $L$, $\mu_{\mathsf P}$, $\mathsf M$, $\mathsf L$}{objects of Section~\ref{s:ampl} and of the counting lemmas}
\nt{$\psi$, $\pi$, $\nu$, $\varpi$, $\mathcal C$, $H$, $O$, $R$, $A$, $U$, $W$, $Y$, $Z$, $K$, $I$, $e$, $h$, $z$, $d$}{letters inside the notations above ($\psi_f$, $2\pi i$, $\ucov{S}$, $\vpS{S}$, a compact set, a matrix, rings, and the symbols of the macros)}
\end{multicols}
\endgroup

\end{document}
